\chapter{Lean Export of Path Sums, Symbolic Sums, Inductions and Refutations}\label{ch:lean-export}\chaptermark{Lean Export of Path Sums and Refutations}

% Listing styles of this chapter: the shared literate list plus the characters the exporters emit.
\lstdefinestyle{lexlean}{style=lean,
  literate={→}{{$\to$}}1 {←}{{$\leftarrow$}}1 {≤}{{$\le$}}1 {≥}{{$\ge$}}1 {≠}{{$\ne$}}1
           {ω}{{$\omega$}}1 {⟩}{{$\rangle$}}1 {⟨}{{$\langle$}}1 {∀}{{$\forall$}}1
           {∃}{{$\exists$}}1 {λ}{{$\lambda$}}1 {∑}{{$\sum$}}1 {·}{{$\cdot$}}1
           {ℂ}{{$\mathbb{C}$}}1 {ℕ}{{$\mathbb{N}$}}1 {ℤ}{{$\mathbb{Z}$}}1 {ℚ}{{$\mathbb{Q}$}}1
           {↦}{{$\mapsto$}}1 {⊗}{{$\otimes$}}1 {∘}{{$\circ$}}1 {≫}{{$\gg$}}1
           {α}{{$\alpha$}}1 {β}{{$\beta$}}1 {π}{{$\pi$}}1 {∈}{{$\in$}}1 {×}{{$\times$}}1
           {₀}{{$_0$}}1 {₁}{{$_1$}}1 {₂}{{$_2$}}1 {⁻}{{$^{-}$}}1 {¹}{{$^{1}$}}1
           {√}{{$\surd$}}1 {∧}{{$\wedge$}}1 {∨}{{$\vee$}}1 {¬}{{$\neg$}}1 {↔}{{$\leftrightarrow$}}1
           {⟶}{{$\longrightarrow$}}1 {ε}{{$\varepsilon$}}1 {μ}{{$\mu$}}1 {σ}{{$\sigma$}}1
           {Δ}{{$\Delta$}}1 {Φ}{{$\Phi$}}1 {Ψ}{{$\Psi$}}1 {θ}{{$\theta$}}1 {ℝ}{{$\mathbb{R}$}}1
           {⊢}{{$\vdash$}}1 {∏}{{$\prod$}}1 {₃}{{$_3$}}1 {φ}{{$\varphi$}}1 {τ}{{$\tau$}}1
           {ρ}{{$\rho$}}1 {ζ}{{$\zeta$}}1 {Σ}{{$\Sigma$}}1 {Π}{{$\Pi$}}1 {−}{{$-$}}1 {∉}{{$\notin$}}1
,
}
\lstdefinestyle{lexinline}{style=inline,
  literate={→}{{$\to$}}1 {←}{{$\leftarrow$}}1 {≤}{{$\le$}}1 {≥}{{$\ge$}}1 {≠}{{$\ne$}}1
           {ω}{{$\omega$}}1 {⟩}{{$\rangle$}}1 {⟨}{{$\langle$}}1 {∀}{{$\forall$}}1
           {∃}{{$\exists$}}1 {λ}{{$\lambda$}}1 {∑}{{$\sum$}}1 {·}{{$\cdot$}}1
           {ℂ}{{$\mathbb{C}$}}1 {ℕ}{{$\mathbb{N}$}}1 {ℤ}{{$\mathbb{Z}$}}1 {ℚ}{{$\mathbb{Q}$}}1
           {↦}{{$\mapsto$}}1 {⊗}{{$\otimes$}}1 {∘}{{$\circ$}}1 {≫}{{$\gg$}}1
           {α}{{$\alpha$}}1 {β}{{$\beta$}}1 {π}{{$\pi$}}1 {∈}{{$\in$}}1 {×}{{$\times$}}1
           {₀}{{$_0$}}1 {₁}{{$_1$}}1 {₂}{{$_2$}}1 {⁻}{{$^{-}$}}1 {¹}{{$^{1}$}}1
           {√}{{$\surd$}}1 {∧}{{$\wedge$}}1 {∨}{{$\vee$}}1 {¬}{{$\neg$}}1 {↔}{{$\leftrightarrow$}}1
           {⟶}{{$\longrightarrow$}}1 {ε}{{$\varepsilon$}}1 {μ}{{$\mu$}}1 {σ}{{$\sigma$}}1
           {Δ}{{$\Delta$}}1 {Φ}{{$\Phi$}}1 {Ψ}{{$\Psi$}}1 {θ}{{$\theta$}}1 {ℝ}{{$\mathbb{R}$}}1
           {⊢}{{$\vdash$}}1 {∏}{{$\prod$}}1 {₃}{{$_3$}}1 {φ}{{$\varphi$}}1 {τ}{{$\tau$}}1
           {ρ}{{$\rho$}}1 {ζ}{{$\zeta$}}1 {Σ}{{$\Sigma$}}1 {Π}{{$\Pi$}}1 {−}{{$-$}}1 {∉}{{$\notin$}}1
,
}
\newcommand{\lexlean}{\lstinline[style=lexlean,frame=none,basicstyle=\ttfamily,breaklines=true,breakatwhitespace=false]}
\newcommand{\lexpy}{\lstinline[style=lexinline,language=Python]}
\newcommand{\lexcode}{\lstinline[style=lexinline]}

This chapter specifies the Python modules that turn the engine's verdicts and proof-kernel
theorems into \tool{Lean~4}{https://lean-lang.org}~\cite{demoura2021lean4} source files over
the ArchytasZX Lean library and \tool{Mathlib}{https://github.com/leanprover-community/mathlib4}~\cite{mathlib2020}.
Each exporter writes one self-contained file: literal data (path sums, networks, step lists,
tables), the side conditions of that data as Boolean checks closed by \lexlean!decide! or
\lexlean!decide +kernel!, scalar identities closed by \lexlean!linear_combination!, and a final
theorem whose proof applies a soundness lemma of the library. The Lean kernel then checks the
file; an exported verdict is Lean-certified only when the file compiles. A Lean-certified
verdict is about the engine's diagrams only when the export is diagram-level
(\py!kernel.diagram_level!: its final theorem states \lexlean!DEqual! or an \lexlean!Equal! of
\lexlean!Diagram! literals' \lexlean!toNet!); otherwise \code!certificate lean! adds the
encoding-level line. Exports that end without a diagram-level theorem: the path-sum networks
(\py!lean_pathsum.export_net!, \py!export_qudit_net!, \py!export_mixed_net!), the qudit digit
networks (\py!lean_qdigits.export_qudit_digits_net!, when the symbolic normaliser gives no
diagram-level export), the restricted chain and pair families (\py!lean_chainfam!,
\py!lean_pairs!), \py!lean_net.export! when no diagram literal is built, and a restricted
corollary (\py!lean_restrict!) over a bare final theorem. The exporters never
enumerate a path sum or a tensor contraction to produce an equality; the refutation exporters
evaluate at one point, which is their purpose.

The library itself (the semantics \lexlean!Net!, \lexlean!Equal!, \lexlean!PSem!, \lexlean!QSem!, the
symbolic normalizer, the linked diagrams) is the subject of \cref{ch:lean-core},
\cref{ch:lean-boxes}, \cref{ch:lean-pathsum} and \cref{ch:lean-sym}; the general network
exporter \module{archytaszx.semantics.lean_net} is \cref{ch:lean-export-net}. This chapter names
Lean declarations only as the exporters emit them.

\section{Coverage}\label{sec:lean-export-coverage}

\Cref{tab:lean-export-coverage} lists, for each verdict or theorem class, the exporter that
writes it, the Lean module the file imports, and whether the result is certified when the file
compiles. \enquote{Certified} means that the final theorem states the claimed equality (or
inequality) of the diagrams' Lean semantics; \enquote{trace only} means that the file checks
data but proves no value equality. The column \enquote{called by} names the engine module that
calls the exporter; \enquote{tests only} marks an exporter with no caller in the
\module{archytaszx} package.

{\footnotesize
\begin{longtable}{@{}>{\raggedright\arraybackslash}p{0.2\textwidth}>{\raggedright\arraybackslash}p{0.29\textwidth}>{\raggedright\arraybackslash}p{0.23\textwidth}>{\raggedright\arraybackslash}p{0.18\textwidth}@{}}
\caption{Coverage of the Lean exporters.}\label{tab:lean-export-coverage}\\
\toprule
Verdict or theorem class & Exporter & Lean module (imported) & Status \\
\midrule
\endfirsthead
\toprule
Verdict or theorem class & Exporter & Lean module (imported) & Status \\
\midrule
\endhead
Path-sum \azx!equal!, qubit or qudit, one canonical form &
\lexpy!export_pair! & \file{ArchytasZX.PathSum} & trace only; tests only (Lean-checked) \\
Path-sum \azx!equal!, qubit, no atoms, no [Quarter] &
\lexpy!export_proof!, \lexpy!export_sums! & \file{ArchytasZX.PathSumChain} & certified (value equality of path sums); tests only \\
Path-sum \azx!equal!, qubit, amplitude atoms or [Quarter], [Turn], [Average] on a context &
\lexpy!export_proof! (amp) & \file{ArchytasZX.PathSumChain} (\lexlean!PathSumAmp!) & certified (values); tests only \\
Path-sum \azx!equal!, odd prime $p$ &
\lexpy!export_qudit_proof! & \file{ArchytasZX.GaussSign}, \file{ArchytasZX.PathSumQuditChain} & certified (values); tests only \\
Diagram \azx!equal!: Z, X, F, B, S on $2^{k}$, W on 2 &
\lexpy!export_net! & \file{ArchytasZX.PathSumNetW} (+\,\file{PathSumSupport}) & certified \lexlean!Equal!; \lexcode!certify_decision!, \lexcode!kernel! \\
the same with amplitude atoms &
\lexpy!export_net! (amp) & \file{ArchytasZX.PathSumNetAmp} & certified \lexlean!Equal! \\
the same through [Quarter], [Turn] &
\lexpy!export_net! (quarter) & \file{ArchytasZX.PathSumNetAmp}, \file{PathSumChain} & certified \lexlean!Equal! \\
the same, both sides zero &
\lexpy!export_net! (zero plan) & \file{ArchytasZX.PathSumNetW} & certified \lexlean!Equal! \\
Diagram \azx!equal!: Z, X, F on one odd prime $p$ &
\lexpy!export_qudit_net! & \file{ArchytasZX.PathSumQuditNet} & certified \lexlean!Equal!; \lexcode!certify_decision!, \lexcode!kernel! \\
Diagram \azx!equal!: Z, X, F, B, S on $p^{k}$, some $k>1$ &
\lexpy!export_qudit_digits_net! & \file{ArchytasZX.PathSumQuditDigits} & certified \lexlean!Equal!; \lexcode!certify_decision!, \lexcode!lean_net! \\
Symbolic \azx!equal! (dimensions, phases, multiplicities) &
\lexpy!export_symbolic! & \file{ArchytasZX} (\lexlean!SymNet!, \lexlean!Sym!) & certified for every value; \lexcode!kernel!, \lexcode!certify_decision! \\
Qubit networks of Z, X, F without phases (inside \lexpy!export_symbolic!) &
\lexpy!q_source! & \file{ArchytasZX} (\lexlean!SymF2!, \lexlean!SymF2Bridge!) & certified \lexlean!Equal! \\
Diagram \azx!equal! on dimensions 2, 3 and 6 (mixed) &
\lexpy!export_mixed_net! & \file{ArchytasZX.PathSumMixedNet}, \file{PathSumChain} (+\,\file{PathSumSupport}) & certified \lexlean!Equal!; \lexcode!certify_decision! \\
Path-sum \azx!equal! whose canonical-form lift raises (parity route) &
\lexpy!export_sums! (parity) & \file{ArchytasZX.PathSumChain} & certified (values); tests only \\
\azx!unequal! at a point in $\Q(\zeta_M)$, by residues modulo a prime $p \equiv 1 \pmod M$ &
\lexpy!export_refutation! & \file{ArchytasZX} (\lexlean!Refute!, \lexlean!eval_res!, \lexlean!refute_res!) & certified $\lnot\forall$; \lexcode!decide!, \lexcode!kernel! \\
\azx!unequal! outside $\Q(\zeta_M)$ &
\lexpy!export_numeric_refutation! & \file{ArchytasZX} (\lexlean!Ball!, \lexlean!Refute!) & certified $\lnot\forall$ \\
\azx!unequal! by boundary dimensions &
\lexpy!export_interface_refutation! & \file{ArchytasZX} & certified $\lnot\forall$ \\
\azx!unequal! of \bbox families &
\lexpy!_instantiated_refutation! & as the instance's refutation & certified at one instance \\
Kernel peel step &
\lexpy!export_peel! & \file{ArchytasZX} (\lexlean!LinkedPeel!, \lexlean!LinkedGhost!) & certified (digit sums); \lexcode!kernel! \\
Kernel identity slide &
\lexpy!export_slide! & \file{ArchytasZX} (\lexlean!LinkedSlide!) & certified (digit sums); \lexcode!kernel! \\
Kernel regroup step &
\lexpy!export_regroup! & \file{ArchytasZX} (\lexlean!LinkedRegroup!) & certified (layout); \lexcode!kernel! \\
Kernel power split &
\lexpy!export_split! & \file{ArchytasZX} (\lexlean!LinkedSplit!) & certified (digit sums); \lexcode!kernel! \\
Induction through a chain cut &
\lexpy!export_chain! & \file{ArchytasZX.ChainInduct} (\file{ChainMixed}) & certified for every $k$; tests only \\
Induction through a tree window &
\lexpy!export_tree_theorem!, \lexpy!export_tree! & \file{ArchytasZX.TreeInduct} & certified for every depth; tests only \\
Periodic theorem &
\lexpy!export_periodic! & \file{ArchytasZX.PeriodicChain} & certified for every index; tests only \\
Restricted theorem &
\lexpy!corollary!, \lexpy!shift_corollary! & the source's & certified when the source is; \lexcode!kernel! \\
\bottomrule
\end{longtable}
}

Verdicts outside every row (for example a qubit diagram with a W node on dimension $2^{k}$,
$k>1$, a weighted qudit spider, or a refutation whose cyclotomic field has degree above 512) are
reported uncertified by the caller with the exporter's message.

\section{Conventions}\label{sec:lean-export-conventions}

\paragraph{Exceptions.} The path-sum, digit, chain and periodic exporters raise
\lexpy!PathSumUnsupported! (a subclass of \lexpy!PathSumError!, \module{archytaszx.algebra.pathsum});
the symbolic, refutation and ball exporters raise \lexpy!NetUnsupported!
(\module{archytaszx.semantics.lean_net}); scalar failures raise \lexpy!ScalarUnsupported!
(\module{archytaszx.semantics.lean_scalar}); the restriction reader raises
\lexpy!RestrictUnsupported!; the linked-diagram exporters return \lexpy!None!. Each message quoted in
\lexcode!code! type in this chapter is the exact text of the exception, with
\lexcode!<...>! standing for interpolated values.

\paragraph{Determinism.} Every exporter is a pure function of its inputs: dictionaries are
iterated in insertion order, nodes in sorted id order, wires in \lexpy!sort_key! order, and sets are
sorted before emission. Two runs on equal inputs give byte-identical files.

\paragraph{File shape.} A file starts with its imports, a module comment \lexcode!/-! ... -/!
stating what is proved, a \lexcode!namespace!, \lexcode!open! lines and \lexcode!set_option! lines, and ends
with \lexcode!end <namespace>!. Unless stated otherwise, a network proof file uses the options
\lexcode!maxRecDepth 100000!, \lexcode!maxHeartbeats 4000000!, and the linters
\lexcode!unusedSimpArgs!, \lexcode!unusedTactic! off.

% ---------------------------------------------------------------------------
\section{Exact Scalars and Scalar Identity Proofs}\label{sec:lean-export-scalar}

The module \module{archytaszx.semantics.lean_scalar} writes exact complex scalars as Lean
terms and proves identities between them. Every exporter of this chapter that states a
constant (a diagram scalar, a product of step factors, a weight, a coefficient of an atom) uses
it. A scalar is a \tool{SymPy}{https://www.sympy.org}~\cite{meurer2017sympy} expression over a
fixed set of uninterpreted atoms; an identity $a=b$ is proved by Mathlib's
\lexcode!linear_combination! tactic over polynomial relations that the atoms satisfy, with
cofactors found by Buchberger's algorithm~\cite{buchberger2006thesis}, so that Lean's
\lexcode!ring! normal form closes the goal. The module never evaluates a scalar numerically.

\subsection{Atoms}\label{sec:lean-export-scalar-atoms}

\begin{definition}[Scalar atoms]\label{def:lean-export-atoms}
The module defines the following SymPy function symbols and one symbol. Their Lean meanings
are those of the Lean library (\cref{ch:lean-core}).
\begin{center}
\begin{tabular}{@{}llp{0.4\textwidth}@{}}
\toprule
Python atom & Meaning & Lean text \\
\midrule
\lexpy!Rt(x)! & $\sqrt{x}$, $x$ a natural number & \lexlean!(rt x)! \\
\lexpy!Om(m)! & $\omega_m = e^{2\pi i/m}$ & \lexlean!(om m)! \\
\lexpy!Inv(e)! & $e^{-1}$ & \lexlean!(e)⁻¹! \\
\lexpy!Ex(t)! & $e^{t}$ & \lexlean!(Complex.exp t)! \\
\lexpy!GS(D, θ)! & $\sum_{k<D} e^{2\pi i\theta k}$ & \lexlean!(∑ k ∈ Finset.range D, ex (θ * (k : ℝ)))! \\
\lexpy!NC(n)! & a natural number cast to $\C$ & \lexlean!((n : ℕ) : ℂ)! \\
\lexpy!II! (symbol \lexcode!ii!) & the imaginary unit $i$ & \lexlean!Complex.I! \\
\lexpy!ModDelta(j, m)! & $[\,m \mid j\,]$ & \lexlean!(if j % m = 0 then (1 : ℂ) else 0)! \\
\lexpy!ModGcd(c, m)! & $\gcd(c,m)$ & \lexlean!((Nat.gcd c m : ℕ) : ℂ)! \\
\bottomrule
\end{tabular}
\end{center}
\lexpy!ModDelta! and \lexpy!ModGcd! come from \module{archytaszx.algebra.scalar}.
\end{definition}

\begin{signature}
\lexpy!class ScalarUnsupported(Exception)!: a scalar or identity outside what the module writes or
proves. Every exporter re-raises it unchanged or converts it to its own \lexcode!Unsupported!
exception.
\end{signature}

\subsection{The \texttt{Field} record}

\begin{signature}
\lexpy!@dataclass class Field!: \lexpy!theta: sp.Symbol!, \lexpy!poly: sp.Expr!,
\lexpy!relation: str!, \lexpy!forms: dict[sp.Expr, sp.Expr]!,
\lexpy!inverse: Callable[[sp.Expr], sp.Expr]!.
\end{signature}
A \lexcode!Field! describes a real algebraic number $\theta$ that is a root of the integer
polynomial \lexpy!poly! (an expression in \lexpy!theta!). \lexpy!relation! is the Lean proof text of
$\mathrm{poly}(\theta)=0$ over $\C$; \lexpy!forms! maps an engine scalar to its form over
$\theta$ and \lexpy!II!; \lexpy!inverse! maps a nonzero form to the form of its inverse. With a field
attached, \lexpy!model! replaces every key of \lexpy!forms! by its form, a negative integer power of a
form containing $\theta$ becomes a positive power of its inverse, and \lexpy!prove! adds the
relation $\mathrm{poly}(\theta)=0$.

\subsection{The \texttt{Scalars} renderer}

\begin{signature}
\lexpy!@dataclass class Scalars!:\\
\quad \lexpy!nat: Callable[[sp.Symbol], str]! --- the Lean binder of a natural-number symbol;\\
\quad \lexpy!other: Callable[[sp.Symbol], tuple[str, str]]! --- the Lean name and Lean type
(\lexlean!ℝ! or \lexlean!ℂ!) of any other symbol;\\
\quad \lexpy!positive: Callable[[str], str]! --- a Lean proof that a natural term is positive;\\
\quad \lexpy!nats: set[str]! (default empty) --- names of the natural-number symbols;\\
\quad \lexpy!opaques: dict[str, tuple[sp.Symbol, sp.Expr]]! (default empty) --- opaque symbols by
\lexpy!sp.srepr! key;\\
\quad \lexpy!alg: Field | None = None!.
\end{signature}
Exporters over path sums construct \lexpy!Scalars(nat=str, other=lambda s: (str(s), "ℂ"), positive=lambda x: "(by norm_num)")! (the helper \lexpy!_scalars()! of
\module{archytaszx.semantics.lean_pathsum}).

\paragraph{\texttt{model(expr)}.} Returns \lexpy!expr! rewritten over the atoms. It calls the
internal \lexpy!_model! and then \lexpy!text! on the result; if either raises
\lexcode!ScalarUnsupported!, the whole expression becomes an opaque symbol through
\lexpy!opaque(expr)!. \lexpy!opaque! keys an expression by \lexpy!sp.srepr! and returns the symbol
\lexcode!opaque_<k>!, $k$ counting from 1 in order of first use; equal expressions share one
symbol. An exporter that finds \lexpy!sc.opaques! nonempty after modeling a scalar it must
state exactly raises its own unsupported error (for example \lexcode!the scalar ... has no exact Lean form!).

\begin{algorithm}[\texttt{\_model}]\label{alg:lean-export-model}
For a SymPy expression $e$, the first matching case applies.
\begin{enumerate}
\item With a field and $e$ a key of \lexpy!alg.forms!: its form.
\item A number or a symbol: $e$ unchanged.
\item \lexpy!ModDelta! or \lexpy!ModGcd!: unchanged when every argument is a natural-number
expression over the symbols of \lexpy!nats! (\lexpy!_nat_arg!: \lexpy!nat_text! succeeds, every free
symbol is in \lexpy!nats!, and the argument is not a negative integer); otherwise
\lexcode!ScalarUnsupported("... has a non-natural argument")!.
\item $i$: \lexpy!II!. $\pi$ alone: \lexcode!ScalarUnsupported("a bare π")!. $e$ (Euler's number):
\lexpy!Ex(1)!.
\item A sum or product: the sum or product of the models of the arguments (through
\lexpy!model!, so a non-modelable argument becomes opaque).
\item A power $b^{x}$ whose exponent has free symbols: write $x = c + \sum_t c_t s_t$ (constant
plus terms); the result is $\mathrm{model}(b^{c}) \cdot \prod_t \mathrm{model}(b^{c_t})^{s_t}$;
a term that is not a coefficient times a symbol raises \lexcode!the exponent ...!.
\item $(-1)^{q}$, $q$ rational: with $r = (q/2) \bmod 1$, the result is
\lexpy!Om(r.q)**r.p!, or $1$ when $r=0$.
\item $b^{n}$, $n$ an integer: $\mathrm{model}(b)^{n}$ for $n\ge 0$; for $n<0$ the power of the
field inverse when the model contains $\theta$, else \lexpy!Inv(model(b))**(-n)!.
\item $b^{p/2}$ with $b$ natural (\lexpy!_natural!: a positive integer, or a non-numeric sum,
product, power or symbol all of whose symbols are in \lexpy!nats!): \lexpy!Rt(b)**p! for $p>0$,
\lexpy!Inv(Rt(b))**(-p)! otherwise.
\item Any other power: \lexcode!the power ...!.
\item A one-index \lexpy!sp.Sum! over $k$ from 0 to $h$ of $\exp(a)$ with
$a/(2\pi i) = \theta k$ and $\theta$ free of $k$: \lexpy!GS(h + 1, θ)! (checked with \lexpy!text!).
Any other sum: \lexcode!the sum ...!.
\item $\exp(a)$: with $t = a/(2\pi i)$ rational, \lexpy!Om(q)**p! for $t \bmod 1 = p/q \neq 0$,
$1$ when $t \bmod 1 = 0$; otherwise \lexpy!Ex(a)!.
\item Anything else: \lexcode!the scalar ... (<type>)!.
\end{enumerate}
\end{algorithm}

\paragraph{Rendering.} \lexpy!text(expr)! writes a $\C$ term: $1$ as \lexlean!1!; an integer $v$ as
\lexlean!(v : ℂ)! or \lexlean!(-|v| : ℂ)!; a rational $p/q$ as \lexlean!((p : ℂ) / q)!; \lexpy!II! as
\lexlean!Complex.I!; a natural symbol as \lexlean!(name : ℂ)!; another symbol as its name, cast
\lexlean!(name : ℂ)! when its type is not $\C$; sums and products parenthesized with
\lexcode! + ! and \lexcode! * !; a nonnegative integer power as \lexlean!(b ^ n)!; a power with a
symbolic exponent as \lexlean!(b ^ e)! with $e$ by \lexpy!nat_text!; then the atoms as in
\cref{def:lean-export-atoms}. Anything else raises \lexcode!cannot write ... in Lean!.
\lexpy!nat_text! writes integers, symbols (through \lexpy!nat!), sums, products and powers with a
nonnegative integer or non-numeric exponent as $\N$ terms; anything else raises \lexcode!... is not a natural number!. \lexpy!real_text! writes integers \lexlean!(v : ℝ)!, rationals
\lexlean!((p : ℝ) / q)!, real symbols (a symbol whose \lexpy!other! type is not \lexlean!ℝ! raises),
sums and products. \lexpy!complex_text! writes the argument of \lexpy!Ex!: $\pi$ as
\lexlean!(Real.pi : ℂ)!, $i$ as \lexlean!Complex.I!, sums, products and integer powers (negative
powers as \lexlean!(b ^ n)⁻¹!), and falls back to \lexpy!text!.

\paragraph{Nonzero proofs.} \lexpy!nonzero(atom)! returns a Lean proof of $\mathit{atom}\neq 0$:
\lexlean!(rt_ne x pos)! for \lexpy!Rt!, \lexlean!(om_ne_zero x pos)! for \lexpy!Om!,
\lexlean!(Complex.exp_ne_zero _)! for \lexpy!Ex!, \lexlean!Nat.cast_ne_zero.2 (Nat.pos_iff_ne_zero.1 pos)!
for \lexpy!NC! and natural symbols, \lexlean!(by norm_num : (v : ℂ) ≠ 0)! for integers, and
\lexlean!(by norm_num [Complex.ext_iff] : t ≠ 0)! otherwise. \lexpy!nonzero_expr(expr)! extends it to
products (\lexlean!mul_ne_zero!), nonnegative powers (\lexlean!pow_ne_zero!), inverses
(\lexlean!inv_ne_zero!), $1$ (\lexlean!one_ne_zero!), non-integer rationals (\lexlean!norm_num!), $i$
(\lexlean!Complex.I_ne_zero!) and, with a field, a form containing $\theta$ (through
\lexlean!left_ne_zero_of_mul_eq_one! and a proof that the form times its inverse is 1).

\subsection{Proving an identity}\label{sec:lean-export-scalar-prove}

\begin{algorithm}[\texttt{prove(lhs, rhs)}]\label{alg:lean-export-prove}
Returns a tactic block text \lexcode!by ...! proving \lexpy!text(lhs) = text(rhs)!.
\begin{enumerate}
\item \emph{Polynomial model.} Each atom (\lexpy!Rt!, \lexpy!Om!, \lexpy!Ex!, \lexpy!Inv!, \lexpy!GS!,
\lexpy!ModDelta!, \lexpy!ModGcd!, \lexpy!II!, every free symbol, and every power $b^{s}$ with a symbolic
exponent stripped of its integer constant part) becomes a fresh generator $a_0, a_1, \ldots$ in
order of first appearance; \lexpy!NC(x)! is modeled as the model of $x$; a power with symbolic
exponent $c+r$ requires $c$ a nonnegative integer and becomes $\mathrm{gen}(b^{r})\cdot
\mathrm{model}(b)^{c}$. Let $D = \mathrm{model}(\mathit{lhs}) - \mathrm{model}(\mathit{rhs})$.
\item \emph{Relations per atom} (a work list, processing atoms added on the way):
\begin{itemize}
\item $i$: $a^2+1$, proof \lexlean!Complex.I_sq!.
\item $\theta$ of the field: $\mathrm{poly}(\theta)$, proof \lexpy!alg.relation!.
\item $\sqrt{x}$: $a^2 - \mathrm{model}(x)$, proof \lexlean!(rt_sq' x)!.
\item $\omega_m$ with a non-numeric order: no relation. $\omega_m$ with a numeric order (a non-integer one raises \lexcode!a root of unity of symbolic order!): $a^m - 1$ (\lexlean!(om_pow_self m (by norm_num))!) and, for each prime $p \mid m$ with
$s=m/p$, the cyclotomic relation $\sum_{j<p} a^{js}$, proved by \lexlean!om_geom m p!.
\item $e^{-1}$: $\mathrm{model}(e)\cdot a - 1$, proof \lexlean!(mul_inv_cancel₀ (nonzero e))!.
\end{itemize}
\item \emph{Cross relations.}
\begin{itemize}
\item For orders $s<b$ of two \lexpy!Om! atoms with $s \mid b$ (the least such $b$):
$\omega_s - \omega_b^{b/s}$, proof \lexlean!(om_sub b s (b/s) ...)!.
\item With $\sqrt2$ present, for every $\omega_m$ with $8\mid m$, $q=m/8$:
$\sqrt2 - \omega_m^{q} - \omega_m^{7q}$ (\lexlean!rt_sub2!). With $\sqrt3$ present, for
$12 \mid m$, $q=m/12$: $\sqrt3 - \omega_m^{q} - \omega_m^{11q}$ (\lexlean!rt_sub3!).
\item With $i$ present, for $4\mid m$: $i - \omega_m^{m/4}$ (\lexlean!I_sub!).
\item With $\sqrt5$ present, for $5 \mid m$, $q=m/5$: $\sqrt5 - 1 - 2\omega_m^{q} -
2\omega_m^{4q}$ (\lexlean!rt_sub5!).
\item For every $\sqrt p$ with $p$ a prime other than 2, 3, 5 and every $\omega_m$ with $p\mid m$:
the quadratic Gauss-sum relation~\cite{berndt1998gauss} of \lexpy!_gauss_relation! below, when it
exists.
\item For every pair $e^{s}$, $(e^{-1})^{s}$ with the same symbolic exponent:
$e^{s}(e^{-1})^{s} - 1$, proved by \lexlean!← mul_pow, mul_inv_cancel₀, one_pow!.
\end{itemize}
\item If there are no generators or no relations: \lexcode!by ring! when $D$ expands to 0, else
\lexcode!ScalarUnsupported("... does not hold")!.
\item \emph{Ideal membership.} \lexpy!_lift(D, relations, generators)! returns cofactors $c_i$ with
$D = \sum_i c_i r_i$, or \lexpy!None!, which raises
\lexcode!ScalarUnsupported("... is not implied by the relations")!.
\item \emph{Output.} For each nonzero cofactor $c_i$: a line \lexcode!have r<i> := <proof>! (or
\lexcode!have r<i> : <statement> := (by ...)! for a proof that carries its statement), and the term
\lexcode!(<text(c_i)>) * r<i>!. With no nonzero cofactor the result is \lexcode!by ring!; otherwise
\lexcode!by have r0 ...; ...; linear_combination <terms>!, the parts joined by \lexcode!; !.
\end{enumerate}
\end{algorithm}

\begin{algorithm}[\texttt{\_lift}: cofactors by Buchberger's algorithm]\label{alg:lean-export-lift}
Polynomials are \lexpy!sp.Poly! over $\Q$ in the generators, ordered by \lexcode!grevlex!.
\begin{enumerate}
\item The basis starts as the relations $r_i$, each with the unit cofactor vector $e_i$.
\item \emph{Reduction} of $(f, c)$ by the basis: while $f\neq 0$, take its leading monomial
$m$ with coefficient $\lambda$; the first nonzero basis element $(g, c_g)$ whose leading
monomial divides $m$ gives $t = (\lambda/\mathrm{LC}(g))\, m/\mathrm{LM}(g)$, $f \gets f - tg$,
$c \gets c - t c_g$; without a divisor, the leading term moves to the remainder.
\item \emph{Completion.} Pairs $(i,j)$, $j<i$, are popped from the end of a list. A pair whose
leading monomials are coprime (product criterion) is skipped. Otherwise the S-polynomial with
its cofactors is reduced; a nonzero remainder joins the basis with pairs against every earlier
element. The loop stops when no pair is left or after 400 popped pairs.
\item Reduce the target with zero cofactors; a nonzero remainder returns \lexpy!None!; otherwise
the result is the negated cofactors (the invariant is $f = D + \sum c_i r_i$).
\end{enumerate}
The bound of 400 pairs makes the procedure incomplete: an identity implied by the relations
can fail when the basis is not complete at the bound.
\end{algorithm}

\begin{algorithm}[\texttt{\_gauss\_relation(p, m, atoms, poly)}]\label{alg:lean-export-gauss}
For an odd prime $p$ and $p \mid m$, let $h=(p-1)/2$, $q=m/p$,
$g=\#\{\,j<h : p < 2(2j+1)\,\}$ and $\left(\frac{k}{p}\right)$ the Legendre symbol (0 at
$k\equiv 0$). The relation is
\[
  \sum_{k=0}^{p-1} \left(\tfrac{k}{p}\right)\omega_m^{kq} \;=\; i^{h}\,(-1)^{g}\,\sqrt{p}.
\]
$i^{h}$ is written $(-1)^{h/2}$ for even $h$; for odd $h$ through the atom $i$ when present,
else as $\omega_{m'}^{(m'/4)h}$ for the first atom $\omega_{m'}$ with $4 \mid m'$; with none of
these the relation is omitted (\lexpy!None!). The proof applies
\lexlean!gauss_rt p (h p + 2p)! (two arguments by \lexlean!decide +kernel!), rewrites $\omega_p$ by
\lexlean!om_sub m p q!, expands the sum with \lexlean!Finset.sum_range_succ! and the values
\lexlean!QSem.lg p k! and \lexlean!gneg p! by \lexlean!decide!, rewrites $i^{h}$ as needed, then
\lexlean!push_cast! and \lexlean!linear_combination g!.
\end{algorithm}

\begin{example}
For the qubit proof of $\sqrt2\cdot\omega_8 \cdot \omega_8^{7} = \sqrt2$ the model is
$D = a_0a_1a_1^{7}-a_0$ with $a_0=\sqrt2$, $a_1=\omega_8$; the relations include
$a_1^{8}-1$, $a_1^4+1$ (the prime 2 of 8) and $a_0^2-2$. The cofactor of $a_1^8-1$ is $a_0$,
and the output is
\lexcode!by have r1 := (om_pow_self 8 (by norm_num)); linear_combination ((rt 2)) * r1!.
\end{example}

% ---------------------------------------------------------------------------
\section{Path-Sum Exports}\label{sec:lean-export-pathsum}

The module \module{archytaszx.semantics.lean_pathsum} writes Lean files for a path-sum
verdict \azx!equal!~\cite{amy2018pathsums}: two diagrams whose path sums
(\cref{ch:semantics}) reduce to one canonical form. It offers three levels of statement.
\begin{enumerate}
\item \lexpy!export_pair! writes each side's reduction trace as one instance of a rule lemma of
\file{ArchytasZX/PathSum.lean} per step and the equality of the two canonical forms by
\lexlean!rfl!. It is a trace certificate, not a theorem about values.
\item \lexpy!export_proof! and \lexpy!export_sums! (qubits) and \lexpy!export_qudit_proof! (odd primes)
prove that the two path sums have one value at every boundary assignment: each side's trace
replayed as a step list whose side conditions Lean checks by \lexlean!decide!, the canonical
pivots as canon steps, and the constants' identity by \lexlean!linear_combination!
(\cref{sec:lean-export-scalar}). When a canonical form's XOR lift raises, the parity route
(\cref{sec:lean-export-parity}) compares the two phases through their difference.
\item \lexpy!export_net!, \lexpy!export_qudit_net! and \lexpy!export_mixed_net! prove the Lean
proposition \lexlean!Equal! of the two diagrams' networks (\cref{ch:lean-pathsum}): each
network's value is written as its fragments' path sum, pin steps take it onto
\lexpy!from_diagram!'s variables, and the steps of level~2 follow. \lexpy!export_mixed_net!
(\cref{sec:lean-export-mixed}) takes a left diagram on dimensions 2, 3 and 6.
\end{enumerate}
No export enumerates a sum. Every unsupported input raises
\lexpy!PathSumUnsupported! (\module{archytaszx.algebra.pathsum}) with the message quoted below;
an unprovable constant identity raises \lexpy!ScalarUnsupported!.

\subsection{Constants and literals}

\begin{center}
\begin{tabular}{@{}lll@{}}
\toprule
Name & Value & Meaning \\
\midrule
\lexpy!PIN_OFFSET! & $2^{20}$ & first digit variable of a summed wire pinned to another variable \\
\lexpy!QPIN_OFFSET! & $2^{18}$ & first label of a summed $\Z_p$ wire whose class variable is taken \\
\lexpy!CHAIN_BLOCK! & 20 & steps between two checkpoints of a \lexlean!run_chain! (quarter networks) \\
\lexpy!RUN_BLOCK! & 4 & steps between two checkpoints of a side's reduction chain \\
\lexpy!NODE_CHUNK! & 25 & node facts per chunk theorem of a network fact \\
\lexpy!MIX_FRESH! & $2^{12}$ & offset of a mixed side's X-spider variables past its label base \\
\lexpy!MIX_ENT! & $2^{13}$ & offset of a mixed side's indicator-entry variables past its label base \\
\lexpy!MIX_SPAN! & $2^{16}$ & labels of a mixed side past its label base \\
left label base & $2^{16}$ & label of the left network's summed wire $w$ is $2^{16}+w$
(\lexpy!MIX_LABEL0! for a mixed side) \\
right label base & $2^{17}$ & the same for the right network
(\lexpy!MIX_LABEL0!${}+{}$\lexpy!MIX_SPAN! after a mixed side) \\
\bottomrule
\end{tabular}
\end{center}
\lexpy!DK! is the text \lexcode!(by decide +kernel)!, the proof of the decided side conditions of
the Z, X, F, B and S fragment lemmas. \lexpy!MIX_DIMS! is the text \lexcode!if k = 1 then 2 else k!,
the body of \lexlean!Δ! in a mixed network proof.

\paragraph{Options.} The list \lexpy!PROOF_OPTIONS! holds
\lexcode!set_option maxRecDepth 100000! and \lexcode!set_option maxHeartbeats 0!. \lexpy!NET_OPTIONS! is \lexcode!maxRecDepth 100000!,
\lexcode!maxHeartbeats 4000000!, \lexcode!linter.unusedSimpArgs false! and
\lexcode!linter.unusedTactic false!. \lexpy!QNET_OPTIONS! adds \lexcode!linter.unreachableTactic false!
and \lexcode!linter.unusedVariables false! to \lexpy!NET_OPTIONS!.

\paragraph{Literals.} A rational $v$ is \lexlean!(p / q : ℚ)!, or \lexlean!(v : ℚ)! when integral. An
$\mathbb{F}_2$ polynomial is the list of its monomials, each a sorted variable list, monomials sorted
by their sorted lists (\lexpy!_f2!, \lexpy!f2_text!); \lexpy!None! is \lexlean![]!. A qubit path sum with
phase modulus $M$ is a \lexlean!PSem.Data! literal (\lexpy!_data_term!):
\begin{lstlisting}[style=lexlean]
{ path := [v1, ..., vk], phase := [([vars], c), ...], cons := [[monos], ...], rows := [] }
\end{lstlisting}
with the path sorted, each phase term's coefficient written as the integer
$\lfloor cM\rfloor \bmod M$ (units of $1/M$ turn), terms ordered by monomial degree then sorted
variables, and each constraint as its monomial list. \lexpy!_data_text(ps, m, path)! is the same
literal with a given path order. A qubit canonical form is a \lexlean!PSCanon! literal
\lexcode!{ rows := [(p, [monos])...], phase := [([vars], q)...], const := [q0, q1, ...], n := n }!,
the constant as its reduced coefficient list over $\zeta_n$ (\lexpy!Cyclo.reduced(n)!); a qudit
canonical form is a \lexlean!QCanon! literal with monomials as lists of
(variable, exponent) pairs and an extra field \lexcode!p := p!. A qudit path sum is a
\lexlean!QSem.Data! literal whose polynomials are lists \lexcode!([(v, e), ...], c)! from
\lexpy!ModPoly.key()! (\lexpy!_qdata_term!).

\subsection{The trace certificate \texttt{export\_pair}}

\begin{signature}
\lexpy!export_pair(left: Diagram, right: Diagram, name: str = "pair") -> str!
\end{signature}
\begin{algorithm}[\texttt{export\_pair}]
\begin{enumerate}
\item Compute the raw path sums by \lexpy!from_diagram!. When both are qudit sums: reduce
(\lexpy!reduce_qudit!), take canonical forms (\lexpy!canonical_qudit!); raise
\lexcode!the qudit path sums do not reduce to one canonical form! if either is \lexpy!None!, they
differ (\lexpy!same!) or the form is zero. When both are qubit sums: the same with \lexpy!reduce! and
\lexpy!canonical! and the message \lexcode!the path sums do not reduce to one canonical form!.
Mixed bases raise \lexcode!the two path sums have different bases!.
\item $n = \mathrm{lcm}$ of the two constants' orders; write both canonical forms over $\zeta_n$.
\item Emit, under \lexlean!import ArchytasZX.PathSum! and namespace
\lexlean!ArchytasZX.PathSumCert.<name>!: the left trace, the definitions \lexlean!canonL!,
\lexlean!canonR! of type \lexlean!PSCanon! or \lexlean!QCanon!, the theorem
\lexlean!canon_eq : canonL = canonR := rfl!, and the right trace.
\end{enumerate}
\end{algorithm}
Each trace step becomes a comment line and one \lexlean!example! (\lexpy!_step_line!):
\begin{lstlisting}[style=lexlean]
-- L step 3: [hh_pivot] var 5, pivot 2, value [[1], [4]]
example := @ps_hh_pin
\end{lstlisting}
The comment lists the variable, then (when set) the pivot, the value, the digit group with its
fresh variables and digit polynomials, and the qudit polynomial. The body is
\lexlean!example := And.intro ps_omega0 ps_omega1! for \lexcode!omega!,
\lexlean!And.intro ps_omega_conj0 ps_omega_conj1! for \lexcode!omega_conj!,
\lexlean!example (S : ℕ) := ps_dyadic <|group|> S! for \lexcode!dyadic!,
\lexlean!example := @<lemma> p! (and the pivot for \lexcode!q_gauss!) for a qudit rule, and
\lexlean!example := @<lemma>! otherwise, with the lemma table
\begin{center}
\begin{tabular}{@{}ll@{\qquad}ll@{}}
\toprule
Rule & Lemma & Rule & Lemma \\
\midrule
\lexcode!elim! & \lexlean!ps_elim! & \lexcode!q_elim! & \lexlean!q_elim! \\
\lexcode!hh_pivot! & \lexlean!ps_hh_pin! & \lexcode!q_hh_pivot! & \lexlean!q_hh! \\
\lexcode!hh_gen! & \lexlean!ps_hh_gen! & \lexcode!q_hh_constraint! & \lexlean!q_hh! \\
\lexcode!rem! & \lexlean!ps_rem! & \lexcode!q_gauss! & \lexlean!q_gauss! \\
\lexcode!hh_constraint! & \lexlean!ps_hh! & & \\
\lexcode!hh_nl! & \lexlean!ps_hh_nl! & & \\
\bottomrule
\end{tabular}
\end{center}
A rule outside the table and outside \lexcode!omega!, \lexcode!omega_conj!, \lexcode!dyadic! (for
example \lexcode!bsub!, \lexcode!average!, \lexcode!pin!, \lexcode!quarter!) raises
\lexpy!PathSumUnsupported! (\lexcode!no PathSum.lean lemma for the R step!).

\subsection{Replaying a trace as \texttt{PSem.Step} constructors}\label{sec:lean-export-trace-step}

The value-level proofs run a list of \lexlean!PSem.Step! constructors
(\file{ArchytasZX/PathSumRules.lean}; \cref{ch:lean-pathsum}) on a \lexlean!Data! state whose path
is an ordered list. A step that sums out a variable acts on the head of the path, so each trace
step is preceded by \lexlean!.front x! steps, each moving one variable $x$ to the head of the path
(the qudit steps of \cref{alg:lean-export-qudit-proof} still use \lexlean!.perm!).
\lexpy!trace_step(step, cur)! returns the step texts and the path after them:
\begin{center}
\begin{tabular}{@{}p{0.2\textwidth}p{0.47\textwidth}p{0.25\textwidth}@{}}
\toprule
Trace rule & Emitted steps & Path after \\
\midrule
\lexcode!elim!, \lexcode!hh_constraint!, \lexcode!omega!, \lexcode!omega_conj! &
\lexlean!.front y!, then \lexlean!.elim!, \lexlean!.hhCons!, \lexlean!.omega!, \lexlean!.omegaConj! &
\lexcode!rest! (path without $y$) \\
\lexcode!rem! & \lexlean!.front y!, \lexlean!.rem! & \lexcode![y, rest]! \\
\lexcode!bsub! & \lexlean!.bsub y! & unchanged \\
\lexcode!dyadic! & \lexlean!.front g! for each group variable $g$ from the last to the first, then
\lexlean!.dyadicB group [(vars, c)...] fresh bits! & \lexcode![fresh, rest]! \\
\lexcode!dyadic_columns! & \lexlean!.front g! for each group variable from the last to the first,
then \lexlean!.columns group! & \lexcode!rest! (path without the group) \\
\lexcode!hh_pivot!, \lexcode!hh_gen!, \lexcode!hh_nl! (pivot $z$ set) &
\lexlean!.front z!, \lexlean!.front y!, then \lexlean!.hhPivot!, \lexlean!.hhGen!, \lexlean!.hhNl! &
\lexcode!rest! for \lexcode!hh_pivot!, else \lexcode![y, rest]! \\
any other & \lexcode!PathSumUnsupported("the rule <r> has no PathSumRules step")! & \\
\bottomrule
\end{tabular}
\end{center}
The \lexlean!.front! steps leave the group, or $y$ followed by $z$, at the head of the path in
order. In the dyadic step, the polynomial is the step's $S$ (each monomial's variable list
with its coefficient modulo $2^k$) and \lexcode!bits! the list of digit polynomials
\lexpy!step.values!. The step \lexlean!.columns ds! sums the head digits \lexcode!ds! of a dyadic
group as constraints (\lexlean!columnsS!, \file{ArchytasZX/PathSumRules.lean}).

\paragraph{Canon steps.} \lexpy!_steps(raw, red, order, boolean)! starts from the sorted path of
the raw sum, appends the trace steps, and ends with \lexlean!.canon v! for each pivot $v$ of
\lexcode!order!, or, when \lexcode!boolean!, with \lexlean!.gather!, \lexlean!.bcanon v! per pivot and
\lexlean!.pdrop!.

\begin{algorithm}[\texttt{\_elimination\_order}]\label{alg:lean-export-elim-order}
The pivots of \lexpy!pathsum.eliminate! in the order it takes them. Start from the nonzero
constraints. Repeat while constraints remain and the last sweep made progress: for each
boundary variable $v$ not yet taken, in increasing order, find the first constraint holding the
monomial $\{v\}$ and no other monomial containing $v$; if found, remove it, substitute
$v \mapsto$ (constraint without $\{v\}$) in the others by \lexpy!f2_substitute!, drop zero
constraints, and append $v$. \lexpy!_qudit_order! is the analogue over $\mathbb{F}_p$: a constraint
yields $v$ when \lexpy!_q_linear_pivot(c, {v}, -1)! returns a coefficient $a$, and the
substituted value is $-(c - a v)/a$.
\end{algorithm}

\subsection{Planning a qubit proof}

\begin{algorithm}[\texttt{\_plan(raw\_l, raw\_r, m\_min=1, support=False)}]\label{alg:lean-export-plan}
\begin{enumerate}
\item A side with amplitude atoms raises \lexcode!export_proof has no amplitude atoms!.
\item Reduce and canonicalize both sides. When \lexpy!canonical! raises
\lexpy!PathSumUnsupported! (an XOR lift that raises), raise \lexpy!_LiftRaises! holding both reduced
sums (\cref{sec:lean-export-parity}). Unequal, missing or zero canonical forms raise
\lexcode!the path sums do not reduce to one canonical form!.
\item The phase modulus $M$ is the least common multiple of 4, $m_{\min}$ and every
phase-coefficient denominator of both raw sums, and also of 8 when either trace contains
\lexcode!omega! or \lexcode!omega_conj!.
\item The elimination is Boolean when \lexpy!eliminate! returns \lexpy!None! on either reduced sum.
Then \lexpy!boolean_eliminate! yields rows and a leftover support $g$ per side; the pivots are
the left rows' variables; unequal supports, a support other than $1$ when \lexpy!support! is
false, or unequal pivot lists raise \lexcode!the Boolean eliminations differ or leave a support!.
Otherwise the pivots are \lexpy!_elimination_order! of the left side.
\item Return both reduced sums, the pivots, the Boolean flag, $M$ and both step lists
(\lexpy!_steps!).
\end{enumerate}
\end{algorithm}

\paragraph{Factor model.} The Lean function \lexlean!PSem.fac! assigns $2$ to \lexlean!.elim!,
\lexlean!.hhCons!, \lexlean!.hhPivot! and \lexlean!.hhNl!, $\sqrt2\,\omega_8$ to \lexlean!.omega!,
$\sqrt2\,\omega_8^{7}$ to \lexlean!.omegaConj!, $2^{|ds|}$ to \lexlean!.columns ds! and $1$ to every
other step. \lexpy!_fac_model(steps, M)! is the same product with $\omega_8$ written
$\omega_M^{M/8}$; it reads $|ds|$ of a step text \lexcode!.columns [...]! as its number of commas
plus one. \lexpy!_fac_lines(side, steps, M, text)! emits
\begin{lstlisting}[style=lexlean]
theorem facL : facs stepsL = <text> := by
  simp only [stepsL, facs, fac]        -- adds List.length_cons, List.length_nil after a .columns step
  rw [om_sub M 8 (M/8) (by norm_num) (by norm_num)]   -- only when M ≠ 8 and an omega step occurs
  ring
\end{lstlisting}
\lexpy!_const_turn(red, M, boolean)! substitutes the elimination rows into the reduced phase and
returns its constant term in units of $1/M$, modulo $M$. \lexpy!_const_model(c, k)! is
$c\cdot\sqrt2^{k}$ over the atoms (\lexpy!Inv(Rt(2))**(-k)! for $k<0$). \lexpy!_leftover(red)! is true
when \lexpy!eliminate! fails and \lexpy!boolean_eliminate! leaves a support other than $1$.

\subsection{The value-level proof \texttt{export\_proof}}

\begin{signature}
\lexpy!export_proof(left: Diagram, right: Diagram, name: str = "pair") -> str!\\
\lexpy!export_sums(raw_l: PathSum, raw_r: PathSum, name: str = "pair") -> str!
\end{signature}
\lexpy!export_proof! computes both raw sums by \lexpy!from_diagram!; a sum that is not a qubit
\lexpy!PathSum! raises \lexcode!export_proof covers qubit path sums only!. It returns
\lexpy!export_sums! of the two sums, which takes the remaining steps.
\begin{algorithm}[\texttt{export\_sums}]\label{alg:lean-export-proof}
\begin{enumerate}
\item If either sum has amplitude atoms, or either reduction takes a \lexcode!quarter!, \lexcode!turn!,
\lexcode!turn_held! or \lexcode!average_ctx! step (\lexpy!_quartered!), delegate to the amplitude proof
(\cref{sec:lean-export-amp-proof}).
\item Plan with \lexpy!support=True!. When the plan raises \lexpy!_LiftRaises!, return the parity
proof \lexpy!_export_parity_proof! (\cref{sec:lean-export-parity}). The support flag is
\lexpy!_leftover! of the left side.
\item Let $f_L, f_R$ be the factor models, $k_L, k_R$ the constant turns, $c_L, c_R$ the
constant models. Prove $f_L c_L \omega_M^{k_L} = f_R c_R \omega_M^{k_R}$ by \lexpy!prove!.
\item Emit the file below.
\end{enumerate}
\end{algorithm}
\begin{lstlisting}[style=lexlean]
import ArchytasZX.PathSumChain
namespace ArchytasZX.PathSumProof.<name>
open ArchytasZX ArchytasZX.PSem
set_option maxRecDepth 100000
set_option maxHeartbeats 0
def sL : Data := { path := [...], phase := [...], cons := [...], rows := [] }
def sR : Data := ...
def stepsL : List Step := [.front 3, .elim, ..., .canon 0]
def stepsR : List Step := [...]
theorem facL : facs stepsL = ... := by ...
theorem facR : facs stepsR = ... := by ...
theorem consts : (fL) * cL * om M ^ kL =
    (fR) * cR * om M ^ kR := by
  have key : <lhs> = <rhs> :=
    by ...                                  -- Scalars.prove
  linear_combination key
run_chain runL M stepsL sL 4
run_chain runR M stepsR sR 4
theorem main (σ : ℕ → ZMod 2) :
    PSem.val M cL sL σ = PSem.val M cR sR σ := by
  rw [run_sound M (by norm_num) stepsL sL _ σ runL_ok, runL_eq,
    run_sound M (by norm_num) stepsR sR _ σ runR_ok, runR_eq]
  apply final_sound M (by norm_num) _ _ _ _ σ kL kR (by decide) (by decide)
    (by decide)
  simpa only [facL, facR] using consts
end ArchytasZX.PathSumProof.<name>
\end{lstlisting}
\lexlean!finalS_sound! replaces \lexlean!final_sound! when a support is left. The command
\lexlean!run_chain! (\file{ArchytasZX/PathSumChain.lean}) elaborates the run of a step list in
blocks of \lexpy!RUN_BLOCK! steps and defines \lexlean!runL!, \lexlean!runL_ok! (every side condition
holds) and \lexlean!runL_eq! (the final state); \lexpy!_chain_lines(kind, M, startL, startR)! writes
the two commands for \lexcode!kind! \lexcode!run!, \lexcode!aruns! or \lexcode!qrun!.

\subsection{The parity route}\label{sec:lean-export-parity}

The parity route proves a qubit pair whose canonical form takes an XOR lift that raises
\lexpy!PathSumUnsupported! (\lexpy!canonical!, \cref{ch:semantics}). Its Lean side is
\file{ArchytasZX/PathSumParity.lean}, imported by \file{ArchytasZX/PathSumChain.lean}: the step
\lexlean!.canonH v! makes the row of pivot $v$ with the phase kept unsubstituted, and
\lexlean!finalP_sound M hM cL cR vs sL sR σ k h hk hc! concludes \lexlean!val M cL sL σ = val M cR sR σ!
from the decided check \lexlean!finalPOK M vs sL sR = true!, the constant turn
\lexlean!constK (diffP M vs sL sR) = k! of the phase difference with the rows of the pivots
\lexcode!vs! substituted, and \lexlean!cL * om M ^ k = cR!.

\begin{signature}
\lexpy!class _LiftRaises(PathSumUnsupported)!: message \lexcode!an XOR lift of the canonical form raises!;
fields \lexpy!red_l!, \lexpy!red_r! (the two reduced sums). \lexpy!_plan! raises it.
\end{signature}

\begin{algorithm}[\texttt{\_parity\_turn(red\_l, red\_r, M)}]\label{alg:lean-export-parity-turn}
\begin{enumerate}
\item Take \lexpy!canonical(red, parity=True)! of both sides; a missing or zero form raises
\lexcode!the path sums do not reduce to one canonical form!.
\item Let $D$ = \lexpy!parity_difference! of the two forms (left minus right, parity terms
included, shared rows substituted). When $D$ is \lexpy!None! or holds a nonconstant monomial,
raise \lexcode!the phase difference of the canonical forms is not constant!. When
$c_L\cdot e^{2\pi i D_0} \neq c_R$ for the forms' constants $c_L$, $c_R$ and the constant term
$D_0$ of $D$ (in turns), raise \lexcode!the canonical forms differ in their constants!.
\item Let the rows be \lexpy!eliminate! of the left constraints; \lexpy!None! or \lexpy![None]! raises
\lexcode!a constraint has no pivot!.
\item Form the left reduced phase minus the right reduced phase and substitute each row
$(\mathit{pivot}, \mathit{rest})$ in order. A nonconstant monomial left raises
\lexcode!the phase difference is not constant!.
\item Return $D'_0 M \bmod M$ (\lexpy!int! truncation of $D'_0 M$, then the remainder), with
$D'_0$ the constant term of that difference in turns.
\end{enumerate}
\end{algorithm}

\begin{algorithm}[\texttt{\_parity\_plan}]\label{alg:lean-export-parity-plan}
\lexpy!_parity_plan(raw_l, raw_r, held, m_min=1)! returns
$(M, k, \mathit{order}, \mathit{steps}_L, \mathit{steps}_R)$.
\begin{enumerate}
\item With the reduced sums of \lexcode!held!: when \lexpy!eliminate! returns \lexpy!None! on either side,
raise \lexcode!the parity route has no Boolean elimination!.
\item $M$ is the least common multiple of 4, $m_{\min}$ and every phase-coefficient
denominator of both raw sums, and also of 8 when either trace contains \lexcode!omega! or
\lexcode!omega_conj!.
\item $k$ = \cref{alg:lean-export-parity-turn}; \lexcode!order! = \lexpy!_elimination_order! of the
left reduced sum (\cref{alg:lean-export-elim-order}).
\item Each side's steps are \lexpy!_steps(raw, red, (), False)! (the trace steps without canon
steps) followed by \lexlean!.canonH v! for each $v$ of \lexcode!order!; both sides use the left
order.
\end{enumerate}
\end{algorithm}

\paragraph{\texttt{\_export\_parity\_proof(raw\_l, raw\_r, held, name)}.} Plan with
$m_{\min}=1$; let $f_L, f_R$ be the factor models and $c_L, c_R$ the constant models; prove
$f_L c_L\,\omega_M^{k} = f_R c_R$ by \lexpy!prove!. The file has the import, namespace, options and
definitions of \cref{alg:lean-export-proof}, with these differences:
\begin{lstlisting}[style=lexlean]
theorem consts : (fL) * cL * om M ^ k =
    (fR) * cR := by
  have key : <lhs> = <rhs> :=
    by ...
  linear_combination key
run_chain runL M stepsL sL 4
run_chain runR M stepsR sR 4
theorem main (σ : ℕ → ZMod 2) :
    PSem.val M cL sL σ = PSem.val M cR sR σ := by
  rw [run_sound M (by norm_num) stepsL sL _ σ runL_ok, runL_eq,
    run_sound M (by norm_num) stepsR sR _ σ runR_ok, runR_eq]
  apply finalP_sound M (by norm_num) _ _ [v1, ..., vr] _ _ σ k (by decide) (by decide)
  simpa only [facL, facR] using consts
\end{lstlisting}

\subsection{Proofs with amplitude atoms}\label{sec:lean-export-amp-proof}

A qubit path sum may carry amplitude atoms $a^{F}$ (a cyclotomic base $a$ raised to an $\mathbb{F}_2$
polynomial $F$), and the reduction rules [Quarter], [Turn] and [Average] on a context create
them. These sums are proved over \file{ArchytasZX/PathSumAmp.lean}, whose state \lexlean!AData!
holds an inner \lexlean!Data!, a list of held variables, and the atom list; \cref{sec:lean-export-amp}
defines the step writer.

\begin{signature}
\lexpy!@dataclass class AmpPlan!: \lexpy!coefs!, \lexpy!runs: (AmpRun, AmpRun)!, \lexpy!starts! (each side's
initial inner path, held variables and atoms), \lexpy!m: int!, \lexpy!turns: (int, int)!,
\lexpy!support: bool = False!; property \lexpy!final! is \lexcode!"afinalS_perm"! when \lexpy!support!, else
\lexcode!"afinal_perm"!.
\end{signature}

\begin{algorithm}[\texttt{amp\_plan(sides)}]\label{alg:lean-export-amp-plan}
\lexpy!sides! is \lexpy!((raw_l, run_l), (raw_r, run_r))! with initial runs.
\begin{enumerate}
\item Reduce and canonicalize both sides; a missing or zero form raises \lexcode!the path sums do not reduce to one canonical form!. The \emph{atom form} applies when the two canonical forms
differ (\lexpy!same! false); then their rows, supports and boundaries must agree, else the same
error.
\item Eliminate as in \cref{alg:lean-export-plan}; in the Boolean case unequal supports or pivot
lists raise \lexcode!the Boolean eliminations differ!. The canon steps are \lexlean!.gather!,
\lexlean!.bcanon v!\ldots, \lexlean!.pdrop! (Boolean) or \lexlean!.canon v!\ldots.
\item Snapshot both runs and the number of coefficient bases. Try the modes
(drop, split) $=(\mathrm{F},\mathrm{F})$, and additionally $(\mathrm{T},\mathrm{F})$,
$(\mathrm{T},\mathrm{T})$ when the elimination is Boolean with a support $g\neq 1$. For each
mode and side: set \lexpy!run.support! to $g$ when dropping, record the start state, replay the
trace (\lexpy!reduce_run!), append the canon steps, substitute the rows, and in the atom form
normalize to the canonical phase and support (with \lexcode!split!). Stop at the first mode whose
two \lexpy!final_key()! values agree; otherwise restore the snapshots (and truncate the bases) and
try the next mode. If none agrees, raise \lexcode!the reduced sums carry different amplitude atoms!.
\item Return the plan with $M$ = the coefficients' modulus, both constant turns, and
\lexpy!support! = Boolean with $g\neq 1$, plus the two reduced sums.
\end{enumerate}
\end{algorithm}
\lexpy!amp_modulus(raws, m_min=1)! is $\mathrm{lcm}(8, m_{\min}, \text{all phase denominators})$.

\paragraph{\texttt{\_export\_amp\_proof}.} With $M$ from \lexpy!amp_modulus!, a fresh
\lexpy!Coefs(M)! and \lexpy!initial! runs on both raw sums (held: the path variables an atom reads),
it plans, writes the factor theorems \lexlean!facL!, \lexlean!facR! by \lexpy!afac_lines!, proves
$f_L c_L\omega_M^{k_L} = f_R c_R\omega_M^{k_R}$, and emits under
\lexlean!import ArchytasZX.PathSumChain! and \lexpy!NET_OPTIONS!: the coefficient vector \lexlean!A!,
\lexlean!sL!, \lexlean!tL : AData := ⟨sL, held, atoms⟩!, \lexlean!stepsL : List AStep! (and the right
side), the coefficient identities \lexlean!reqsL!, \lexlean!reqsR!, the root identities
\lexlean!qreqsL!, \lexlean!qreqsR!, \lexlean!consts!, two \lexlean!aruns_chain! commands, and
\begin{lstlisting}[style=lexlean]
theorem main (σ : ℕ → ZMod 2) :
    valA M cL A tL σ = valA M cR A tR σ := by
  rw [aruns_sound M (by norm_num) A stepsL tL _ σ reqsL qreqsL arunsL_ok, arunsL_eq,
    aruns_sound M (by norm_num) A stepsR tR _ σ reqsR qreqsR arunsR_ok, arunsR_eq]
  apply afinal_perm M (by norm_num) _ _ A _ _ σ kL kR (by decide) (by decide)
    (by decide) (by decide) (by decide) (by decide)
  simpa only [facL, facR] using consts
\end{lstlisting}

\subsection{Qudit proofs \texttt{export\_qudit\_proof}}

\begin{signature}
\lexpy!export_qudit_proof(left: Diagram, right: Diagram, name: str = "pair") -> str!
\end{signature}
\begin{algorithm}[\texttt{export\_qudit\_proof}]\label{alg:lean-export-qudit-proof}
\begin{enumerate}
\item Both raw sums must be qudit sums (\lexcode!export_qudit_proof covers odd prime path sums only!), over one prime $p$ with one boundary (\lexcode!the two path sums have different bases or boundaries!), without atoms or table atoms (\lexcode!export_qudit_proof has no atoms!).
\item Reduce and canonicalize; failure raises \lexcode!the qudit path sums do not reduce to one canonical form!. The two \lexpy!_qudit_order! lists must agree (\lexcode!the qudit elimination orders differ!).
\item Each side's \lexlean!QSem.Step! list (\lexpy!_qudit_steps!): per trace step, \lexcode!q_elim! and
\lexcode!q_hh_constraint! give \lexlean!.perm [y, rest]! and \lexlean!.elim! or \lexlean!.hhCons!;
\lexcode!q_gauss! gives \lexlean!.perm [y, rest]!, \lexlean!.gauss z!; \lexcode!q_hh_pivot! gives
\lexlean!.perm [y, z, rest]!, \lexlean!.hhPivot!; another rule raises \lexcode!the rule <r> has no PathSumQudit step!. Then \lexlean!.canon v! per pivot and \lexlean!.norm!.
\item The constant turn (\lexpy!_qudit_const_turn!) substitutes the rows of \lexpy!eliminate_qudit!;
its pivots must equal the sorted order (\lexcode!the qudit elimination pivots differ!).
\item $n = \mathrm{lcm}(p, \text{both constant orders})$. The step factor of \lexlean!.elim!,
\lexlean!.hhCons!, \lexlean!.hhPivot! is $p$; of \lexlean!.gauss a! the Gauss sum
\lexlean!qg p a!, modeled in Legendre form
$G(a) = \left(\frac{a}{p}\right)\sum_{k=1}^{p-1}\left(\frac{k}{p}\right)\omega_n^{kn/p}$
(\lexpy!_legendre_model!)~\cite{berndt1998gauss}.
\item Prove $f_L c_L\omega_n^{k_L n/p} = f_R c_R \omega_n^{k_R n/p}$.
\end{enumerate}
\end{algorithm}
The file imports \lexlean!ArchytasZX.GaussSign! and \lexlean!ArchytasZX.PathSumQuditChain!, opens
the namespace \lexlean!ArchytasZX.PathSumQuditProof.<name>!, sets \lexcode!maxRecDepth 100000!,
\lexcode!maxHeartbeats 4000000!, turns off the linters \lexcode!unusedTactic!,
\lexcode!unusedSimpArgs! and \lexcode!unreachableTactic!, declares \lexlean!instance : Fact (Nat.Prime p) := ⟨by norm_num⟩!,
and contains: \lexlean!sL!, \lexlean!sR!, \lexlean!stepsL!, \lexlean!stepsR!; one theorem
\lexlean!lg<k> : lg p k = v := by decide! per residue $k<p$; per Gauss argument $a$ a theorem
\lexlean!qg<a> : qg p a = <Legendre form>! proved by \lexlean!qg_legendre!, the sum expansion, the
\lexlean!lg! facts, \lexlean!om_sub n p (n/p)! when $n\neq p$, \lexlean!push_cast! and \lexlean!ring!;
\lexlean!facL!, \lexlean!facR! (\lexlean!QSem.facs!); \lexlean!consts!; two \lexlean!qrun_chain! commands; and
\begin{lstlisting}[style=lexlean]
theorem main (σ : ℕ → ZMod p) :
    QSem.val p (cL) sL σ = QSem.val p (cR) sR σ := by
  rw [QSem.run_sound stepsL sL _ σ qrunL_ok, qrunL_eq,
    QSem.run_sound stepsR sR _ σ qrunR_ok, qrunR_eq]
  apply QSem.final_sound _ _ _ _ σ kL kR (by decide) (by decide) (by decide)
  first | (rw [facL, facR]; exact consts) | rw [facL, facR]
\end{lstlisting}

\subsection{Network proofs over qubit digits}\label{sec:lean-export-net}

\lexpy!export_net! proves \lexlean!Equal E sL NL sR NR! for two diagrams built from Z, X, F, B and S
nodes on dimensions $2^{k}$ and W nodes on dimension 2: the Lean network \lexlean!N! of each
diagram (\cref{ch:lean-pathsum}) has a value that equals its fragments' path sum over digit
variables (\lexlean!NetPSD.eval_eq!). A port of dimension $2^{k}$ carries $k$ digit variables.

\begin{algorithm}[\texttt{\_net\_side}: one diagram's network over digit variables]\label{alg:lean-export-net-side}
Inputs: the diagram, $M$, the shared phase-vector list \lexcode!phases!, the label base
\lexcode!label0!, the coefficient table \lexcode!coefs!, optionally \lexcode!amp! (a \lexpy!Coefs!), the
index tables \lexcode!tabs! and a preferred phase-index map \lexcode!prefer!.
\begin{enumerate}
\item Build the raw sum by \lexpy!from_diagram(diagram, lay, w_nodes=True)!; a non-qubit sum
raises \lexcode!the network export covers qubit path sums only!. \lexcode!lay! supplies the
normalized diagram, each port's digit variables (\lexcode!digit!), the fresh variables of X and
legless spiders (\lexcode!fresh!, \lexcode!zdigits!) and the W pins (\lexcode!wpins!).
\item \emph{Boundary labels.} The boundary ports (outputs, then inputs) get labels $0,1,\ldots$;
label $i$ gets the next $|\mathit{digit}|$ consecutive digit variables from 0, each mapped to
its class variable. Their total is the number of boundary variables.
\item \emph{Summed labels.} For each wire $w$ in \lexpy!sort_key! order (a wire on a boundary port
raises \lexcode!a wire on a boundary port!), the label is $\mathit{label0}+w$. Its digit $j$ is the
class variable $c$ of the port's digit $j$ when $c$ is a path variable (not boundary) not yet
taken; otherwise the next variable from \lexpy!PIN_OFFSET!.
\item \emph{Nodes.} The generators Z, X, F, B, S, W map to \lexlean!.Z!, \lexlean!.X!, \lexlean!.F!,
\lexlean!.B!, \lexlean!.S!, \lexlean!.W!; another generator raises \lexcode!node <id> is a <name>!. A legless
node other than Z or X raises \lexcode!node <id> has no legs!.
\item \emph{Constraint hyperedges.} A Z spider with ports joins, per digit $j$, the $j$-th
digits of all its labels (all labels must have one digit count, else \lexcode!node <id> joins legs of different dimensions!). A binder B (whole output, parts $a$, $b$ inputs) or splitter S (whole
input, parts outputs) joins each digit of the whole with the digit of $b$ followed by $a$ in
order.
\item \emph{Anchors.} Breadth-first search from each class variable that is its own class, in
increasing order: each unvisited hyperedge on the current variable is entered there; its other
variables are reached with that variable as parent. Every hyperedge must be entered
(\lexcode!a digit constraint not reached from its class's variable!). The anchor of a hyperedge is
the label owning its entry variable.
\item \emph{Fragments}, per node in sorted order (\lexcode!common! stands for the argument string
\lexcode!E σ ownL labsL digL hown hσ _ rfl!):
\begin{itemize}
\item B and S: node \lexlean!⟨.B, outs, ins, p, q⟩! ($p$, $q$ the digit counts of $a$, $b$),
fragment \lexlean!bFrag digL whole a b!, proof \lexlean!b_frag! or \lexlean!s_frag! with the decided
side conditions and $M$. The decided side conditions of the B, S, F, Z and X fragment lemmas
(\lexlean!b_frag!, \lexlean!s_frag!, \lexlean!f_frag_dig!, \lexlean!z_frag_dig!, \lexlean!x_frag_dig!) are
each \lexpy!DK!; those of \lexlean!za_frag! and \lexlean!xa_frag! are \lexlean!(by decide)!.
\item W: every port must have one digit (\lexcode!node <id> is a W on dimension 2^k, k != 1!);
node \lexlean!⟨.W, outs, ins, 1, 0⟩!, fragment \lexlean!wFrag digL outs in0!, proof
\lexlean!wnode_frag!.
\item F: node \lexlean!⟨.F, outs, ins, k, 0⟩!, fragment \lexlean!fFragD M 2^k k digL out in!, proof
\lexlean!f_frag_dig! with $M$, $2^{k}$, $M/2^{k}$; it contributes $k$ factors $1/\sqrt2$.
\item Z and X: the phase table over $2^{k}$ values (\lexpy!_turns_table!, or \lexpy!_weights! for a
weighted spider, which returns the table, the zero indicator list, one pair (amplitude,
indicator list) per magnitude ratio $r\neq 1$ in increasing $r$ (amplitude $\sqrt2$ for $r=2$,
else $\sqrt r$), and the base magnitude $\mathit{mag}\cdot\sqrt2^{\mathit{mag\_k}}$). A Z spider with ports reads its table at
its anchors' digits, X and legless Z at their own fresh digits. \lexpy!_table_poly! interpolates
the table into a phase polynomial in units of $1/M$; \lexpy!_table_terms! returns the same
polynomial as its list of (sorted monomial, nonzero coefficient) terms, and \lexpy!_table_poly! is
its \lexlean!Poly! text (an entry or coefficient off the $1/M$
grid raises \lexcode!a phase <t> off the 1/<M> grid!). A nonzero table or a weight appends the
vector \lexcode!((ts : List ℕ).getD v 0 : ℕ)! to \lexcode!phases! and takes its 1-based index as the
phase index, unless \lexcode!prefer! names an index whose vector is equal (and whose coefficient
slot is free or the node is weighted). Node \lexlean!⟨.Z, outs, ins, k, ph⟩! (\lexlean!.ZW!,
\lexlean!.XW! when weighted). Fragment \lexlean!zFragD digL anchors legs poly! with proof
\lexlean!z_frag_dig!, or \lexlean!xFragD M 2^k k digL poly ys outs ins! with proof \lexlean!x_frag_dig!
(contributing $k\cdot|\mathit{ports}|$ factors $1/\sqrt2$). A legless Z wraps the proof in
\lexlean!(val_z_legless E σ _ rfl rfl rfl).trans!.
\item Weights: the base coefficient $w$ is stored at the phase index in \lexcode!coefs! (a shared
index with another coefficient raises \lexcode!node <id> shares a phase index of another weight!;
no table raises \lexcode!node <id> carries weights and no coefficient table!). With a zero entry
or an amplitude pair and no \lexcode!amp!, raise \lexcode!node <id> has weights of different magnitudes!. With \lexcode!amp!, the node gets one atom per nonempty indicator (base 0 for the
zero list, then each pair's amplitude) on the indicator polynomial of its anchor variables,
the index table \lexcode!tbl! (per value $v$ the atom indices that hold at $v$) is stored in
\lexcode!tabs! (no \lexcode!tabs! raises \lexcode!node <id> carries atoms and no index table!; a shared
phase index with another table raises \lexcode!node <id> shares a phase index of other atoms!),
the fragment is scaled by $w$, and the proof is \lexlean!za_frag! (Z with ports) or
\lexlean!xa_frag!. A weight without atoms scales the fragment
and wraps the proof in \lexlean!w_frag! with \lexlean!val_zw! or \lexlean!val_xw!.
\item With \lexcode!amp!, nodes without atoms wrap the proof in \lexlean!atoms_nil A M _ _!, and
nodes that are Z with ports or atom-free in \lexlean!fvalA_of M A _ _ _ (by decide)!.
\end{itemize}
\item \emph{Pins.} The current path is every summed digit followed by the nodes' own digits.
Without \lexcode!amp!: for each reached variable $x \ge$ \lexpy!PIN_OFFSET! in reverse reach order
(deepest first), remove $x$ and emit \lexlean!.front x!, \lexlean!.pin [[parent(x)]]!; then
\lexlean!.norm!; then for each W pin $(x, F)$ emit \lexlean!.front x!, \lexlean!.pin F!. The
remaining path must be the raw path (\lexcode!the pinned network keeps other path variables!) and
is finally sorted by a \lexlean!.perm!. With \lexcode!amp!, the side keeps the pins and the W pins
and writes no steps; every digit is held initially.
\item The side's scalar is the diagram scalar as a constant model.
\end{enumerate}
\end{algorithm}

\lexpy!_table_modulus(diagram)! is the least common denominator, in turns, of the phase tables of
all Z and X spiders on dimensions $2^{k}$; \lexpy!net_body! takes $m_{\min}$ as its lcm with the
caller's value.

\paragraph{Owner lines (\texttt{\_own\_lines}).} For side tag $T$ (\lexcode!L! or \lexcode!R!):
\begin{lstlisting}[style=lexlean]
def ownT : Own := [(x, (lab, j)), ...]
def labsT : List ℕ := [labels]
def digTT : List (ℕ × List ℕ) := [(lab, [digit variables]), ...]
def digT (l : ℕ) : List ℕ := (digTT.lookup l).getD []
theorem ownTOK : ownOK ownT labsT digT = true :=
  ownOK_tab ownT labsT digTT (by decide +kernel)
\end{lstlisting}
The owner table lists, per label in order and per digit $j$, the digit variable $x$ with its
owner $(\mathit{lab}, j)$; \lexlean!digTT! stands for \lexlean!digTL! or \lexlean!digTR!, the
table behind \lexlean!digL! or \lexlean!digR!. \lexpy!_anet_lines! (atom networks) starts with the same lines.

\paragraph{Per-side lines (\texttt{\_net\_lines}).} After the owner lines: the network
\lexlean!def NT : Net := ⟨[nodes], summed, [0, ..., b-1]⟩!; per node $i$ (0-based) the pair
\lexlean!noncomputable def qTi : Node × Frag := (node, frag)! and its node fact
\begin{lstlisting}[style=lexlean]
theorem ndTi (σ : ℕ → ℕ) (hσ : (∀ l ∈ labsT, σ l < 2 ^ (digT l).length)) :
    ArchytasZX.val E σ qTi.1 = fval M qTi.2 (dbits ownT σ) := by
  have hown := ownTOK
  exact <fragment proof>
\end{lstlisting}
then the node names split into chunks of \lexpy!NODE_CHUNK! in order (one empty chunk when there
are no nodes); per chunk $j$ the list \lexlean!psTcj : List (Node × Frag)! (\lexlean!([] : List (Node × Frag))!
when empty) and the theorem
\begin{lstlisting}[style=lexlean]
theorem chTj : ∀ p ∈ psTcj, ∀ σ, (∀ l ∈ labsT, σ l < 2 ^ (digT l).length) →
    ArchytasZX.val E σ p.1 = fval M p.2 (dbits ownT σ) := by
  intro p hp
  simp only [psTcj, List.mem_cons, List.not_mem_nil, or_false] at hp
  rcases hp with rfl | ... | rfl
  exacts [ndTa, ..., ndTb]
\end{lstlisting}
(the last two lines only for a nonempty chunk); then \lexlean!psT := psTc0 ++ (psTc1 ++ (...))!,
right-nested, and
\begin{lstlisting}[style=lexlean]
theorem netT : NetPSD E NT M ownT labsT digT psT where
  ownok := ownTOK
  nodes := rfl
  node := by
    intro p hp
    simp only [psT, List.mem_append] at hp     -- with more than one chunk
    rcases hp with h | ... | h
    exacts [chT0 p h, ..., chTr p h]           -- one chunk: exact chT0 p hp
  disj := disjB_pairwise _ (by decide +kernel)
  dims := dims_tab E labsT digTT (by decide +kernel)
  cover := cover_of NT (by decide +kernel)
  sub := sub_of NT labsT (by decide +kernel)
theorem cT : (joinAll (psT.map Prod.snd)).c = w1 * ... * (rt 2)⁻¹ ^ j := by
  simp only [psT, psTc0, ..., qT0, ..., List.cons_append, List.nil_append, List.map_cons,
    List.map_nil, joinAll, Frag.join, Frag.unit, zFragD, xFragD, fFragD, bFrag, wFrag,
    Frag.scale, List.length_cons, List.length_nil]
  ring
\end{lstlisting}

\paragraph{The environment (\texttt{net\_env}).} Every network proof starts with
\begin{lstlisting}[style=lexlean]
def Δ (k : ℕ) : ℕ := 2 ^ k
noncomputable def Φ : ℕ → ℕ → ℝ
  | 1 => fun v => (([0, 1, 0, 3] : List ℕ).getD v 0 : ℕ) / ((8 : ℕ) : ℝ)
  | _ => fun v => ((([] : List ℕ).getD v 0 : ℕ) : ℝ) / ((M : ℕ) : ℝ)
def labs : List (ℕ × ℕ) := [(label, digit count), ...]
noncomputable def Ψ : ℕ → ℕ → ℂ          -- only when coefficients exist
  | i => fun v => w_i * ex (Φ i v) * (((tbl_i : List (List ℕ)).getD v []).map A).prod
  | _ => fun _ => 1
noncomputable def E : Env := { dimv := Δ, phase := Φ, labs := labs, coef := Ψ }
\end{lstlisting}
The atom factor of \lexlean!Ψ! is present only at indices of \lexcode!tabs!; without coefficients
\lexlean!E! omits \lexlean!coef!.

\begin{algorithm}[\texttt{net\_body}]\label{alg:lean-export-net-body}
\begin{enumerate}
\item Both raw sums (with W nodes) must be qubit sums with one boundary
(\lexcode!the boundaries differ!).
\item $m_{\min} \gets \mathrm{lcm}(m_{\min}, \text{both table moduli})$; plan with
\lexpy!support=True!. If the plan raises \lexpy!_LiftRaises!, take the parity plan
(\cref{alg:lean-export-parity-plan}) with this $m_{\min}$. If the plan, or that parity plan,
raises another \lexpy!PathSumUnsupported!, try the zero plan (below); if that is \lexpy!None!,
re-raise.
\item Build both sides with label bases $2^{16}$, $2^{17}$ (\lexcode!prefer! on the left only);
the boundary digit counts must agree (\lexcode!the boundary ports carry different dimensions!).
Append the plan's steps to each side's pin steps.
\item The constant turns are $k_L, k_R$ of \lexpy!_const_turn!; on the parity route $k_L$ is the
parity plan's $k$ and $k_R = 0$.
\item Prove $f_L s_L w_L (1/\sqrt2)^{j_L}\omega_M^{k_L} = f_R s_R w_R
(1/\sqrt2)^{j_R}\omega_M^{k_R}$ with $s$ the diagram scalars, $w$ the weight products and $j$
the counts of $1/\sqrt2$ factors. The statement of \lexlean!consts! omits the factor
\lexlean!om M ^ kR! on the parity route.
\item Emit both sides' lines, the step lists, \lexlean!facL!, \lexlean!facR!, \lexlean!consts!, two
\lexlean!run_chain! commands from \lexlean!(netDataD digT NT (psT.map Prod.snd))!, and the theorem
below, named \lexcode!name! (\lexcode!main! by default).
\end{enumerate}
\end{algorithm}
\begin{lstlisting}[style=lexlean]
theorem main : Equal E (sL) NL (sR) NR := by
  apply equal_psd2 _ _ netL netR rfl
  intro σ
  rw [cL, cR, ← val_smul, ← val_smul,
    run_sound M (by norm_num) stepsL (netDataD digL NL (psL.map Prod.snd)) _ _ runL_ok, runL_eq,
    run_sound M (by norm_num) stepsR (netDataD digR NR (psR.map Prod.snd)) _ _ runR_ok, runR_eq,
    val_agree M _ runL [0, ..., nb-1] _ _ (by decide +kernel) (by decide +kernel)
      (dbits_agree ownL ownR _ [0, ..., nb-1] (by decide +kernel))]
  apply final_sound M (by norm_num) _ _ _ _ _ kL kR (by decide +kernel) (by decide +kernel)
    (by decide +kernel)
  first | (rw [facL, facR]; exact consts) | rw [facL, facR]
\end{lstlisting}
On the parity route the \lexlean!apply! line is
\lexlean!apply finalP_sound M (by norm_num) _ _ [v1, ..., vr] _ _ _ kL (by decide +kernel) (by decide +kernel)!
with the parity plan's pivot order.
The returned \lexpy!NetBody! holds the lines, the labels' digit counts, both scalar texts and $M$
(fields \lexpy!head! and \lexpy!tabs! default empty).

\paragraph{Zero plans.} \lexpy!_zero_plan(left, right, m_min)! applies when each side either has
scalar zero (by \lexpy!sp.simplify!) or reduces to a sum containing the constraint $1=0$. For a
side of the second kind its steps are the trace steps followed by \lexlean!.norm!; $M$ is
$\mathrm{lcm}(4, m_{\min}, \text{phase denominators})$, with 8 for an \lexcode!omega! step. It returns
\lexpy!None! if a side has atoms, is not a qubit sum, or neither holds. \lexpy!_zero_body! emits both
sides' lines, for each reduced side a \lexlean!run_chain! from its network data with
\lexpy!RUN_BLOCK!, and the theorem \lexlean!Equal E (sL) NL (sR) NR! proved by \lexlean!equal_psd2!, the
rewrites \lexlean!run_sound!, \lexlean!run_eq! and \lexlean!PSem.val_of_zeroOK M _ runT _ (by decide +kernel)!
per reduced side, and \lexlean!simp!.

\paragraph{Networks with atoms.} \lexpy!net_amp_body! uses $M = $ \lexpy!amp_modulus! with the table
moduli, one \lexpy!Coefs(M)! (or the shared one; a modulus mismatch raises \lexcode!the shared coefficients are over zeta_<a>, not zeta_<M>!), sides built with \lexcode!amp!, and runs from
\lexpy!_net_run!: start from the held digits and the nodes' atoms, release, then per pin $(x, p)$:
when an atom reads $x$, hold $p$ if it is on the path (inner \lexlean!.front p!, then
\lexlean!.hold!), emit \lexlean!.asubB x [[p]]!, rename $x$ to $p$ in the atoms, release; then
inner \lexlean!.front x!, \lexlean!.pin [[p]]!; finally \lexlean!.norm! and the W pins through
\lexpy!AmpRun.pin!. \lexpy!amp_plan! completes the runs. Every decided side condition of the final
\lexlean!valA_agree! and the plan's final lemma is \lexlean!(by decide +kernel)!. The
lines use \lexlean!NetPSDA! (\lexpy!_anet_lines!, fragment triples \lexlean!(Node × Frag × Atoms)!),
\lexlean!netDataA!, \lexlean!aruns_chain!, \lexlean!equal_psda2!, \lexlean!valA_smul!, \lexlean!valA_agree! and
the plan's final lemma. \lexpy!_export_amp_net! wraps it under \lexlean!import ArchytasZX.PathSumNetAmp!
(and \lexlean!ArchytasZX.PathSumSupport! when \lexlean!afinalS_perm! occurs).

\paragraph{Networks through [Quarter].} \lexpy!net_quarter_body! applies to atom-free networks
whose reductions take [Quarter] or [Turn] steps: each side's pin steps run as plain
\lexlean!Step!s (\lexlean!stepsL!, chained with \lexpy!CHAIN_BLOCK!), and its reduction runs as
\lexlean!AStep!s from an empty-held, atom-free run (\lexlean!stepsAL!, chained from
\lexlean!⟨runL, [], []⟩!). The theorem rewrites by \lexlean!run_sound!, then \lexlean!← valA_nil M _ A!,
\lexlean!aruns_sound!, \lexlean!valA_agree! and closes with \lexlean!rw [facL, facR, facAL, facAR]!,
\lexlean!exact consts!. \lexpy!_export_quarter_net! imports \lexlean!ArchytasZX.PathSumNetAmp! and
\lexlean!ArchytasZX.PathSumChain!.

\begin{signature}
\lexpy!export_net(left: Diagram, right: Diagram, name: str = "main") -> str!
\end{signature}
Dispatch: atoms on either raw sum select the atom network proof; two qubit sums of which a
reduction is quartered select the quarter network proof; otherwise \lexpy!net_body! under
\lexlean!import ArchytasZX.PathSumNetW! (plus \lexlean!ArchytasZX.PathSumSupport! when
\lexlean!finalS_sound! occurs), namespace \lexlean!ArchytasZX.PathSumNetProof!, \lexpy!NET_OPTIONS!, and
\lexpy!net_env(phases, labs, "2 ^ k", M, coefs)!.

\subsection{Qudit network proofs}

\begin{algorithm}[\texttt{\_qnet\_side}]\label{alg:lean-export-qnet-side}
The diagram must give a qudit sum (\lexcode!the qudit network export covers odd prime path sums only!) without atoms (\lexcode!the qudit network export has no atoms!), every port of dimension
$p$ (\lexcode!a leg of dimension p^k!), and only Z, X, F nodes (\lexcode!node <id> is a <name>!),
none weighted (\lexcode!node <id> carries weights!) or legless.
\begin{enumerate}
\item Boundary port $i$ is label $i$. A wire's label is its class variable when that is a path
variable not yet taken, else the next label from \lexpy!QPIN_OFFSET!.
\item Anchors: breadth-first search over Z spiders from each self-class variable; every Z spider
must be entered (\lexcode!a Z spider not reached from its class's variable!). X spiders take the
last $|X|$ sorted path variables as fresh variables, in node order.
\item Fragments: F (not on a self-loop, else \lexcode!node <id> is an F on a self-loop!):
\lexlean!fFragQ p out in!, proof \lexlean!f_frag_q!, phase $+\,x_{\mathrm{out}}x_{\mathrm{in}}$, one
$1/\sqrt p$. Z: \lexlean!zFragQ p var legs poly!, proof \lexlean!z_frag_q!, constraints
$\mathit{var}-\ell$ per port $\ell$. X: \lexlean!xFragQ p var poly outs ins!, proof
\lexlean!x_frag_q!, phase $+\,\mathit{var}\cdot o$ per output and $-\,\mathit{var}\cdot i$ per
input, $|\mathit{ports}|$ factors $1/\sqrt p$. A nonzero turns table (entries on the $1/p$ grid,
else \lexcode!node <id> has a phase <t> off the 1/<p> grid!) appends a phase vector.
\item Pins: for each reached label $x \ge$ \lexpy!QPIN_OFFSET! deepest first, find a constraint
that sets $x$ to its parent $a$ (\lexcode!no constraint sets the label <x> to <a>!), emit
\lexlean!.perm [x, path]!, \lexlean!.pin a!, substitute. Then \lexlean!.norm!; the path must equal the
raw path. The pinned sum is reduced and its steps appended (\lexpy!_qudit_steps!).
\item The scalar must model without opaques (\lexcode!the scalar <s> has no exact Lean form!).
\end{enumerate}
\end{algorithm}
\lexpy!qnet_body! checks one prime and one boundary, one canonical form and one elimination order,
uses $n = 12$ for $p=3$ and $n=p$ otherwise, and emits \lexlean!QNetPS! facts (fields \lexlean!nodes!,
\lexlean!node!, \lexlean!disj!, \lexlean!dims!, \lexlean!cover!), \lexlean!cT : (qjoinAll ...).c = (rt p)⁻¹ ^ j!,
the \lexlean!lg! and \lexlean!qg! theorems, \lexlean!facL!, \lexlean!facR!, \lexlean!consts!,
\lexlean!qrun_chain! commands from \lexlean!(qnetData NT (psT.map Prod.snd))!, and the theorem
\lexlean!Equal E (sL) NL (sR) NR! by \lexlean!equal_qps!, \lexlean!QSem.val_smul!,
\lexlean!QSem.run_sound! and \lexlean!QSem.final_sound!. Every label has digit count 1.
\lexpy!export_qudit_net! wraps it under \lexlean!import ArchytasZX.PathSumQuditNet!, namespace
\lexlean!ArchytasZX.PathSumQuditNetProof!, \lexpy!QNET_OPTIONS!, the prime instance, and
\lexpy!net_env(phases, {}, str(p), p)!.

\subsection{Mixed network proofs \texttt{export\_mixed\_net}}\label{sec:lean-export-mixed}

\begin{signature}
\lexpy!export_mixed_net(left: Diagram, right: Diagram, name: str = "main") -> str!\\
\lexpy!mixed_net_body(left: Diagram, right: Diagram, name: str, phases: list[tuple[str, int]]) -> tuple[NetBody, Coefs]!
\end{signature}
\lexpy!export_mixed_net! proves \lexlean!Equal E sL NL sR NR! for a left diagram on dimensions 2, 3
and 6 (\lexpy!is_mixed!: no \bboxes, every port of dimension 1, 2, 3 or 6, with a factor 2 and a
factor 3 present) against a qubit right diagram. A left side that is not mixed raises
\lexcode!the mixed network export takes a mixed left side!. The left network's value is its
$\Z_3$ and $\mathbb{F}_2$ fragments' mixed path sum (\lexlean!NetPSM.eval_mval!,
\file{ArchytasZX/PathSumMixedNet.lean}); the $\Z_3$ replay \lexlean!z3L! leaves a qubit path sum
with zero atoms (\cref{sec:lean-export-amp-proof}); the $\mathbb{F}_2$ pins and both sides'
reduction steps follow as \lexlean!AStep! runs, and \lexlean!valA_agree! with \lexlean!mbits_dbits! and
the plan's final lemma give \lexlean!Equal!.

A label $\ell$ of the left network is its own $\mathbb{F}_2$ and $\Z_3$ variable. Its value
$j$ on dimension 6 has the bit $j \bmod 2$ and the trit $\lfloor j/2\rfloor$ (\lexlean!mbit!,
\lexlean!mtrit!); on dimension 2 only the bit, on dimension 3 only the trit. An \term{indicator
entry} $(u, z, a, b, k)$ (Lean \lexlean!Ent!) states that the $\mathbb{F}_2$ variable $u$ is
$[\,a\rho(z)+b = k\,]$ for the $\Z_3$ assignment $\rho$.

\begin{signature}
\lexpy!@dataclass class _MixSide!: \lexpy!nodes: list[str]!, \lexpy!frags3: list[str]! ($\Z_3$
fragments), \lexpy!frags2: list[str]! ($\mathbb{F}_2$ fragments), \lexpy!proofs: list[str]!,
\lexpy!dims: dict[int, int]! (label to dimension), \lexpy!summed: list[int]!, \lexpy!boundary: int!,
\lexpy!ents: list[tuple[int, int, int, int, int]]!, \lexpy!inv2: int!, \lexpy!inv3: int! (powers of
$1/\sqrt2$ and $1/\sqrt3$), and the side's sums \lexpy!path3: list[int]!, \lexpy!phase3: ModPoly!
(over 3), \lexpy!cons3: list[ModPoly]!, \lexpy!path2: list[int]!, \lexpy!phase2: PhasePoly!,
\lexpy!cons2: list[F2Poly]!, each defaulting to empty.
\end{signature}

\lexpy!_mixed_modulus(diagram)! is the least common multiple of 2 and the denominators of
$t \bmod 1$ over the phase tables (\lexpy!_turns_table!, in turns) of the Z and X spiders with
ports whose first port has dimension 2 or 6.

\begin{algorithm}[\texttt{\_mixed\_net\_side(diagram, tag, label0, M, phases)}]\label{alg:lean-export-mixed-side}
\begin{enumerate}
\item A diagram that is not \lexpy!is_mixed! raises
\lexcode!the mixed network export covers diagrams on 2, 3 and 6!.
\item \emph{Labels.} Boundary port $i$ (outputs, then inputs) is label $i$ of dimension 2; a
boundary port of another dimension raises
\lexcode!a boundary leg of the mixed diagram is not of dimension 2!. Wire $w$ in \lexpy!sort_key!
order is label $\mathit{label0}+w$ with its ports' dimension, summed; a wire on a boundary port
raises \lexcode!a wire on a boundary port!, a wire whose ends differ in dimension
\lexcode!a wire joins two dimensions!, and at least \lexpy!MIX_FRESH! wires
\lexcode!too many wires for the mixed label offsets!, as does a boundary of more than
$\mathit{label0}$ ports. Fresh X variables start at $\mathit{label0}+$\lexpy!MIX_FRESH!, entry
variables at $\mathit{label0}+$\lexpy!MIX_ENT!; after the nodes, X variables past
$\mathit{label0}+$\lexpy!MIX_ENT! raise \lexcode!too many X spiders for the mixed label offsets!
and entry variables past $\mathit{label0}+$\lexpy!MIX_SPAN! raise
\lexcode!too many indicator entries for the mixed label offsets!.
\lexcode!path3! starts as the summed labels of dimension 3 or 6.
\item \emph{Nodes} in sorted order. A node without ports raises \lexcode!node <id> has no legs!; a
weighted phase raises \lexcode!node <id> carries weights!. Let $k$ be the dimension of the first
port and \lexcode!ls! the output labels followed by the input labels.
\begin{itemize}
\item B or S: the whole port (B's output, S's input) on 6 with parts $a$ on 3 and $b$ on 2,
else \lexcode!node <id> is not a (3, 2) <name> on 6!. Constraints $\mathit{whole}-a$ over
$\Z_3$ and $\mathit{whole}=b$ over $\mathbb{F}_2$. Node \lexlean!⟨.B, outs, ins, 3, 2⟩! (or
\lexlean!.S!), fragments \lexlean!bFrag3 whole a! and \lexlean!bFrag2 whole b!, proof \lexlean!node_b! with
\lexlean!val_binder! or \lexlean!val_splitter!.
\item Any other node with ports of different dimensions raises
\lexcode!node <id> joins legs of different dimensions!; a generator other than Z, X, F raises
\lexcode!node <id> is a <name>!.
\item $k=6$: only Z (\lexcode!node <id> is a <name> on dimension 6!). With the six-entry table
$t$, extend it to eight entries
$t' = (t_0,\ldots,t_5,\ t_2+t_4-t_0,\ t_3+t_5-t_1)$. With $\ell_0$ the first label and fresh
entry variables $u_1, u_2$, add the entries $(u_1, \ell_0, 1, 0, 1)$ and
$(u_2, \ell_0, 1, 0, 2)$, the $\Z_3$ constraints $\ell_0 - \ell$ and the $\mathbb{F}_2$
equalities $\ell_0 = \ell$ for every label $\ell$; interpolate $t'$ over
$[\ell_0, u_1, u_2]$ (\lexpy!_table_terms!) into the $\mathbb{F}_2$ phase. A nonzero table over
the first six entries appends a phase vector over $M$. Node \lexlean!⟨.Z, outs, ins, 6, ph⟩!,
fragments \lexlean!z6Frag3 l0 ls! and \lexlean!z6Frag2 l0 ls poly!, proof \lexlean!node_z6!.
\item $k=3$, F: a self-loop raises \lexcode!node <id> is an F on a self-loop!; phase
$+\,x_{\mathrm{out}}x_{\mathrm{in}}$, fragment \lexlean!QSem.fFragQ 3 out in!, one $1/\sqrt3$.
\item $k=3$, Z or X: the variable is the first label (Z) or the next fresh variable (X, added
to \lexcode!path3!). The phase is \lexpy!_qudit_table! over that variable; Z adds the constraints
$\mathit{var}-\ell$, X the phase $+\,\mathit{var}\cdot o$ per output and
$-\,\mathit{var}\cdot i$ per input and $|\mathit{ls}|$ factors $1/\sqrt3$. The turns table on
the $1/3$ grid (else \lexcode!node <id> has a phase <t> off the 1/3 grid!) always appends a phase
vector over 3. Fragment \lexlean!QSem.zFragQ 3 var ls poly! or
\lexlean!QSem.xFragQ 3 var poly outs ins!; the $\mathbb{F}_2$ fragment is \lexlean!Frag.unit!; proof
\lexlean!node_q3! around \lexlean!QSem.z_frag_q! or \lexlean!QSem.x_frag_q! (F: \lexlean!QSem.f_frag_q!).
\item $k=2$: F gives \lexlean!fFragD M 2 1 digT out in!, phase $+\frac12 x_{\mathrm{out}}x_{\mathrm{in}}$
and one $1/\sqrt2$. Z reads its table at its first label and adds the equalities
$\ell_0=\ell$; X reads it at a fresh variable $y$ (added to \lexcode!path2!), adds the phase
$+\frac12\ell y$ per label and $|\mathit{ls}|$ factors $1/\sqrt2$. A nonzero table appends a
phase vector over $M$. Fragments \lexlean!zFragD! or \lexlean!xFragD! as in
\cref{alg:lean-export-net-side}; the $\Z_3$ fragment is \lexlean!QSem.QFrag.unit!; proof
\lexlean!node_q2! around \lexlean!fvalA_of! and \lexlean!atoms_nil! of \lexlean!z_frag_dig! or
\lexlean!x_frag_dig!.
\end{itemize}
\item The side keeps the nonzero $\Z_3$ constraints.
\end{enumerate}
\end{algorithm}

\paragraph{Per-side lines (\texttt{\_mix\_lines}).} With fragment type
\lexlean!(Node × QSem.QFrag) × Frag × Atoms! and the qubit labels \lexcode!qlabs! (labels of
dimension 2, sorted): \lexlean!msT : List MLab! (every label $\ell$ of dimension $k$ as
$(\ell, k, \ell, \ell)$), \lexlean!miT! (the summed labels likewise), \lexlean!esT : List Ent!,
\lexlean!ownT! ($(\ell, (\ell, 0))$ per qubit label), \lexlean!qlabsT!, \lexlean!digTT! and \lexlean!digT!
($[\ell]$ per qubit label); the facts \lexlean!msTOK : mOK msT = true!, \lexlean!ownTOK! by
\lexlean!ownOK_tab!, \lexlean!qlT! (each qubit label is in \lexlean!msT! with dimension 2, has digits
$[\ell]$ and is no entry variable) and \lexlean!esTnd : (esT.map (·.1)).Nodup!, all by
\lexlean!decide +kernel!; the network \lexlean!NT!; per node $i$ the pair \lexlean!pTi! with empty atoms
and its fact \lexlean!ndTi! under \lexlean!MIn msT σ!,
\begin{lstlisting}[style=lexlean]
ArchytasZX.val E σ pTi.1.1 = QSem.fval 3 pTi.1.2 (mtrits msT σ) *
  fvalA M A pTi.2.1 pTi.2.2 (liftA esT (mtrits msT σ) (mbits msT σ))
\end{lstlisting}
chunk lists and theorems \lexlean!psTcj!, \lexlean!chTj! of \lexpy!NODE_CHUNK! nodes as in
\lexpy!_net_lines!, \lexlean!psT!, and the theorem \lexlean!netT : NetPSM E NT M A msT miT esT psT! with
fields \lexlean!ok!, \lexlean!nodes := rfl!, \lexlean!node!, \lexlean!disj3 := disj3_of _ (by decide +kernel)!,
\lexlean!disj2 := disj2_of _ (by decide +kernel)!, \lexlean!dims!, \lexlean!sub!, \lexlean!bnd! by
\lexlean!decide +kernel!, \lexlean!int := rfl! and \lexlean!cover := cover_of NT (by decide +kernel)!.
\lexpy!_mix_c_lines! adds \lexlean!cT : (joinAll (psT.map (·.2.1))).c = (rt 2)⁻¹ ^ inv2! and
\lexlean!c3T : (QSem.qjoinAll (psT.map (·.1.2))).c = (rt 3)⁻¹ ^ inv3!, each by \lexlean!simp only!
over the definitions and \lexlean!ring!.

\begin{signature}
\lexpy!@dataclass class _Z3Run!: \lexpy!defs: list[str]!, \lexpy!thm: list[str]!, \lexpy!q: str! (the name of
the final \lexlean!AData!), \lexpy!fac: int! (the power of 3 gained),
\lexpy!subs: list[tuple[int, int]]! (duplicate entries substituted), \lexpy!pairs: list[tuple[int, int]]!
(the entry pairs summed); the last two default to empty.
\end{signature}

\begin{algorithm}[\texttt{\_z3\_lines(tag, side, M)}: the $\Z_3$ replay]\label{alg:lean-export-z3}
The state is a qubit \lexlean!AData! \lexcode!q!, a \lexlean!QSem.Data! \lexcode!t! and an entry list
\lexcode!e!, starting at \lexlean!qT0!, \lexlean!tT0!, \lexlean!esT!; every change defines a new numbered
state (\lexlean!tT1!, \lexlean!eT2!, \ldots) and records one step lemma. A \emph{perm} to a path
order records \lexlean!mval_perm! and a state with that path.
\begin{enumerate}
\item \emph{Pins.} Repeat until no change: take the first $\Z_3$ constraint with exactly two
variables of which one, $x$, is on the path (the first such); let $y$ be the other. Perm $x$ to
the head if needed, record \lexlean!mval_pin!, set \lexlean!t := QSem.pinS 3 y t! and
\lexlean!e := e.map (pinE x y 1 0)!, drop the constraint, substitute $x \mapsto y$ in the others,
the phase and the entries. A constraint left raises
\lexcode!a Z_3 constraint of the mixed side is not a pin!.
\item Record \lexlean!mval_norm!; \lexlean!t := QSem.normS 3 t!.
\item \emph{[HH] pivots.} Repeat until no change: take the first path variable $y$ read by no
entry, of degree 1 in the phase, whose coefficient is linear without constant and without
powers, and the least variable $z$ of that coefficient read by no entry and on the path. Perm
to $[y, z, \mathit{rest}]$, record \lexlean!mval_hhPivot! with factor 3, substitute
$z \mapsto -c_z^{-1}(\mathit{coefficient} - c_z z)$ and drop $y$; \lexlean!t := QSem.hhPivotS 3 t!.
\item \emph{Duplicates.} For each entry whose $(z, a, b, k)$ equals an earlier entry's $e_1$:
record \lexlean!mval_dup!, set \lexlean!q := substA M u [[u1]] q! and erase the entry; record
$(u, u_1)$ in \lexcode!subs!.
\item \emph{[Count3].} For each remaining read variable $s$ in increasing order: its entries
sorted by $k$ must be exactly $(\cdot, s, 1, 0, 1)$, $(\cdot, s, 1, 0, 2)$ with $s$ on the
path, else \lexcode!a Z_3 variable of the mixed side is not one entry pair!. Perm $s$ to the
head, record \lexlean!mval_count3!, and set \lexlean!q := ⟨q.s, u1 :: u2 :: q.held, (0, [[u1, u2]]) :: q.atoms⟩!;
record $(u_1, u_2)$ in \lexcode!pairs!.
\item A path or phase left raises
\lexcode!the Z_3 part of the mixed side does not reduce to a constant!.
\end{enumerate}
\end{algorithm}
The emitted theorem rewrites each recorded step by \lexlean!rw [show <before> = <after> from <lemma>]!:
\begin{lstlisting}[style=lexlean]
theorem z3L (c : ℂ) (τ : ℕ → ZMod 2) (ρ : ℕ → ZMod 3) :
    mval M c A qL0 tL0 esL τ ρ =
      3 ^ fac * valA M c A qLn τ := by
  rw [show mval M c A qL0 tL0 esL τ ρ = mval M c A qL0 tL1 esL τ ρ from mval_perm ...]
  ...
  rw [show eLn = [] from rfl,
    mval_end M (by norm_num) c A qLn tLn τ ρ (by decide +kernel)
      (by decide +kernel) (by decide +kernel) (by decide +kernel)]
  ring
\end{lstlisting}
with the initial states
\lexlean!qT0 := ⟨⟨[], J.phase, J.cons, []⟩, bvars miT ++ J.path, (psT.map (·.2.2)).flatten⟩! for
\lexlean!J = joinAll (psT.map (·.2.1))! and
\lexlean!tT0 := ⟨tvars miT ++ K.path, K.phase, K.cons, []⟩! for
\lexlean!K = QSem.qjoinAll (psT.map (·.1.2))!.

\begin{algorithm}[\texttt{\_mix\_qubit(side, z3, amp, boundary)}]\label{alg:lean-export-mix-qubit}
Returns the qubit \lexpy!PathSum! of the left side after its $\Z_3$ replay and its
$\mathbb{F}_2$ pins, and its \lexpy!AmpRun!.
\begin{enumerate}
\item Held variables: the summed labels not of dimension 3, then \lexcode!path2!. In the
$\mathbb{F}_2$ phase and constraints substitute each \lexcode!subs! pair $(u, u_1)$. For each
entry pair $(u_1, u_2)$ in order, prepend $u_1, u_2$ to the held variables and the atom
$(0, u_1u_2)$ to the atoms (the zero coefficient's index).
\item Join the two variables of each equality constraint $x = y$ by an edge. The start order is
the boundary labels on an edge, then every variable on an edge in increasing order. For each
start variable $v_0$ not yet seen, take its component; the roots are the component's variables
that are not held, in start order (or $v_0$ alone when there are none); search breadth-first
from the roots, neighbors in increasing order, recording each newly reached
$(x, \mathit{parent})$.
\item Start an \lexpy!AmpRun! with the held variables and atoms, released. For each reached
$(x, p)$ deepest first: $x$ must be held and read by no atom, else
\lexcode!the F_2 variable <x> of the mixed side is not a free pin!; substitute $x \mapsto p$ and
emit inner \lexlean!.front x!, \lexlean!.pin [[p]]!. Then inner \lexlean!.norm!.
\item Return the path sum over the boundary with the remaining held variables, the phase,
constant 1, the distinct nonzero constraints and one zero-base amplitude atom per pair.
\end{enumerate}
\end{algorithm}

\begin{algorithm}[\texttt{mixed\_net\_body}]\label{alg:lean-export-mixed-body}
\begin{enumerate}
\item The right raw sum (\lexpy!from_diagram(right, w_nodes=True)!) must be a qubit sum, else
\lexcode!the right side of the mixed export is not a qubit path sum!.
\item $M = \mathrm{lcm}(8, \text{\lexpy!_mixed_modulus(left)!}, \text{\lexpy!_table_modulus(right)!})$.
Build the left side with label base $2^{16}$ and its $\Z_3$ replay; \lexpy!Coefs(M)! must place
the zero coefficient at index 0 (\lexcode!the zero coefficient is not at index 0!).
\item The right boundary must be $(0, \ldots, b-1)$ for the left boundary size $b$
(\lexcode!the right side's boundary digits are not its port labels!). Build the left qubit sum
(\cref{alg:lean-export-mix-qubit}) and the right side by \lexpy!_net_side! with label base
$2^{17}$ and \lexcode!amp!; a different boundary size or a boundary label with digits other than
$[i]$ raises \lexcode!the right side's boundary labels are not one bit each!.
\item Run the right side by \lexpy!_net_run!; plan both runs by \lexpy!amp_plan!
(\cref{alg:lean-export-amp-plan}), giving $k_L, k_R$.
\item Prove $s_L (1/\sqrt3)^{\mathit{inv3}}\,3^{\mathit{fac}} f_L (1/\sqrt2)^{\mathit{inv2}}
\omega_M^{k_L} = f_R s_R w_R (1/\sqrt2)^{j_R}\omega_M^{k_R}$, with $s_L$ the left diagram
scalar as a constant model and $f$ the \lexpy!afac_lines! factors.
\item Emit the left lines, its constants, the $\Z_3$ replay, the right \lexpy!_anet_lines!, both
\lexlean!AStep! lists, the factor, coefficient and root identities, the theorem \lexlean!key!, two
\lexlean!aruns_chain! commands from the replay's final state and
\lexlean!(netDataA digR NR (psR.map (·.2.1)) (psR.map (·.2.2)))!, and the theorems below.
\item Return the \lexpy!NetBody! (labels: the left dimensions and the right labels as $2^{k}$) and
the \lexpy!Coefs!.
\end{enumerate}
\end{algorithm}
\begin{lstlisting}[style=lexlean]
theorem evalL (σ : ℕ → ℕ) (hσ : InRange E NL.bnd σ) :
    eval E NL σ = (QSem.qjoinAll (psL.map (·.1.2))).c *
      (3 ^ fac * valA M (joinAll (psL.map (·.2.1))).c A qLn
        (mbits msL (clamp NL.bnd σ))) := by
  rw [netL.eval_mval (by decide +kernel) (by decide +kernel) σ hσ]
  exact congrArg _ (z3L _ _ _)

theorem main : Equal E (sL) NL (sR) NR := by
  refine ⟨List.Perm.refl _, fun σ hσ => ?_⟩
  have hσR : InRange E NR.bnd σ := hσ
  rw [evalL σ hσ, netR.eval_eq σ hσR, show clamp NL.bnd σ = clamp NR.bnd σ from rfl,
    c3L, cL, cR, aruns_sound M (by norm_num) A stepsL qLn _ _ reqsL qreqsL
    arunsL_ok, arunsL_eq, aruns_sound M (by norm_num) A stepsR tR _ _ reqsR
    qreqsR arunsR_ok, arunsR_eq, valA_agree M _ A arunsL [0, ..., b-1] _ _ (by decide +kernel)
      (by decide +kernel) (by decide +kernel) (by decide +kernel)
      (mbits_dbits msL ownR labsR digR ownROK (by decide +kernel) _ [0, ..., b-1]
        (by decide +kernel))]
  simp only [← valA_smul]
  apply afinal_perm M (by norm_num) _ _ A _ _ _ kL kR
    (by decide +kernel) (by decide +kernel) (by decide +kernel) (by decide +kernel)
    (by decide +kernel) (by decide +kernel)
  rw [facL, facR]
  linear_combination key
\end{lstlisting}
\lexlean!afinalS_perm! replaces \lexlean!afinal_perm! when the plan leaves a support.
\lexpy!export_mixed_net! wraps the body under \lexlean!import ArchytasZX.PathSumMixedNet!,
\lexlean!import ArchytasZX.PathSumChain! (and \lexlean!ArchytasZX.PathSumSupport! when
\lexlean!afinalS_perm! occurs), namespace \lexlean!ArchytasZX.PathSumMixedProof!,
\lexlean!open ArchytasZX ArchytasZX.PSem ArchytasZX.MSem!, \lexpy!QNET_OPTIONS!, the coefficient
vector \lexlean!A! (\lexpy!Coefs.lean_def!), and \lexpy!net_env(phases, labs, MIX_DIMS, M)!.

% ---------------------------------------------------------------------------
\section{Amplitude-Atom Step Writer}\label{sec:lean-export-amp}

The module \module{archytaszx.semantics.lean_amp} follows the reduction trace of a qubit path
sum with amplitude atoms and writes the \lexlean!AStep! list of
\file{ArchytasZX/PathSumAmp.lean} that replays it. The Lean state \lexlean!AData! is a triple
$\langle s, H, A\rangle$: an inner \lexlean!Data! $s$ over the \emph{inner path}, a list $H$ of
\emph{held} variables (summed outside the inner sum), and a list of atoms $(j, F)$ standing for
$A_j^{F}$, where $A$ is the coefficient vector of the file and $F$ an $\mathbb{F}_2$
polynomial. An atom $(j,F)$ contributes $A_j$ where $F=1$ and $1$ where $F=0$.

\begin{signature}
\lexpy!Atom = tuple[int, F2Poly]!: a coefficient index and an $\mathbb{F}_2$ polynomial.\\
\lexpy!Req = tuple[int, int, tuple[tuple[int, ...], ...]]!: $(k, j, ps)$ asserts
$k\cdot A_j = \sum_{P\in ps}\prod_{i\in P} A_i$.\\
\lexpy!f2_text(poly)!: the polynomial as a Lean list of sorted monomials.\\
\lexpy!atoms_text(atoms)!: \lexcode![(j, [monos]), ...]!.
\end{signature}

\subsection{Coefficient bases}

\begin{signature}
\lexpy!@dataclass class Coefs!: \lexpy!m: int!, \lexpy!values: list[Cyclo]!.\\
\lexpy!index(c) -> int!: the index of the first equal base, appending \lexpy!c! when new.\\
\lexpy!model(i) -> sp.Expr!: base $i$ over \lexpy!Om(m)!; raises
\lexcode!PathSumUnsupported("an atom base <c> off zeta_<m>")! when the base's order does not
divide \lexpy!m!.\\
\lexpy!lean_def(sc) -> list[str]!: the vector
\lexlean!noncomputable def A : ℕ → ℂ | 0 => ... | _ => 0!.
\end{signature}

\subsection{The run}

\begin{signature}
\lexpy!@dataclass class AmpRun!: \lexpy!coefs!, \lexpy!cur: list[int]! (inner path), \lexpy!held: list[int]!,
\lexpy!atoms: list[Atom]!, \lexpy!steps: list[str]!, \lexpy!reqs: list[Req]!,
\lexpy!facs: list[sp.Expr]! (step factors), \lexpy!omega: bool! (an omega step occurred),
\lexpy!qreqs: list[tuple[int,int,int]]!, \lexpy!support: F2Poly | None!.
\end{signature}
\lexpy!snapshot()! copies the state (sharing \lexpy!coefs!); \lexpy!restore(other)! takes it back.
\lexpy!avars()! is the set of variables the atoms read. \lexpy!initial(coefs, path, amps)! returns
the run of a raw sum: held are the path variables an atom reads (sorted), inner the others.

\paragraph{Elementary operations.}
\begin{itemize}
\item \lexpy!inner(texts)!: each \lexlean!PSem.Step! text $t$ becomes \lexlean!.inner (t)!; factor 2
for \lexlean!.elim!, \lexlean!.hhCons!, \lexlean!.hhPivot!, \lexlean!.hhNl!; $\sqrt2\,\omega_m^{m/8}$ for
\lexlean!.omega! and $\sqrt2\,\omega_m^{7m/8}$ for \lexlean!.omegaConj! (both set \lexpy!omega!).
\item \lexpy!trace(step)!: \lexpy!trace_step! (\cref{sec:lean-export-trace-step}) on the inner path.
\item \lexpy!release()!: every held variable no atom reads is moved to the front of the held list
(\lexlean!.hperm [y, rest]! when not already first) and released (\lexlean!.release!) onto the front
of the inner path.
\item \lexpy!_pick(atom)!: moves an atom to the head of the list, emitting \lexlean!.apick i! when its
position $i$ is not 0.
\item \lexpy!hold(vs)!: each inner variable of \lexcode!vs!, lowest first: \lexlean!.inner (.perm [v, rest])!,
\lexlean!.hold!.
\item \lexpy!fold(path, unit=False)!: if any atom has polynomial 0, emit \lexlean!.adrop! and remove
those atoms; then, in list order, an atom on $1$ is picked and folded into the constant
(\lexlean!.aconst j!, factor $A_j$); with \lexcode!unit!, an atom of base 1 is dropped
(\lexlean!.aone j!, requirement $(1, j, ((),))$); an atom of base 0 reading no variable of
\lexcode!path! becomes a constraint (\lexlean!.azero j!, requirement $(1, j, ())$).
\end{itemize}

\paragraph{Reduction rules on held variables.}
\begin{itemize}
\item \lexpy!average(y, g)! ([Average] with guard $g$): every atom reading $y$ must be on
$(1+y)g$ or on $yg$ (\lexcode!an atom on y is not on (1 + y) g or y g!). With $\alpha$ the indices
on $(1+y)g$ and $\beta$ the others, $a = \prod_\alpha A$, $b=\prod_\beta A$, the new base
$(a+b)/2$ gets index $j$; emit \lexlean!.average y g j α β!, requirement $(2, j, (\alpha,\beta))$,
factor 2. The atoms reading $y$ are replaced by $(j, g)$, and $y$ leaves the held list.
\item \lexpy!average_half(y, q)! ([Average] through the half phase $yq/2$): atoms reading $y$ must
be on $1+y$ or $y$. With $j_0$ for $(a+b)/2$, $j_1$ for $(a-b)/2$ and $j_n$ for $-1$: emit
\lexlean!.inner (.norm)!, \lexlean!.averageH y j0 j1 jn α β!, requirements $(2, j_0, (\alpha,\beta))$,
$(2, j_1, (\alpha, \beta\,j_n))$, $(0, j_n, ((), (j_n)))$, factor 2; the new atoms are
$(j_0, q+1)$ and $(j_1, q)$.
\item \lexpy!average_ctx(y, q, ctx)! ([Average] on a context $c_1..c_r$ of $y$): per atom $(j,F)$
reading $y$, with $F_0 = F|_{y=0}$ and $F_1 = F|_{y=1}$, the lists $z_0$, $z_1$ of context
values $n<2^r$ where $F_0$ resp.\ $F_1$ is 1 (\lexpy!ctx_value!; an undefined value raises
\lexcode!an atom on y off its context!); $Q$ the values where $q$ is 1. Per $n$, $p_0$ = atoms with
$n\in z_0$, $p_1$ = atoms with $n\in z_1$ plus $j_n$ when $n\in Q$, and a new index $j_n'$ for
$(\prod_{p_0} + \prod_{p_1})/2$. Emit \lexlean!.inner (.norm)!,
\lexlean!.avctx y [ctx] [(j, z0, z1)...] Q [j'...] jn!, the requirements $(0, j_n, ((),(j_n)))$ and
$(2, j_n', (p_0, p_1))$ per $n$, factor 2; the new atoms are $(j_n', [\mathit{ctx}=n])$ by
\lexpy!turn_indicator!.
\end{itemize}

\paragraph{Rules on inner variables.}
\begin{itemize}
\item \lexpy!quarter(y, b0, b1)! ([Quarter]): requires $4 \mid m$
(\lexcode![Quarter] over zeta_<m>!); emit \lexlean!.inner (.perm [y, rest])!,
\lexlean!.quarter b0 b1 j! with $j$ the index of $(1+i)/2$ (\lexpy!QUARTER_BASE!), root requirement
$(m/4, j, 1)$, factor 2; append the atom $(j, b_0)$.
\item \lexpy!turn(y, bs)! ([Turn] with $k=|bs|$ digits): requires $2^{k}\mid m$; per
$s\in$ \lexpy!turn_values(k)! ($1\le s<2^k$, $s\ne 2^{k-1}$) the index $j_s$ of
$(1+e^{2\pi i s/2^k})/2$; emit \lexlean!.turn [bs] [(s, j_s)...]!, root requirements
$(s\,m/2^{k}, j_s, 1)$, factor 2; append atoms $(j_s, [\mathit{bs}=s])$.
\item \lexpy!turn_held(y, bs)!: hold the variables of \lexcode!bs!, apply \lexpy!turn!; when the
indicator of $s=2^{k-1}$ reads a live variable, prepend the zero atom $(j_0, \cdot)$ with
\lexlean!.acons j0! and requirement $(1, j_0, ())$.
\item \lexpy!pin(y, value)! ([Pin] along $y+\mathit{value}$): emit \lexlean!.inner (.norm)!; if an atom
reads $y$, hold the inner variables of \lexcode!value!, emit \lexlean!.asubB y value!, substitute in
the atoms and release; then \lexlean!.inner (.perm [y, rest])!, \lexlean!.inner (.pin value)!.
\end{itemize}

\paragraph{Final normalization.}
\begin{itemize}
\item \lexpy!substitute(rows)!: for each row (pivot, value) whose pivot an atom reads, emit
\lexlean!.asub pivot! and substitute; with a support $g$, drop every atom monomial whose product
with $g$ is 0 (\lexlean!.adropC!, only when something changes); emit \lexlean!.anorm! when atoms
remain; fold with \lexcode!unit!; merge atoms on equal polynomials (pick both, emit
\lexlean!.amerge j j1 j2! with $j$ the index of the product, requirement $(1, j, ((j_1, j_2)))$,
fold again) until no two share a polynomial; then \lexpy!complement()!.
\item \lexpy!complement()!: two atoms of one index whose polynomials add to $1$ are picked and
folded into the constant (\lexlean!.acomp j!, factor $A_j$), until none is left.
\item \lexpy!split()!: with a support $g$, every atom with at least two monomials whose pairwise
products vanish on $g$ is split into one atom per monomial (\lexlean!.asplit!); then
\lexpy!substitute([])!.
\item \lexpy!normalize(phase, support, split=False)!: with \lexpy!support_lift(support)! $=(a,b)$
emit \lexlean!.psup a b! and replace every phase monomial containing $a$ by its union with $b$,
then \lexlean!.inner (.norm)!; move every nonconstant phase term $c\cdot\mathit{mono}$ to an atom
(\lexlean!.aphase mono e j!, $e = cm \bmod m$, $j$ the index of $e^{2\pi i c}$, root requirement
$(e, j, 0)$); rewrite every atom $(j, F)$ with a constant term, $F \neq 1$, as
$A_j\cdot(A_j^{-1})^{F+1}$ (\lexlean!.aflip j inv one!, requirements $(1, \mathit{one}, ((),))$ and
$(1, \mathit{one}, ((j,\mathit{inv})))$, factor $A_j$); then \lexpy!substitute([])! and, with
\lexcode!split!, \lexpy!split()!.
\item \lexpy!final_key()!: the atoms sorted by index and sorted monomials.
\end{itemize}

\begin{algorithm}[\texttt{reduce\_run(run, trace)}]\label{alg:lean-export-reduce-run}
Fold on the live variables (inner and held) and release. For each trace step: \lexcode!average!
(its \lexpy!value! is the guard; missing raises \lexcode!an [Average] step without its guard!) by
\lexpy!average!; \lexcode!average_half! and \lexcode!pin! (missing value raises \lexcode!a [<rule>] step without its polynomial!) by \lexpy!average_half! or \lexpy!pin!; \lexcode!quarter! with
\lexpy!values = (b0, b1)!; \lexcode!turn!; \lexcode!turn_held!; \lexcode!average_ctx! with
\lexpy!value! (or 0) and \lexpy!values! as context; each followed by a fold on the live variables;
any other rule through \lexpy!trace!. Release after every step. Held variables left at the end
raise \lexcode!the held variables <H> remain!.
\end{algorithm}

\subsection{Requirement theorems}

\lexpy!reqs_lines(tag, reqs, coefs, sc, steps)! writes, per requirement $i$,
\begin{lstlisting}[style=lexlean]
theorem reqL0 : reqHolds A (2, 5, [[1, 2], [3]]) := by
  have key : (<k * A_j>  : ℂ) = <Σ Π A> :=
    by ...
  simp only [reqHolds, A, List.map_cons, List.map_nil, List.prod_cons,
    List.prod_nil, List.sum_cons, List.sum_nil]
  all_goals push_cast
  all_goals linear_combination key
\end{lstlisting}
and the collector \lexlean!reqsL : ∀ r ∈ areqs stepsL, reqHolds A r!, which states the decided list
\lexlean!areqs stepsL = [...]! and closes each case. \lexpy!qreqs_lines(tag, run, sc, steps)! writes,
per distinct triple $(e, j, b)$, the theorem \lexlean!qreq<tag><j>e<e>b<b> : qreqHolds m A (e, j, b)!
asserting $(b+1)A_j = b + \omega_m^{e}$, and the collector \lexlean!qreqsL : ∀ p ∈ aqreqs m stepsL, qreqHolds m A p! by peeling \lexlean!List.mem_cons! once per listed triple, closing with
\lexlean!simp at hp!. \lexpy!afac_lines(tag, run, sc)! writes
\lexlean!facL : afacs A stepsL = <Π facs>! by \lexlean!simp only [stepsL, afacs, AStep.fac, fac, A]!,
\lexlean!om_sub m 8 (m/8)! when $m\neq 8$ and an omega step occurred, and \lexlean!all_goals ring!,
returning the lines and the product model.

% ---------------------------------------------------------------------------
\section{Symbolic-Dimension Exports by the Normalizer}\label{sec:lean-export-sym}

The module \module{archytaszx.semantics.lean_sym} proves an equality of two diagrams for
every value of their symbolic dimensions, phase parameters and, for families with \bboxes,
multiplicities. Its single entry point \lexpy!export_symbolic! writes both diagrams as Lean
networks (through the network exporter of \module{archytaszx.semantics.lean_net},
\cref{ch:lean-export-net}), their dimensions as monomials in the dimension symbols (table
\lexlean!DT!), their phase vectors as roots of unity times atoms (table \lexlean!PT!) and their
scalars as constant term lists (\lexlean!S0!, \lexlean!S1!), and proves \lexlean!Equal! by one of the
\lexlean!equal_sym! theorems of \file{ArchytasZX/SymNet.lean} and its companion files
\file{SymHint.lean}, \file{SymFourier.lean}, \file{SymRange.lean}, \file{SymCarry.lean} and
\file{SymFold.lean} (\cref{ch:lean-sym}). The Lean kernel normalizes the scaled difference of
the two networks' index-sum expressions (\lexlean!decide +kernel!) and accepts it when it is the
empty sum, or when a certificate shows that the remaining terms sum to zero: a character
certificate (\lexlean!FzCert!), a range certificate (\lexlean!RgCert!), or a list of rewriting hints
over the labels' ranges (\lexlean!RHint!). The Python functions of the module mirror their Lean
namesakes on tuples, so that \lexpy!sym_check! predicts the Lean result before the file is
written and computes every certificate and hint the file states. Every unsupported input
raises \lexpy!NetUnsupported! (from \module{archytaszx.semantics.lean_net}).

\subsection{Constants}\label{sec:lean-export-sym-constants}

\begin{longtable}{@{}lp{0.17\textwidth}p{0.55\textwidth}@{}}
\toprule
Name & Value & Meaning \\
\midrule
\endhead
\lexpy!OFF! & $10^{6}$ & first summation variable of a node ($\mathit{OFF}+i+k$ for node $k$) \\
\lexpy!RBASE! & $2\cdot 10^{6}$ & first canonical bound-variable name \\
\lexpy!CP_BASE! & $3\cdot 10^{6}$ & digit variables of a compact table: $\mathit{CP\_BASE}+2s$ and $\mathit{CP\_BASE}+2s+1$ for the node variable $s$ (Lean \lexlean!cpX!) \\
\lexpy!ROUNDS! & 3 & elimination and canonicalization rounds per term \\
\lexpy!FULL_TERMS! & 4096 & most terms of the unnormalized translation \\
\lexpy!NORM_TERMS! & $2^{14}$ & term budget of the normalized translation \\
\lexpy!BS_NODES! & \lexpy!False! & whether every translation writes binder and splitter nodes (\lexlean!bsTerm!) \\
\lexpy!BS_SMALL! & 64 & most nodes of both networks for which an export writes binder and splitter nodes with \lexpy!BS_NODES! off \\
\lexpy!_BS_ON! & \lexpy!False! & whether the current export writes binder and splitter nodes (module state) \\
\lexpy!_FOLD! & \lexpy!False! & whether the compact check folds binder, compact spider and splitter blocks (\lexlean!trNetK!; module state) \\
\lexpy!BLOCK_MAX! & 8 & most nodes of one \lexlean!trBlockM! theorem \\
\lexpy!BLOCK_TERMS! & 24 & most nodes times terms of the largest stated expression in one \lexlean!trBlockM! theorem \\
\lexpy!TR_BLOCK! & 32 & most nodes of one \lexlean!trBlockM! theorem of an unnormalized translation (\lexlean!trNetG! with the identity) \\
\lexpy!RG_CASES! & 729 & most cases \lexpy!rg_cert! splits the labels into \\
\lexpy!RG_WORK! & $2^{16}$ & most cases times residual terms \lexpy!rg_cert! normalizes \\
\lexpy!RG_FUEL! & 256 & most pins and eliminations \lexpy!rg_cert! runs on one term \\
\lexpy!RX_ROUNDS! & 256 & most rounds \lexpy!rx_hints! runs on one residual term \\
\lexpy!RX_STEPS! & 32 & most steps of one round of \lexpy!rx_hints!, and the steps of one \lexlean!decide! in a chunked round \\
\lexpy!C_STEPS! & 4096 & most steps of one round of \lexpy!rx_cert_c! and of the first try of \lexpy!rx_cert_f! \\
\lexpy!CP_MARK! & \lexpy!(0, ())! & the modulus of a rule marking a compact table (\lexlean!cpData!) \\
\lexpy!CP_ROW_MIN! & 256 & a compact table with more entries is stored as \lexlean!cpRow! instead of entry by entry \\
\lexpy!TABLE_CHUNK! & 256 & (of \module{archytaszx.semantics.lean_net}) a flat phase table with more entries is read in chunks \\
\bottomrule
\end{longtable}
Integer division and remainder follow Lean's \lexlean!Int!: \lexpy!_emod(a, b)! is $a \bmod |b|$
(nonnegative), and $a$ itself for $b=0$; \lexpy!_ediv(a, b)! is $(a - \_emod(a,b))/b$, and 0 for
$b=0$. The caches are \lexpy!functools.lru_cache! of size $2^{18}$ on \lexpy!ip_varset!,
\lexpy!ip_lin! and \lexpy!elim_var!, and of size $2^{16}$ on the per-term normalization
\lexpy!_norm_cached(fl, t)!. \lexpy!_BS_ON! and \lexpy!_FOLD! are module globals that
\lexpy!_export_symbolic! and \lexpy!_SymExport.compact_check! set and restore.

\subsection{The data mirror}\label{sec:lean-export-sym-data}

\begin{definition}[Mirror types]\label{def:lean-export-sym-types}
\begin{itemize}
\item \lexpy!DMon = (c, e)!: the dimension monomial $c\prod_i d_i^{e_i}$, $c\in\Z$, $e$ a trimmed
exponent tuple (no trailing zeros) over the dimension symbols $d_0, d_1, \ldots$. Lean text
\lexlean!⟨c, [e0, e1, ...]⟩!.
\item \lexpy!IPoly!: an integer polynomial as a tuple of (monomial, coefficient) pairs, monomials
as sorted variable tuples with repetition, sorted increasingly, no zero coefficient.
\item \lexpy!Fac = (DMon, IPoly)!: a factor with modulus $D$ and polynomial $p$. In a delta list
it means $[\,p \equiv 0 \pmod D\,]$, written $[D \mid p]$; in a phase list (a \emph{root of
unity}) $e^{2\pi i p/D}$.
\item \lexpy!CMon = (q, r, sq, rc, am, gc)!: the coefficient
\[
q\, e^{2\pi i r} \prod_i \sqrt{d_i}^{\,sq_i} \prod_{c\in rc}\sqrt{c}\ \prod_j g_j^{\,am_j}
\prod_{(a, D)\in gc}\gcd(a, D),
\]
$q\in\Q$, $r\in[0,\tfrac12)$, $rc$ a sorted list of distinct integers, $g_j = e^{2\pi i\,
\mathrm{GS}_j}$ the phase atoms, $gc$ a sorted list of gcd factors.
\item \lexpy!STm = (CMon, sums, dels, phs)!: the term $\mathit{CMon}\cdot\sum_{x\in
\mathit{sums}}\prod \mathit{dels}\prod\mathit{phs}$, each summed variable $x$ with its
dimension monomial, ranging over $0 \le x < D(x)$.
\end{itemize}
An expression is a list of terms (their sum). A \emph{label} is a network label below
\lexpy!OFF!; a \emph{sum} is an entry of a term's sum list.
\end{definition}

\paragraph{Arithmetic.} Dense vectors add (\lexpy!nadd!), subtract with truncation at zero
(\lexpy!nsub!), compare pointwise (\lexpy!nle!) and \lexpy!trim!. Dimension monomials multiply
(\lexpy!dm_mul!), divide (\lexpy!dm_dvd!: $a_c\neq0$, $a_c\mid b_c$, exponents $\le$; \lexpy!dm_div!), test
constancy (\lexpy!dm_const!) and evaluate at symbol values (\lexpy!dm_den!). Polynomials insert a term
keeping order and dropping zeros (\lexpy!ip_ins!), add, scale, multiply, raise to powers,
substitute a variable (\lexpy!ip_subst!), split linearly in a variable (\lexpy!ip_lin(x, p)! returns
$(a, b)$ with $p = ax + b$ when $x$ has degree $\le 1$ in every monomial), read a constant
(\lexpy!ip_const!), reduce coefficients modulo $k$ to $[0,|k|)$ (\lexpy!ip_modc!), divide every
coefficient by $k$ with floor (\lexpy!ip_quo!), test and perform exact division of coefficients
(\lexpy!ip_divby!, \lexpy!ip_divc!), symmetric reduction (\lexpy!sym_mod(c, k)! is $r$ or $r-k$ with
$r=c\bmod k$, the latter when $2r>k$; \lexpy!ip_mods! applies it to every coefficient) and
normalize the sign so the first coefficient is positive (\lexpy!ip_sign!). \lexpy!pivot_of(p)! is the
largest variable $x$ with $p = sx + b$, $s = \pm1$, as $(x, s, b)$; \lexpy!pin_val(s, b)! is $-sb$.

\paragraph{Coefficient monomials.} \lexpy!rnorm(q, r)! reduces $r$ to $[0,1)$ and, when the
result is at least $\tfrac12$, negates $q$ and subtracts $\tfrac12$. \lexpy!rc_mul! multiplies two
square-root lists: equal entries $c$ cancel into a factor $c$. \lexpy!cm_mul! multiplies all
components (gcd factors merged sorted by $(a, c, e)$). Constructors: \lexpy!of_q(q)!;
\lexpy!rt_c(c, k)! $=\sqrt{c}^{\,k}$ (exact for even $k$ or a perfect square, else $c^{\lfloor
k/2\rfloor}$ with $\sqrt c$ in $rc$); \lexpy!rt_d(d, k)! $=\sqrt{d}^{\,k}$ for a monomial $d$;
\lexpy!gcd_a(a, d)! $=\gcd(|a|, d)$, rational when $d$ is constant or $|a|=1$; \lexpy!ex_q(r)!
$=e^{2\pi i r}$; \lexpy!atom_p(j, k)! $=g_j^{k}$; \lexpy!ph_c(q, ks)! $=e^{2\pi i q}\prod_j g_j^{k_j}$
(through \lexpy!rnorm(1, q)!).

\subsection{The normalizer}\label{sec:lean-export-sym-normalizer}

The four elimination rules (\lexpy!elim_unused!, \lexpy!elim_delta1!, \lexpy!elim_char!,
\lexpy!elim_gcd1!) evaluate the sum over one variable $x$ of dimension $d$ ($d_c\ne0$) in a term
with deltas $ds$ and phases $ps$, and return a coefficient and new lists:
\begin{enumerate}
\item \emph{Unused}: $x$ in no factor. Coefficient $d$ (\lexpy!rt_d(d, 2)!).
\item \emph{Delta} at $f=(D, p)$: $p = sx + b$, $s=\pm1$, $D \mid d$, every other factor
containing $x$ has a modulus dividing $D$. Substitute $x \mapsto -sb$ in all other factors;
coefficient $d/D$.
\item \emph{Character}: $x$ in no delta; every phase containing $x$ has modulus exactly $d$ and
is linear in $x$. With $A$ and $B$ the sums of their linear and constant parts,
$\sum_{x<d} e^{2\pi i(Ax+B)/d} = d\,[d \mid A]\,e^{2\pi i B/d}$: append the delta $(d, A)$,
replace those phases by $(d, B)$; coefficient $d$.
\item \emph{Gcd} at $f=(d, \alpha x + b)$: $\alpha$ constant, $\alpha \mid b$, $x$ in no other
factor. Drop $f$; coefficient $\gcd(|\alpha|, d)$.
\end{enumerate}
\lexpy!elim_var(x, d, ds, ps)! (cached) tries unused, delta (each delta reading $x$ in order),
character and gcd (each delta in order).

\begin{definition}[Delta split \texttt{split\_del(R, y, D, f)}]\label{def:lean-export-sym-split}
Mirror of \lexlean!splitDel!. $R$ is a list of (variable, dimension) pairs. Preconditions: $D$ is a
constant monomial with $D_c\ge2$, $D \mid M$ for $f=(M,p)$, $y$ occurs exactly once in $R$ and
$(y, D)\in R$, $p = sy + b$ with $s=\pm1$. Let $c$ be the constant part of $b$, $\mathit{rest}$
its non-constant part, $k = D_c$.
\begin{enumerate}
\item The \emph{partner}: the first $z$, in the order of $\mathit{rest}$, of a single-variable
monomial with coefficient $-s$, such that $z\ne y$, $z$ occurs once in $R$, $(z, D)\in R$ and
$z$'s only monomial in $\mathit{rest}$ is $-s z$. With a partner and
$\mathit{rem} = \mathit{rest} - (-sz)$: if $k \mid \mathit{rem}$ and $k\mid c$ the result is
$[D \mid sy - sz]$ and $[M/D \mid \mathit{rem}/k + c/k]$, else none (no other partner and no
fallback is tried).
\item Without a partner: if $k \mid \mathit{rest}$, with $c' = c \bmod k$ and $m = 1$ when $s=1$
and $c'\ne0$ (else $m=0$), the result is $[D \mid sy + c']$ and
$[M/D \mid \mathit{rest}/k + (c-c')/k + m]$; else none.
\end{enumerate}
\end{definition}
\lexpy!split_step(t)! replaces the first delta (in order) that some sum $y$ of $t$ reading it
splits by \lexpy!split_del(t.sums, y, D(y), f)! with its two deltas, in place.

\begin{algorithm}[Elimination choice \texttt{elim\_choice(t)}]\label{alg:lean-export-sym-elim}
Mirror of \lexlean!elimStepH!; returns the new term and the hint $(i, r, j)$: sum index $i$, rule
$r$ (0 unused, 1 delta, 2 character, 3 gcd, 4 split) and delta index $j$ (0 for rules 0 and 2).
Let $dd(x)$ be the indices of the deltas reading $x$, in order, and $pp$ the set of variables
the phases read (\lexpy!_occ!).
\begin{enumerate}
\item For each sum $(x, d)$ at index $i$ that occurs once in the sum list:
\begin{enumerate}
\item if $dd(x)$ is empty and $x\notin pp$: rule 0;
\item else, with $\mathit{mods}$ the moduli of the deltas $dd(x)$ (and of the phases reading
$x$ when $x \in pp$): for each $j \in dd(x)$ whose modulus every element of $\mathit{mods}$
divides, the delta rule at $j$; the first success is rule 1;
\item else, if $dd(x)$ is empty: the character rule (rule 2);
\item else, if $dd(x) = [j]$ and $x \notin pp$: the gcd rule at $j$ (rule 3).
\end{enumerate}
The first success removes $x$ from the sums and multiplies the coefficient.
\item For each delta $f$ at index $j$ and each sum $(y, d)$ at index $i$ with $y$ in $f$:
\lexpy!split_del(t.sums, y, d, f)!; the first success replaces $f$ in place (rule 4).
\item None.
\end{enumerate}
\lexpy!elim_step(t)! is the term of \lexpy!elim_choice(t)!; \lexpy!elim_all(n, t)! repeats it at most
$n$ times; \lexpy!elim_hints(n, t)! returns the final term and the list of hints.
\end{algorithm}

\begin{algorithm}[Canonical term \texttt{canon\_tm(fl, t)}]\label{alg:lean-export-sym-canon}
$fl$ is the floor vector of the dimension symbols.
\begin{enumerate}
\item Rename the summed variables to $\mathit{RBASE}+k$ in order (\lexpy!ren_list!; a variable that
occurs twice or whose target occurs is skipped).
\item \emph{Pivot pass} (\lexpy!piv_pass!) over the deltas in order: each delta $f$ yields a step
(\lexpy!_step!): a \emph{pivot} on the largest variable with coefficient $\pm1$ (value $-sb$),
applied by substitution to every factor whose modulus divides $f$'s and which contains the
variable; else a \emph{division} on the largest variable with a nonzero constant coefficient
$c$, which rewrites a factor $g = a_g x + b_g$ with $c \mid a_g$ and modulus dividing $f$'s to
$b_g - (a_g/c)\,b$. Each step applies to the deltas already passed, the deltas after it and the
phases.
\item Normalize the deltas (\lexpy!norm_dels!): the polynomial is symmetrically reduced modulo a
constant modulus and sign-normalized; a constant $0$ is dropped; a constant $c$ with constant
modulus $D$ is dropped when $D\mid c$ and makes the term vanish otherwise; a constant $c$ with a
symbolic modulus makes the term vanish when $|c| < D(fl)$ (the modulus at the floors). A
vanishing delta makes the term \lexpy!None!.
\item Merge phases of one modulus (\lexpy!ph_merge!), then normalize (\lexpy!norm_phs!): with a
constant modulus $D$ the polynomial is reduced modulo $D$, and a constant $c$ moves into the
coefficient as $e^{2\pi i c/D}$; an empty polynomial is dropped.
\item Sort deltas (deduplicated) and phases by (modulus coefficient, exponents, polynomial); the
coefficient is multiplied by the phase constants.
\end{enumerate}
\lexpy!norm_tm(fl, r, t)! performs $r$ rounds of \lexpy!elim_all(|sums| + 1)! followed by
\lexpy!canon_tm!. \lexpy!nf_ex(fl, e)! normalizes each term with \lexpy!ROUNDS! rounds (cached per floor
vector and term), drops vanished terms, and adds terms of equal shape (all components but $q$)
by their $q$, dropping zero sums, in insertion order.
\end{algorithm}

\subsection{Translation of networks}\label{sec:lean-export-sym-translate}

A network node is \lexpy!(kind, outs, ins, dim, ph)!: a kind string, output and input labels, a
dimension index into \lexlean!DT! and a phase index into \lexlean!PT! (for a binder or splitter, a
second dimension index). A phase vector row (\lexpy!PhRow!, a list with attribute \lexpy!rule!) lists
its stored entries $(j, a_j)$ (the vector is 1 elsewhere) and an optional \emph{rule}
$(D, (c_0, c_1, \ldots))$ for the extra phase $e^{2\pi i\,p(s)/D}$, $p(s) = \sum c_i s^{i}$
(\lexpy!r_poly!, \lexpy!r_f!). A rule with modulus \lexpy!CP_MARK! and four coefficients marks a
\emph{compact table} instead (below). With $s = \mathit{OFF}+i+k$ the node's variable and $d$
its dimension, \lexpy!tr_node! gives:
\begin{center}
\begin{tabular}{@{}lp{0.68\textwidth}@{}}
\toprule
Kind & Terms \\
\midrule
\lexcode!Z! & \lexpy!spider_c(1, d, s, row, rule, [(d, x - s) for each port x], [])! \\
\lexcode!X! & \lexpy!spider_c(√d^(-|ports|), d, s, row, rule, [], [(d, s·(Σ outs − Σ ins))])! \\
\lexcode!F! (1 in, 1 out) & $\sqrt d^{-1}$, phase $(d, o\cdot i)$ \\
\lexcode!T! (1 in, 1 out) & $[o] + [o-i] - [o][i]$ \\
\lexcode!Ti! (1 in, 1 out) & $[o-i] - [o] + [o][i]$ \\
\lexcode!W! (1 in) & $\sum_p [o_p - i]\prod_{q\ne p}[o_q] + (1-|\mathit{outs}|)[i]\prod_q[o_q]$ \\
\lexcode!I! (1 in, 1 out) & $[o-i]$ \\
\lexcode!B! (1 out, 2 in) & \lexpy!bs_term(d, DT(ph), o, i0, i1, s)!, when binder and splitter nodes are on \\
\lexcode!S! (2 out, 1 in) & \lexpy!bs_term(d, DT(ph), i, o0, o1, s)!, when binder and splitter nodes are on \\
\lexcode!ZP dx dy c!, two legs $x,y$ & \lexpy!zp_term(d, c, (x, y))!: the single term $(1,\{\},[\,],[(d,\ c\,x\,y)])$ (\lexlean!zpTerm!); any other leg count \lexpy!None! \\
other & \lexpy!None! (untranslatable) \\
\bottomrule
\end{tabular}
\end{center}
Here $[p]$ is the delta $(d, p)$. Binder and splitter nodes are on when \lexpy!BS_NODES! or
\lexpy!_BS_ON! holds. \lexpy!spider(k, d, s, row, ds, ps)! is the term
$(k, \{s{:}d\}, ds, ps)$ plus, per stored entry $(j, a)$, the two terms with coefficients
$k\,a$ and $-k$ and the extra delta $[s - j]$, so that the spider sums $1 + \sum_j [s=j](a_j-1)$.
\lexpy!spider_c(k, d, s, row, rule, ds, ps)! is \lexpy!cp_terms! when the rule marks a compact table
that \lexpy!cp_chk! accepts, else \lexpy!spider(k, d, s, row, ds, ps + r_f(s, rule))!.
\lexpy!bs_term(D_s, D_t, k, a, b, s)! is, for a constant $D_t$, the single term
$[D_sD_t \mid k - D_t a - b]$: the label $k$ of range $D_sD_t$ is the pair $(a, b)$ of ranges
$D_s$, $D_t$. For a symbolic $D_t$ (\lexlean!bsSym_den!) it is the character sum over the node's
variable $s$ of range $D=D_sD_t$: the single term with coefficient \lexpy!rt_d(D, -2)!, the sum $\{s{:}D\}$, no delta, and the phases $(D,\ s\,(k-b))$ and $(D_s,\ -s\,a)$.
\lexpy!leg_dms! gives each label's expected dimension: $[D_sD_t, D_s, D_t]$ for the
binder's (output, input, input), $[D_s, D_t, D_sD_t]$ for the splitter's (output, output,
input), $[\mathit{DT}(dx),\mathit{DT}(dy)]$ for a \lexcode!ZP! node, and the node's dimension for
every other port.

\subsubsection{Compact tables}\label{sec:lean-export-sym-compact}

A compact table is a phase vector whose entry $j$ is $e^{2\pi i\, n(j)/M}$ with
\[
n(j) = c\,j \ \text{(kind 1)},\qquad n(j) = c\,\lfloor j/\delta\rfloor\,(j \bmod \delta)\ \text{(kind 2)}.
\]
\begin{itemize}
\item \lexpy!cp_data(rule)! (\lexlean!cpData!): for a rule $(\mathit{CP\_MARK}, (k, m, c, \delta))$ the
quadruple $(\max(k,0), \max(m,0), c, \max(\delta,0))$; else \lexpy!None!.
\item \lexpy!cp_num(q, j)! (\lexlean!cpNum!): $c j$ for kind 1, else $c\lfloor j/\delta\rfloor(j\bmod
\delta)$ (with $\lfloor j/0\rfloor = 0$ and $j \bmod 0 = j$).
\item \lexpy!pt_ok(lst, c, pt)! (\lexlean!ptOK!): $c>0$, the keys of $pt$ are distinct, each entry
$(j, a)$ equals \lexpy!ph_c(lst[j]/c, ())! ($0$ beyond the list), and every $i$ with $c \nmid
lst_i$ is a key.
\item \lexpy!cp_chk(d, pt, q, s, ds, ps)! (\lexlean!cpChk!): $d$ constant positive and
\lexpy!pt_ok([cp_num(q, j) for j < d], M, pt)!; for kind 2 also $\delta^2 = d$ and neither digit
variable $\mathit{CP\_BASE}+2s$, $\mathit{CP\_BASE}+2s+1$ occurs in $ds$ or $ps$.
\item \lexpy!cp_terms(k, d, s, q, ds, ps)! (\lexlean!cpTerms!): kind 1 is the single term
$(k, \{s{:}d\}, ds, ps + [M \mid c\,s])$ (a phase of constant modulus $M$); kind 2, with
$x = \mathit{CP\_BASE}+2s$ and $y = x+1$, is the single term with sums $s{:}d$, $x{:}\delta$,
$y{:}\delta$, the extra delta $[d \mid s - \delta x - y]$ and the extra phase
$e^{2\pi i\, c x y / M}$.
\item \lexpy!cp_row(D, q)! (\lexlean!cpRow!): the entries $(j, e^{2\pi i\, n(j)/M}, n(j)/M)$ for
$j<D$ with $M \nmid n(j)$.
\end{itemize}

\subsubsection{Products}

\lexpy!tm_mul! multiplies coefficients and concatenates lists; \lexpy!mul_ok(s, t)! requires that
neither term's summed variables occur in the other's factors; \lexpy!ex_mul! and \lexpy!ex_mul_ok!
lift them to expressions. The translations differ in normalization and grouping:
\begin{itemize}
\item \lexpy!tr_nodes! (\lexlean!trNodesB!) multiplies the node translations from the last node to
the first; after each product it prefixes the sums of the internal labels whose first
occurrence (in node order) is at that node, and normalizes by \lexpy!nf_ex!; the product sizes are
charged against a budget of \lexpy!NORM_TERMS!, and exhausting it raises
\lexcode!the normaliser's expansion exceeds 16384 terms!. \lexpy!tr_net! (\lexlean!trNet!) adds the sums
of internal labels used by no node.
\item \lexpy!tr_nodes_full! and \lexpy!tr_net_full! (\lexlean!trNetF!) multiply without intermediate
normalization and prefix all internal sums.
\item \lexpy!tr_net_id! (\lexlean!trNetG! with the identity normalizer) multiplies like
\lexpy!tr_nodes! (node $k$ with variable $\mathit{OFF}+k$) without normalizing.
\item \lexpy!tr_net_k! (\lexlean!trNetK!) is \lexpy!tr_net_id! over the blocks of \lexpy!fold_steps!.
\item \lexpy!tr_steps! returns, for every $k$, the normalized expression of the nodes from $k$ on
(entry $k$ the product of \lexpy!nodes[k:]!, entry $m$ the empty product $[1]$).
\end{itemize}

\begin{algorithm}[Folded blocks]\label{alg:lean-export-sym-fold}
\lexpy!fold_bzs(n1, n2, n3)! (\lexlean!foldBZS!) folds a binder, a compact spider and a splitter
into one term. It requires $n_1 = \mathtt{B}\,(a, b) \to k$, $n_2 = \mathtt{Z}\,k \to k'$,
$n_3 = \mathtt{S}\,k' \to (a', b')$, the six labels $k, k', a, b, a', b'$ distinct, the binder's
second dimension and both splitter dimensions equal to the binder's constant dimension $\delta$,
the spider's dimension the constant $\delta^2$, and the spider's phase vector a compact table of
kind 2 with digit range $\delta$ accepted by \lexpy!pt_ok!. The result is the term
$[\delta \mid a' - a]\,[\delta \mid b' - b]\,e^{2\pi i\, c a b/M}$ with $k$ and $k'$.
\lexpy!fold_ok(internal, pre, k, k1, rest)! (\lexlean!foldOK!) requires that $k$ and $k'$ occur
neither in the labels before the block nor in the later nodes, and that the internal labels
among $\{k, k'\}$ are exactly both. \lexpy!fold_steps(dt, pt, nodes, internal)! scans from the
first node: at position $i$ a fold of nodes $i, i+1, i+2$ accepted by \lexpy!fold_ok! forms a
block of three, anything else a block of one; each step is (start, labels before, labels of
the block, folded term or \lexpy!None!). \lexpy!tr_net_k! does not sum $k$ and $k'$ of a folded
block.
\end{algorithm}

\paragraph{Network checks.} \lexpy!node_chk!: the node's dimension is nonzero and each stored phase
index $j < D(fl)$. \lexpy!net_chk! (\lexlean!netChk!): every node passes \lexpy!node_chk!, its dimension
index is in $ix$ and, for a binder or splitter, its second dimension index too; every label is
below \lexpy!OFF!, boundary or internal, has a dimension index in $ix$ and the dimension
\lexpy!leg_dms! gives for its port; every internal label's index is in $ix$.

\begin{algorithm}[\texttt{sym\_check(fl, dt, pt, labs, ix, n0, n1, s0, s1, fz, rg, rx)}]\label{alg:lean-export-sym-check}
A network is (nodes, internal labels, boundary labels). \lexpy!fz!, \lexpy!rg! and \lexpy!rx! are
optional lists that receive certificates.
\begin{enumerate}
\item The sorted boundaries must agree.
\item With binder and splitter nodes off, no node may be a binder or splitter.
\item \lexpy!net_chk! of both networks.
\item Both scalars must be constant (\lexpy!ex_const!: no sums, deltas or variables).
\item With \lexpy!rx!: translate both networks by \lexpy!tr_net_k! when \lexpy!_FOLD! holds, else by
\lexpy!tr_net_id!; both must translate and the products $s_0e_0$, $(-s_1)e_1$ be admissible. The
certificate is \lexpy!rx_cert_f! (with \lexpy!_FOLD!) or \lexpy!rx_cert! of the difference with the
ranges \lexpy!lab_ranges(labs, dt, ix, n0.boundary)!; success appends it and returns true, failure
returns false.
\item For the translations \lexpy!tr_net!, then \lexpy!tr_net_full! (skipped when either network's
full term count, the product of its nodes' term counts, \lexpy!_full_terms!, exceeds
\lexpy!FULL_TERMS!; a \lexpy!NetUnsupported! from \lexpy!tr_net! moves on to \lexpy!tr_net_full!, one from \lexpy!tr_net_full! propagates): if both translate and the products are admissible, let
$\mathit{rest} = \mathit{nf\_ex}(s_0 e_0 + (-s_1) e_1)$. Return true when $\mathit{rest}$ is
empty; with \lexpy!fz!, when \lexpy!fz_cert(rest)! succeeds (appended); with \lexpy!rg!, when
\lexpy!rg_cert(rest, lab_ranges(...))! succeeds (appended).
\item Return false.
\end{enumerate}
\end{algorithm}

\subsection{Character certificates}\label{sec:lean-export-sym-fz}

A normalized residual without sums whose moduli are constant is a finite trigonometric sum.
\lexpy!fz_cert! (Lean \lexlean!fzOK! of \file{SymFourier.lean}) proves it zero by writing every term
as characters $w\,\omega_N^{P}$ and every group of characters as a rational combination of the
\emph{cycles} $\sum_{i<p}\omega_N^{a + iN/p}$, which vanish for each prime $p \mid N$.

\begin{algorithm}[\texttt{fz\_cert(e)}]\label{alg:lean-export-sym-fz}
\begin{enumerate}
\item $e$ is nonempty. The modulus $N$ (\lexpy!_fz_mod!) is the lcm of every factor modulus and of
the denominator of every coefficient's $r$; it exists only when no term has a sum or a
coefficient component $sq$, $rc$, $am$, $gc$, and every modulus is a positive constant. It must
be at least 2.
\item The \emph{pinned} labels (\lexpy!_pinned!): every variable $x$ that some delta reads alone,
as $\pm x + c$. Sorted.
\item \emph{Explode} (\lexpy!_explode!, \lexlean!explode!): for each pinned label in order and each
residue $v < N$, substitute $x \mapsto v$ in all factors and prune (\lexpy!_prune!,
\lexlean!pruneTm!: a term with a constant delta $[D\mid c]$, $D\nmid c$, vanishes; satisfied
constant deltas are dropped). Each surviving term carries its residue vector $vv$.
\item Characters (\lexpy!exp_term!, \lexlean!expTerm!): the term $(q, r)$ starts as $[(q,
\omega_N^{rN})]$ ($rN$ must be an integer); each delta $[m \mid p]$, last first, expands as
$\frac1m\sum_{k<m}\omega_N^{k(N/m)p}$; each phase $e^{2\pi i p/m}$ multiplies by
$\omega_N^{(N/m)p}$. Every modulus must divide $N$.
\item Every character's exponent must be affine (monomials of degree at most 1). Group the
weights by $(vv, \mathit{key})$, $\mathit{key}$ the non-constant part reduced modulo $N$; within
a group, index by the constant part modulo $N$.
\item For each group with weight vector $v\in\Q^N$ (\lexpy!_cycle_cert!): the relations $(p, a)$
for the primes $p\mid N$ and $a < N/p$, each the 0/1 vector with ones at $a + iN/p$; solve
$M\lambda = v$ by Gauss--Jordan elimination (free parameters set to 0); no solution fails. Each
nonzero $\lambda$ gives the row $(vv, \mathit{key}, p, a, \lambda)$.
\item The result is $(N, (\mathit{pinned}, vs, \mathit{rows}))$, $vs$ the sorted variables of
all keys.
\end{enumerate}
\end{algorithm}
\lexpy!_lean_cert! writes \lexlean!((pinned, vs, [(vv, K, nat_lit p, nat_lit a, mkRat ...)]) : FzCert)!.

\subsection{Range certificates}\label{sec:lean-export-sym-rg}

The boundary labels have constant ranges in the theorem (\lexlean!labOK!);
\lexpy!lab_ranges(labs, dt, ix, bnd)! maps each boundary label whose dimension index is in $ix$ and
whose dimension is constant to that constant. A residual that vanishes only on these ranges is
certified by \lexpy!rg_cert! (Lean \lexlean!rgOK! of \file{SymRange.lean}).

\begin{itemize}
\item \lexpy!elim_pin1(x, d, pre, f, post, ps)! (\lexlean!elimPin1!): $f$'s modulus is $d$, a positive
constant, $f = sx + c$ with $s=\pm1$ and $c$ constant; then $x \mapsto (-sc) \bmod d$ in every
other factor, coefficient 1.
\item \lexpy!elim_var_r! (\lexlean!elimVarR!): \lexpy!elim_pin1! at the first delta that accepts, else
\lexpy!elim_var!.
\item \lexpy!split_step_r! (\lexlean!splitStepR!): \lexpy!split_step! without splits of a delta whose
modulus equals the splitting variable's range.
\item \lexpy!elim_step_r! (\lexlean!elimStepR!): \lexpy!elim_var_r! of the first sum occurring once
that accepts, else \lexpy!split_step_r!.
\item \lexpy!sub_tm(x, v, t)! substitutes a value; \lexpy!rsplit(r, e)! (\lexlean!rsplit!) is every
combination of values of the labels of $r = [(x, D_x)]$, the first label varying slowest.
\end{itemize}

\begin{algorithm}[\texttt{rg\_cert(e, ranges)}]\label{alg:lean-export-sym-rg}
\begin{enumerate}
\item $e$ is nonempty. Each term runs \lexpy!elim_step_r! to a fixed point (\lexpy!_pin_count!); it
must end without sums in fewer than \lexpy!RG_FUEL! steps. The fuel is the largest step count, at
least 1.
\item Every factor of the pinned terms must have a positive constant modulus $M$. A label $x$ of
$ranges$ is \emph{split} when some factor has the monomial $c\,x$ with $M \nmid c\,D_x$.
\item $r$ is the split labels, sorted, with their ranges; $\mathit{cases} = \prod D_x$. An empty
$r$ gives one case, the whole pinned residual. Fail when $\mathit{cases} > \mathit{RG\_CASES}$ or $\mathit{cases}\cdot|e| >
\mathit{RG\_WORK}$.
\item For each case of \lexpy!rsplit(r, pinned)!: prune the terms; an empty case has the
certificate $(0, ((), (), ()))$; any other needs \lexpy!fz_cert!.
\item The result is $(\mathit{fuel}, r, \mathit{certs})$.
\end{enumerate}
\end{algorithm}
\lexpy!_lean_rg! writes \lexlean!((fuel, [(x, d), ...], [(N, cert), ...]) : RgCert)!.

\subsection{Normalization with the labels' ranges}\label{sec:lean-export-sym-rx}

The compact check (\cref{alg:lean-export-sym-compactcheck}) cancels the unnormalized difference
by rewriting each term with steps that read the labels' ranges $r$ (a sorted list of (label,
range) pairs). A step is a hint $(i, \mathit{rule}, j)$; the Lean step functions are
\lexlean!elimStepX! (rules 0--4, 5, 6, 20), \lexlean!elimStepC! (adds rules 7--16) and
\lexlean!elimStepF! (adds rules 17 and 18). A term's \emph{range list} for a factor is:
\begin{itemize}
\item \lexpy!rg_of(r, sums)! (\lexlean!rgOf!): the sums followed by the labels of $r$ not among them,
each with its constant range;
\item \lexpy!rg_loc(r, sums, x, f)! (\lexlean!rgLoc!): for each variable $u\ne x$ of $f$'s monomials in
order, $(u, D)$ when the sums bind $u$ exactly once with dimension $D$, nothing when they bind
it more often, else $(u, D_u)$ when $u$ is a label of $r$;
\item \lexpy!rg_u(r, sums, y, f)! (\lexlean!rgU!): the range of $y$ (by \lexpy!rg_loc! of the polynomial
$y$) followed by \lexpy!rg_loc(r, sums, y, f)!.
\end{itemize}
\lexpy!ivl(rg, p)! (\lexlean!ivl!) is the interval $[lo, hi]$ of an affine $p$ when each variable $u$
ranges over $[0, D_u-1]$ ($D_u$ a positive constant in $rg$); \lexpy!None! otherwise.

\paragraph{Rule 5: geometric sums.} \lexpy!geo_data(x, L, lst)! (\lexlean!geoData!) is $(A, B)$ with
$\sum (L/m_f)(a_f x + b_f)$ over the phases $f=(m_f, a_fx+b_f)$ reading $x$, each modulus a
positive constant dividing $L$. \lexpy!elim_geo1(x, d, pre, f, post, ps)! (\lexlean!elimGeo1!): $d$
and $f$'s modulus $k$ positive constants, $x$ in no delta, $L = kd$, $A \equiv \pm p_f \pmod k$
coefficientwise; the result has coefficient $d$, replaces $f$ by $[L \mid A]$, removes the
phases reading $x$ and appends $(L, B)$. This is
$\sum_{x<d}\omega_L^{Ax+B}[k\mid p_f] = d\,[L\mid A]\,\omega_L^{B}$.

\paragraph{Rule 20: pins by range.} \lexpy!elim_xpin1(rg, x, d, pre, f, post, ps)!
(\lexlean!elimXpin1!): $f$'s modulus is $d$, a positive constant; $f = sx + b$, $s = \pm1$;
$v = -sb$ with \lexpy!ivl(rg, v)! $=[lo, hi]$; $K = \lfloor lo/d\rfloor d$; $hi < K + d$. Then
$x \mapsto v - K$ in every other factor, coefficient 1 ($x$ is the unique value in $[0, d)$
congruent to $v$). \lexpy!_xpin_fast! computes the same with the range list of \lexpy!rg_of! read
from dictionaries and fails when $v$ reads $x$.

\begin{algorithm}[\texttt{x\_choice(r, t)}]\label{alg:lean-export-sym-xchoice}
\begin{enumerate}
\item If the sum names are distinct: for each sum $(x, d)$ at index $i$ that a phase reads, no
delta reads, with $d$ a positive constant, and each delta $f$ at index $j$ with a positive
constant modulus $k$: rule 5 when \lexpy!geo_data(x, kd, ps)! exists (cached per $k$), its $A$
satisfies $A \equiv \pm p_f \pmod k$ and \lexpy!elim_geo1! accepts. Hint $(i, 5, j)$.
\item For each delta $f$ at index $j$ with a positive constant modulus, each sum $(x, d)$ at
index $i$ with $d$ equal to $f$'s modulus, $x$ bound once and read by $f$: rule 20 by
\lexpy!_xpin_fast! with the once-bound sums and the unbound labels of $r$. Hint $(i, 20, j)$.
\item \lexpy!elim_choice(t)!.
\end{enumerate}
\end{algorithm}

\paragraph{Carry steps.} \lexpy!c_fac(r, t, i, rule, j)! (\lexlean!carryFac!) rewrites the delta $f
= (M, p)$ at index $j$; $i$ indexes a sum, another delta (in the list without $f$), a phase or
a label, by rule. \lexpy!c_step! applies rule 15 (\lexpy!c_rc!) or \lexpy!c_fac!.
\begin{longtable}{@{}lp{0.10\textwidth}p{0.66\textwidth}@{}}
\toprule
Rule & $i$ & Condition and result \\
\midrule
\endhead
7 & sum $x$ & $M$ a positive constant, $p = sx + b$, $s=\pm1$: $f \mapsto [M \mid x - ((-sb) \bmod M)]$ (coefficients reduced by \lexpy!ip_modc!) \\
8, 9 & delta $g$ & $g$ has modulus $M$, a positive constant: $f \mapsto [M \mid (p \pm p_g) \bmod M]$ ($+$ for 8, $-$ for 9) \\
10 & sum $x$ & $x$'s range $d$ a positive constant, $M$ constant with $d \mid M$: insert $[d \mid p \bmod d]$ after $f$ \\
11 & delta $g$ & $g$'s modulus a positive constant multiple of $f$'s positive constant modulus $k$, $p_g \equiv \pm p \pmod k$: drop $f$ \\
12 & label $y$ & $rg = $ \lexpy!rg_u(r, sums, y, f)!, $D_y$ the range of $y$, $f$'s modulus $\ne D_y$: replace $f$ by \lexpy!split_del(rg, y, D_y, f)! \\
13 & phase $g$ & $f$ has a pivot $(x, s, b)$ and a positive constant modulus $m$, $g$ a positive constant modulus $G$, every monomial of $g$ containing $x$ has a coefficient $c$ with $G \mid cm$: $x \mapsto -sb$ in $g$ \\
14 & -- & $f \mapsto [M \mid -p]$ \\
15 & root $i$, root $j$ & (\lexpy!c_rc!, \lexlean!rcStep!) the square roots $a = rc_i$ and $b = rc_j$ of the coefficient, $ab$ a perfect square: remove both, $q \mapsto q\sqrt{ab}$ \\
16 & label $x$ & $rg = $ \lexpy!rg_u(r, sums, x, f)!, $x$'s range $d_x$ equals $f$'s modulus, a positive constant, $p = sx + b$: with $v = -sb$, \lexpy!ivl(rg, v)! $=[lo,hi]$, $K = \lfloor lo/d_x\rfloor d_x$, $hi < K + d_x$: $x \mapsto v - K$ in every other delta and every phase, $f$ kept \\
\bottomrule
\end{longtable}

\begin{algorithm}[\texttt{c\_choice(r, t)}]\label{alg:lean-export-sym-cchoice}
A step of \lexlean!elimStepC! and its hint; the first that applies:
\begin{enumerate}
\item \lexpy!x_choice(r, t)! when its rule is 5 or 20;
\item \lexpy!_c_pin!: rule 7 at the delta $j$ of modulus $d$ reading a sum $(x, d)$ bound once
(first sum, then first delta), when rule 7 changes the delta and \lexpy!elim_xpin1! with
\lexpy!rg_loc(r, sums without x, x, g)! then pins $x$ by the new delta $g$;
\item the result of \lexpy!x_choice! (rules 0--4);
\item \lexpy!_c_rc!: rule 15 at the first pair $i<j$ of square roots that merges;
\item \lexpy!_c_phase!: rule 13 at the first delta $f$ with constant modulus above 1 and a pivot,
and the first phase reading the pivot whose modulus does not divide $f$'s;
\item \lexpy!_c_tie!: rule 14 at the first delta of even constant modulus $k>2$ whose nonzero
symmetric residues (\lexpy!ip_mods!) start with $k/2$ and whose next residue other than $\pm k/2$
is negative;
\item \lexpy!_c_row!: rule 8, then 9, at the first delta $f$ (constant modulus above 1) and other
delta of the same modulus for which the result reads fewer variables than $f$;
\item \lexpy!_c_absorb!: rule 11 at the first delta (constant modulus above 1) and other delta of a
larger constant modulus;
\item \lexpy!_c_split!: rule 12 at the first delta and label of $r$ not bound in the sums;
\item \lexpy!_c_lpin!: rule 16 at the first delta $f$ with a positive constant modulus whose
largest variable $x$ another factor reads and whose modulus equals $x$'s range by
\lexpy!rg_u!;
\item \lexpy!_c_weaken!: rule 10 at the first delta $f$ (constant modulus $M>1$) and sum $(x, d)$
with $1<d<M$, $d\mid M$, $f$ linear in $x$ with coefficient $\pm1$, such that no delta of
modulus $d$ equals $\pm(p \bmod d)$ and some delta $g$ of modulus $d$ gives
$(p \bmod d) \pm p_g$ fewer variables than $p \bmod d$.
\end{enumerate}
\end{algorithm}

\paragraph{Folded steps.} \lexpy!carry17(i, pre, f, post)! (\lexlean!carry17!): $g$ the delta at index
$i$ of $pre + post$, $g$ with constant modulus $m$ and pivot $(x, s, b)$ ($q = p_g$), $f$ with
positive constant modulus $M$, $x$'s coefficient $c$ in $p$ constant, $M \mid csm$:
$f \mapsto [M \mid (p - c s q) \bmod M]$. \lexpy!_c17! applies it (rule 17) at the first delta $f$
of constant modulus above 1 and other delta $g$ whose pivot is a sum that $f$ reads, when the
result no longer reads the pivot. \lexpy!ph_split(t, i, j)! (\lexlean!phSplit!): the phase $(M, p)$ at
$i$, the sum $(x, d)$ at $j$, $d$ and $M$ positive constants, $M = kd$, every monomial of $p$
reading $x$ with a coefficient divisible by $k$: the phase keeps the monomials without $x$, and
$(d, p_x/k)$ is appended. \lexpy!_c18! applies it (rule 18) when the sum names are distinct, for
the first sum $x$ of range $d>1$ that no delta and some phase reads, at the first phase reading
$x$ with modulus $\ne d$, provided every such phase admits \lexpy!ph_split!. \lexpy!f_choice!
(\lexlean!elimStepF!) is \lexpy!x_choice!, else rule 17, else rule 18.

\begin{algorithm}[Rounds \texttt{rx\_hints(fl, r, t, choice, steps, mids)}]\label{alg:lean-export-sym-rxhints}
Mirror of \lexlean!normTmX! (choice \lexpy!x_choice!, default), \lexlean!normTmC! (\lexpy!c_choice!) and
\lexlean!normTmF! (\lexpy!f_choice!); \lexpy!steps! defaults to \lexpy!RX_STEPS!. At most \lexpy!RX_ROUNDS!
rounds:
\begin{enumerate}
\item Apply \lexpy!choice! up to \lexpy!steps! times to $t_1 = t$, collecting the hints $es$; after
every \lexpy!RX_STEPS!-th step record the term in $\mathit{cut}$; finally record $t_1$ unless it
is the last recorded.
\item $t_2 = $ \lexpy!canon_tm(fl, t1)!. With \lexpy!mids!, append $\mathit{cut}$ when $t_2 \ne t$.
\item Stop when $t_2 = t$. Otherwise record the round hint
$(\mathit{canon\_hint}(t_1), es)$; when $t_2$ is \lexpy!None! return \lexpy!None!; else $t := t_2$.
\end{enumerate}
The result is the final term (or \lexpy!None!), the round hints and the term before each round.
\end{algorithm}
\lexpy!rx_cert(fl, e, ranges)! (\lexlean!rxNf!): \lexpy!None! for an empty $e$ or no ranges; with
$r$ the sorted ranges, \lexpy!rx_hints! per term; when the results add (\lexpy!tm_add!) to nothing,
$(r, \mathit{hints}, \mathit{paths})$ with each path the terms before each round and the
result; else \lexpy!rx_cert_c(fl, e, r)!. \lexpy!rx_cert_c! (\lexlean!rxNfC!) runs \lexpy!rx_hints! with
\lexpy!c_choice!, \lexpy!C_STEPS! steps and \lexpy!mids!, and returns
$(r, \mathit{hints}, \mathit{paths}, \mathit{cuts})$ when the results cancel. \lexpy!rx_cert_f!
(\lexlean!rxNfF!) tries \lexpy!f_choice! with \lexpy!C_STEPS! steps (fourth entry
\lexpy!("F", cuts)!), then with \lexpy!RX_STEPS! steps (fourth entry \lexpy!"F"!). \lexpy!_lean_rx!
writes the ranges as \lexlean![(x, d), ...]! and the rounds as
\lexlean!([[(nat_lit c, [(nat_lit i, nat_lit k, nat_lit j), ...]), ...], ...] : List (List RHint))!.

\subsection{Hints for the stepwise check}\label{sec:lean-export-sym-hints}

The stepwise check (\cref{sec:lean-export-sym-checkdecls}) replaces the normalizer's searches in
the Lean kernel by hints. A \emph{term hint} (\lexlean!THint!) is $(\mathit{mask}, \mathit{ilv},
\mathit{rounds})$, a round hint $(c, es)$ the pivot mask $c$ and the elimination hints $es$.
\begin{itemize}
\item \lexpy!piv_pass_b(m, ds, ps)! (\lexlean!pivPassM!): the pivot pass applying only the steps of
the deltas whose bit is set in $m$ (all for \lexpy!None!), with the mask of the steps that changed
a factor.
\item \lexpy!canon_hint(t)! (\lexlean!canonTmH!): with $t$ renamed, the mask of the changing steps
when the pass restricted to it reproduces the full pass, else $2^{|\mathit{dels}|}-1$.
\item \lexpy!norm_hints(fl, t)!: up to \lexpy!ROUNDS! rounds of \lexpy!elim_hints! and
\lexpy!canon_tm!; each round that changes the term records $(\mathit{canon\_hint}(t_1), es)$,
ending after a vanishing round.
\item \lexpy!ilv_mask(a, b, out)!: the mask (bit $k$ for position $k$ taken from $a$) of the greedy
interleaving of $a$ and $b$, preferring $a$, that gives $out$; \lexpy!None! when none does.
\item \lexpy!term_hint(fl, t)! (\lexlean!normTmK!): let $\mathit{full}$ be $t$'s normal form and
$\mathit{idx}$ the deltas of $t$ that read no sum of $t$ and occur in $\mathit{full}$. When
$\mathit{idx}$ is nonempty and $t$ without them normalizes to the same coefficient, sums and
phases, the full deltas are an interleaving of the kept deltas and the rest's, and no kept delta
reads the rest's sums: $(\mathit{mask}(\mathit{idx}), \mathit{ilv}, \mathit{norm\_hints}(t
\setminus \mathit{idx}))$; else $(0, 0, \mathit{norm\_hints}(t))$. \lexpy!_mask(labels)!
(\lexlean!maskOf!) sets bit $l$ per label.
\item \lexpy!tr_hints(fl, dt, pt, labs, nodes, internal, steps)!: for node $k$, the term hints of
each term of its translation times \lexpy!steps[k+1]!, the labels first met at $k$ summed.
\end{itemize}

\paragraph{Literal writers.} The certificates and hints are written on raw literals:
\begin{center}
\begin{tabular}{@{}lp{0.62\textwidth}@{}}
\toprule
Writer & Lean text \\
\midrule
\lexpy!_rn(x)! & \lexlean!nat_lit x! \\
\lexpy!_rz(x)! & \lexlean!Int.ofNat (nat_lit x)!, or \lexlean!Int.negSucc (nat_lit (-x-1))! for $x<0$ \\
\lexpy!_rq(q)! & \lexlean!mkRat (num) (nat_lit den)! with the numerator by \lexpy!_rz! \\
\lexpy!_rdm(m)! & \lexlean!⟨nat_lit c, [nat_lit e, ...]⟩! \\
\lexpy!_lean_cm(c)! & \lexlean!⟨q, r, [sq], [rc], [am], [(a, D)]⟩!, $sq$, $am$, $a$ by \lexpy!_rz!, $rc$ by \lexpy!_rn! \\
\lexpy!_lean_ip(p)! & \lexlean![([x, ...], c), ...]! \\
\lexpy!_lean_fac(f)!, \lexpy!_lean_facs! & \lexlean!(D, p)!, a list of them (each optionally by a name function) \\
\lexpy!_lean_tm(t)! & \lexlean!⟨c, [(x, D), ...], dels, phs⟩! \\
\lexpy!_lean_ex(e)! & \lexlean!([t, ...] : List STm)! \\
\lexpy!_lean_th(h)! & \lexlean!(nat_lit mask, nat_lit ilv, [(nat_lit c, [(i, r, j), ...]), ...])! \\
\lexpy!_lean_dm!, \lexpy!_lean_q!, \lexpy!_lean_ints! & plain \lexlean!⟨c, [e]⟩!, \lexlean!(a : ℚ)! or \lexlean!((a : ℚ) / b)!, integer lists with \lexlean!(x : ℤ)! for negatives \\
\bottomrule
\end{tabular}
\end{center}

\subsection{The export}\label{sec:lean-export-sym-export}

\begin{signature}
\lexpy!export_symbolic(left, right, name="main", describe=None, floors=None, shift=None, original=None, fixed=None) -> str!
\end{signature}
Lean source proving \lexpy!left = right! for every value of the dimensions and phases;
\begin{itemize}
\item \lexpy!floors! (default \lexpy!dimension_floors(left, right)! of
\module{archytaszx.semantics.decide}) gives lower bounds of the dimension symbols;
\item \lexpy!shift = (n, c)! proves the statement at $n+c$ for every $n$;
\item \lexpy!fixed! is a list of clauses \lexpy!(n, c)! or \lexpy!(n, c, at)!: at $n=c$ when \lexpy!at! (the
default), else at $n + c$;
\item \lexpy!original! restates the theorem over the original diagrams as the network exporter
does (\cref{ch:lean-export-net}).
\end{itemize}
It calls \lexpy!_export_symbolic(..., roots=False)!; when that raises \lexpy!NetUnsupported! and
either diagram's scalar contains an exponential (\lexpy!sp.exp! among \lexpy!_scalar_parts!), it
returns \lexpy!_export_symbolic(..., roots=True)! (whose errors propagate); otherwise it re-raises.
\lexpy!_export_symbolic! builds a \lexpy!_SymExport!, sets \lexpy!shift!, \lexpy!fixed! and the exporter's
\lexpy!original!, sets \lexpy!_BS_ON! to \lexpy!_bs_small(ex)! (true for a pair without multiplicities
with a binder or splitter node and at most \lexpy!BS_SMALL! nodes in all), raises
\lexpy!NetUnsupported! with the recorded reason (default
\lexcode!the normal forms of the two networks differ!) when \lexpy!check()! fails, returns
\lexpy!source(name)!, converts \lexpy!ScalarUnsupported! to \lexpy!NetUnsupported!, and restores
\lexpy!_BS_ON!.

\begin{algorithm}[\texttt{\_SymExport} construction]\label{alg:lean-export-sym-init}
\lexpy!_SymExport(left, right, describe, floors, roots=False)!.
\begin{enumerate}
\item Create the network exporter state (\lexpy!_Export!) with mixed dimensions allowed.
\item \emph{Unnest} each diagram (\lexpy!_unnest!): a \bbox nested in $t>0$ \bboxes of the same
symbolic multiplicity gets the fresh multiplicity \lexcode!<symbols>_nest<t>! (underscores
appended on a clash); the aliases map fresh names to the originals.
\item Integer phase symbols that are not dimension symbols are registered with type
\lexlean!ℤ!.
\item \emph{Family multiplicities} (\lexpy!_family_mults!): every \bbox multiplicity, bare symbols
first, then by text; a node-scoped \bbox whose crossing wires do not all end at Z or X spiders
raises \lexcode!a box meeting the rest at a wire!. The bare multiplicities are the family indices.
\item Label both networks and register them; model both scalars. With \lexpy!roots!, each scalar
model passes through \lexpy!with_roots! (below).
\item \lexpy!tables(m0, m1, floors)! builds the rest (the steps below).
\item When either scalar reads an opaque symbol or a multiplicity and the right scalar is
nonzero, the ratio $m_0/m_1$ (simplified, negative powers written as \lexpy!Inv! by
\lexpy!_no_negative_powers!) free of opaque and multiplicity symbols becomes the left scalar, $1$
the right, and $m_1$ the \emph{common} factor; the stated scalars stay the originals. Opaque
symbols without a common factor raise
\lexcode!the diagrams' scalars differ by a factor outside the exact model!; \bboxes without bare
multiplicities raise \lexcode!bang boxes!.
\item Dimension symbols: the free symbols of every dimension expression, sorted by name, then
every power atom $x^{e}$ with symbolic $e$ as an extra symbol. \lexlean!DT! is each dimension
expression as a monomial (\lexpy!_dmon!: a positive integer coefficient times nonnegative integer
powers of symbols; $\pi$ raises
\lexcode!a phase that is not a rational number of turns (a real angle)!, anything else
\lexcode!the dimension ... is not a monomial!). A mixed pair of dimension indices $(i_q, i_p)$
shares the monomial of $i_p$ and adds the hypothesis \lexlean!(dim_p) = (dim_q)!; with a family
this raises \lexcode!a family whose nodes' legs differ in dimension!.
\item Phase rules (\lexpy!rule_of!): a rule $p(k)/D$ in the phase index $k$ must be a polynomial
with integer coefficients and zero constant term over a monomial $D$
(\lexcode!the phase rule ... is not a polynomial over a dimension!,
\lexcode!... has no integer form!).
\item Atoms (\lexpy!_Atoms!): every phase entry (as \lexpy!phase_value(v)! gives it, the identity) and
every \lexpy!Ex! argument (over $2\pi i$) splits into a rational plus rational multiples of
symbolic bases; within a phase entry, integer multiples of integer symbols drop out. The bases
are sorted by text, each with the lcm $L$ of its coefficient denominators; the atom $g_j$ is
$e^{2\pi i\, b_j/L_j}$ and an entry is $e^{2\pi i q}\prod_j g_j^{k_j}$. A term with a
non-rational coefficient raises \lexcode!the phase term ...!.
\item Floors: the given floors, raised for each Z or X node with stored phase entries on a
dimension that is a bare symbol to one more than its largest stored index; the floor vector
keeps values above 1.
\item Constant scalar terms (\lexpy!scalar_terms!) when no multiplicity enters the scalars: each
additive term loses its roots of unity of symbolic order (\lexpy!root_facs!), which become
constant phase factors, and its rest becomes a coefficient (\lexpy!monomial!): a product of
rationals (\lexlean!CMon.ofQ!), powers of $i$ (\lexlean!CMon.exQ (k/4)!), of $\omega_n$ for an integer
$n$ (\lexlean!CMon.exQ (k/n)!), of square roots of dimension monomials (\lexlean!CMon.rtD!), of
dimension symbols and their inverses (\lexlean!dP d k!), and of \lexpy!Ex! (\lexlean!phC q ks!), combined
by \lexlean!CMon.mul!; any other factor raises \lexcode!the scalar factor ...!.
\item $ix$ is every dimension index not used by a bundle.
\end{enumerate}
\end{algorithm}

\paragraph{Roots of unity of symbolic order.} \lexpy!with_roots(model, dims)! rewrites each
\lexpy!Ex(2πi c/M)! with $c$ an integer and $M$ a monomial over symbols all among the dimension
symbols as \lexpy!Om(M)**c! (\lexpy!Inv(e)! read as $1/e$ first), then applies
\lexpy!_no_negative_powers!. \lexpy!root_facs(term)! removes every factor \lexpy!Om(M)**k! or
\lexpy!Inv(Om(M))**k! with a symbolic $M$ and integer $k$ and returns them as the phase factors
$(\mathit{dmon}(M), \pm k)$. \lexpy!scalar_text(t)! writes a scalar term as
\lexlean!⟨text, [], [], phs⟩!.

\begin{algorithm}[\texttt{check()}]\label{alg:lean-export-sym-checkalg}
Without family indices:
\begin{enumerate}
\item If the $\mathbb{F}_2$ plan applies (\lexpy!q_plan!, the $\mathbb{F}_2$ export section) the
check succeeds.
\item Reset the certificate lists \lexpy!fz!, \lexpy!rg!, \lexpy!rx! and \lexpy!cp_rows!; let $\mathit{args}$
be \lexpy!(fl, dtf, ptf, labs, ix, net(a), net(b), s0, s1)!.
\item $\mathit{first}$ = \lexpy!chained_gadgets()!: some splitter output of either network is a
binder input. If $\mathit{first}$ and \lexpy!compact_check(args)!, succeed.
\item Reset \lexpy!fz!, \lexpy!rg!; \lexpy!sym_check(*args, fz, rg)! succeeds or fails; a
\lexpy!NetUnsupported! it raises is held.
\item If not $\mathit{first}$ and \lexpy!compact_check(args)!, succeed.
\item Raise the held exception, else fail.
\end{enumerate}
With family indices, the family proof is planned once (\lexpy!fixed! by \lexpy!plan_fixed!,
\lexpy!shift! by \lexpy!plan_shift!, else \lexpy!fam! on the two networks), recording
\lexpy!NetUnsupported! messages as the reason.
\end{algorithm}

\begin{algorithm}[Compact check \texttt{compact\_check(args)}]\label{alg:lean-export-sym-compactcheck}
\begin{enumerate}
\item Save the rules and phase table; clear \lexpy!cp_rows!. Fail unless \lexpy!compact_rules()!
marks a table.
\item Reset \lexpy!fz!, \lexpy!rg!; set \lexpy!_BS_ON!.
\item If either network has a folded block (\lexpy!fold_steps!): with \lexpy!_FOLD! set,
\lexpy!sym_check(..., rx=self.rx)!; success sets \lexpy!fold! and succeeds. Clear \lexpy!_FOLD! and
\lexpy!rx!.
\item \lexpy!sym_check(..., rx=self.rx)! succeeds.
\item Restore \lexpy!_BS_ON!, the rules and the table; clear \lexpy!cp_rows! and \lexpy!rx!; fail.
\end{enumerate}
A \lexpy!NetUnsupported! raised by steps 3--4 runs step 5's restore, then propagates.
On success \lexpy!_BS_ON! stays set until \lexpy!_export_symbolic! restores it.
\end{algorithm}

\begin{algorithm}[\texttt{compact\_rules()}]\label{alg:lean-export-sym-compactrules}
\begin{enumerate}
\item The uses of each phase index: per Z or X node (kind, dimension monomial, ports); per other
node with a nonempty phase row, except binders and splitters, a blank use.
\item For each phase vector $p$ without a rule, with a nonempty row and some use: every entry
must be $\pm e^{2\pi i r}$ with no other coefficient component; its turns are
$(r + \tfrac12[q=-1]) \bmod 1$. All uses must share one dimension, none be blank, and the
dimension be a constant $D \ge 2$. $M$ is the lcm of the turns' denominators, $v_j$ the turns
at $j<D$.
\item Candidates: kind 1 $(1, M, v_1 M \bmod M, 0)$; when $D = \delta^2$, $\delta>1$ and every
use is a Z node, also kind 2 $(2, M, v_{\delta+1} M \bmod M, \delta)$. The first that
\lexpy!cp_chk(D, row, q, OFF, (), ())! accepts sets the rule $(\mathit{CP\_MARK}, q)$; when
$D > \mathit{CP\_ROW\_MIN}$ it records \lexpy!cp_rows[p] = (D, q)! and replaces the row by
\lexpy!cp_row(D, q)!.
\item Return whether a table was marked.
\end{enumerate}
\end{algorithm}

\subsection{Families}\label{sec:lean-export-sym-families}

\begin{algorithm}[Family lemma \texttt{fam(x, xt, y, yt, sc)}]\label{alg:lean-export-sym-fam}
Returns the name \lexlean!P<i>! of a lemma
\lexlean!∀ vs : ℕ, Equal (F.env (ksOf vs)) sx X sy Y!; \lexcode!sc! holds both scalar models and texts.
\begin{enumerate}
\item If normalizing the nodes (unbundled tails written plain, \lexpy!_norm!) changes either
network, declare the normalized networks \lexlean!T<i>n!, \lexlean!T<i>m!, recurse on them, and prove
\lexlean!P<i>! by \lexlean!Equal.mul_trans! with two \lexlean!iso_b! isomorphisms and \lexlean!Equal.congr_s!.
\item If no label lies in a \bbox (no chain heads), prove the pair by the normalizer: the
scalar tables \lexlean!SS<n>! and lemmas \lexlean!scalS<n>! (\lexpy!stab!; a scalar depending on a
multiplicity raises \lexcode!a scalar depending on a multiplicity outside its boxes!),
\lexpy!sym_check! (unequal normal forms raise \lexcode!the normal forms of the two networks differ!),
the check declarations of \lexpy!check_decls("C<i>", ..., [("C<i>x", xt, x), ("C<i>y", yt, y)], ...)!
and \lexlean!P<i>! by \lexpy!sym! (\cref{sec:lean-export-sym-checkdecls}).
\item Otherwise choose the \bbox index $m$: the least head that is neither nested in another
nor ported, else the least head not nested (none raises \lexcode!boxes that do not nest!). The
offset is $H = 2^{\max(16, \mathrm{bitlen}(\text{largest label}))}$; register the copies at
$H$ of every label whose chain starts with $m$.
\item A ported $m$ tries \lexpy!port! (below); on failure, re-raise unless a summed label of $x$
lies in the \bbox.
\item The ratios $\alpha$, $\beta$ by which the scalars grow when $m$ grows by 1 (\lexpy!ratio!: a
zero scalar depending on $m$ raises \lexcode!a zero scalar depending on a multiplicity!; a ratio
still depending on a multiplicity raises \lexcode!a scalar not geometric in a multiplicity!).
\item The base lemma on the zero networks (\lexlean!Net.zero (SpC m)!) at $m=0$. The split networks
(\lexlean!Net.split (SpC m) H!) and their peeled copies (\lexpy!_peeled!: the split's nodes and
labels not in the original); when the original is not contained in its split, the
\lexpy!tower! construction applies. Unequal peeled boundaries raise
\lexcode!the peeled copies have different boundaries!.
\item The step lemma \lexlean!Q<i>! proves the split equality at $k+1$ from the hypothesis at $k$
by a chain of \lexlean!Equal.mul_trans! of: an isomorphism from the split to
\lexlean!⟨T.nodes ++ X.nodes, ...⟩!, \lexlean!plug! of the hypothesis, an isomorphism reordering,
\lexlean!plug! of the copies' lemma (recursively \lexpy!fam! on the peeled copies with scalars
$\alpha,\beta$), and an isomorphism to the right split.
\item \lexlean!P<i>! by \lexlean!induction! on the multiplicity: \lexlean!zero! by \lexlean!Equal.of_zero_at!,
\lexlean!succ k_ ih! by \lexlean!Equal.of_split_at! and \lexlean!Q<i>!; split side conditions
\lexlean!Spec.splitOK_struct _ m _ _ H t (by decide +kernel : Spec.splitOK (SpC m) H t = true)!.
\end{enumerate}
\end{algorithm}

\begin{algorithm}[Tower \texttt{tower(i, x, xt, y, yt, sc, m, H, vs, lz, ratios)}]\label{alg:lean-export-sym-tower}
The step for a node-scoped \bbox along $m$ whose copies meet Z or X spiders of $x$ at their
bundled inputs.
\begin{enumerate}
\item The \emph{spiders}: the Z and X nodes of $x$ with a summed input in the \bbox; none
raises \lexcode!a node box meeting no Z or X spider!. A bundle node along another index raises
\lexcode!towers along two multiplicities!. Split both networks at $H$; every other node of $x$
must occur in the left split (\lexcode!the split network does not contain the network!); the rest
of the split is the copy plus the split spiders.
\item \emph{Held} spiders have no boundary leg (\lexpy!tower_leg!: a boundary output, or for a Z
spider a plain boundary input). Each held spider needs a \emph{target} (\lexpy!tower_target!: a
phaseless spider of its kind, not held, joined to it, itself phaseless, at exactly one summed
plain label, \lexpy!tower_shared!; for X spiders the held spider's output and the target's plain
input, neither with bundled outputs); without one, a Z spider with a summed plain leg $u$ that
is the input of an F node whose output is a boundary label is \emph{opened}
(\lexpy!tower_open_leg!); otherwise
\lexcode!the spider meeting the node box has no boundary leg and no neighbour of its colour with one!.
\item Each remaining held spider $z_n$ moves the legs its split adds onto its target $z_m$
(\lexpy!tower_move!): the split spiders of both (unique node in the rest of the same kind and
outputs, else \lexcode!the split spider is not found!; the added inputs must extend $z_n$'s inputs
with plain chains, else \lexcode!the split spider's legs are not the unfused form!) are fused along
the joining label (\lexpy!tower_fuse!: \lexlean!fusionZG! or \lexlean!fusionXG!), and the fused spider
is unfused into $z_n$ and $z_m$'s split with the added legs. Declarations
\lexlean!T<i>m<j>a! to \lexlean!T<i>m<j>d!; steps: isomorphism, \lexlean!plug! of the fusion,
isomorphism, \lexlean!plug! of the reversed fusion.
\item Each other spider is \emph{cut} (\lexpy!tower_cut!, or \lexpy!tower_cut_open! at the opened leg
$u$): the mode is \lexcode!xout! (an X spider's boundary output; an X spider with bundled outputs
raises \lexcode!an X spider meeting the node box with bundled outputs!), \lexcode!zout! (a Z spider's
boundary output) or \lexcode!zin! (a Z spider's plain boundary input); no leg raises
\lexcode!the spider meeting the node box has no boundary leg!. The split spider is found as above
and must gain inputs (\lexcode!the split spider gains no leg!) in the unfused form (inputs plain
inputs, tail, added legs; chains extended by empty chains); a split X spider with bundled
outputs raises \lexcode!the split X spider has bundled outputs!.
\item For each cut $j$ with leg $o$: a fresh wire $w$; \lexpy!tower_fusion! gives $z_a$ ($z_n$ with
$o$ renamed $w$) and $z_b$ (a phaseless spider on $o$, $w$ and the added legs: inputs
$(o, w, \ldots)$ for \lexcode!zin!, output $o$ and inputs $(w, \ldots)$ for \lexcode!zout! and
\lexcode!xout!) with the local equation \lexlean!fusionZG! or \lexlean!fusionXG!. Declarations
\lexlean!T<i>s!, \lexlean!T<i>u! (with suffix $j$ when several cuts); steps: isomorphism to the split
network with its spider last, \lexlean!plug! of the reversed fusion. An opened spider then puts
the four F boxes $w \to w_1 \to w_2 \to w_3 \to w_4$ before $z_b$ (\lexpy!tower_opening!: the
lemma that one spider equals the F chain followed by the spider on $w_4$, by a recursive
\lexpy!fam! with scalars 1, declarations \lexlean!T<i>oa!, \lexlean!T<i>ob!, \lexlean!T<i>o<j>!,
\lexlean!T<i>p<j>!), and its boundary label maps to the first F box's output.
\item Rename: each cut's boundary leg to its fresh wire and every summed label of $y$ to a fresh
label (\lexpy!_perm_to_bijection!); \lexlean!plug! the hypothesis renamed by \lexlean!Equal.rename!;
declare the rest \lexlean!T<i>r! and recurse \lexpy!fam! on it against the right split with the
ratios.
\item \lexlean!Q<i>! is the chain of \lexlean!Equal.mul_trans! of all steps under
\lexlean!Equal.congr_s ... (by ring) (by ring)!; \lexlean!P<i>! by induction as in
\cref{alg:lean-export-sym-fam}.
\end{enumerate}
\end{algorithm}
\lexpy!port! handles a port-scoped \bbox: both networks' port legs (\lexpy!_port_legs!, the mirror of
\lexlean!Net.ports! under \lexlean!Spec.portNet!) must agree and the ratios must be equal
(\lexcode!a port box the port step does not reach!, \lexcode!a port box whose scalars grow unequally!);
the lemma has base cases at 0 and 1 (offset $2H$) and the step \lexlean!Equal.of_port_at! with
\lexlean!Spec.portOK!. \lexpy!plan_shift! splits both networks $c$ times at offsets $H\cdot 2^{t}$
(\lexpy!offset_base!: $2^{\max(16,\mathrm{bitlen})}$ over every label of both networks) and
proves the family lemma at $n+c$, lifted by \lexlean!Equal.of_split_at! once per split
(\lexpy!shift_proof!; \lexcode!the shifted symbol carries no box!,
\lexcode!the shifted index is not a bare multiplicity!). \lexpy!plan_fixed! orders the clauses by
nesting depth, splits each $c$ times and zeroes it when \lexcode!at! (a non-fixed \bbox inside
another at $n \ge c$ raises \lexcode!a lifted box inside a box at n >= c!; also
\lexcode!a fixed symbol carries no box!, \lexcode!a fixed index is not a bare multiplicity!), and lifts
by \lexlean!Equal.of_zero_at! and \lexlean!Equal.of_split_at! (\lexpy!fixed_proof!). A shift or fixed
clause with a common scalar raises \lexcode!a shifted family with a common scalar! or
\lexcode!a fixed family with a common scalar!.

\subsection{The emitted source}\label{sec:lean-export-sym-source}

\lexpy!source(name)! computes the shared texts (\lexpy!texts!) and writes the $\mathbb{F}_2$ proof
(\lexpy!q_source!) when the $\mathbb{F}_2$ plan applies; otherwise the prelude (\lexpy!prelude!),
the check declarations and the theorem. The prelude opens with
\lexlean!open ArchytasZX.Sym!, \lexcode!set_option linter.unusedSimpArgs false!,
\lexcode!set_option Elab.async false! and \lexcode!set_option linter.all false!, then:
\begin{lstlisting}[style=lexlean]
def DT (i : ℕ) : DMon := ([⟨1, [1]⟩, ⟨2, []⟩] : List DMon).getD i ⟨1, []⟩
def PT (i : ℕ) : List (ℕ × CMon) := ([[(1, phC ((1 : ℚ) / 4) [])], cpRow 1024 ((2, 32, 1, 32) : ℕ × ℕ × ℤ × ℕ)] : List (List (ℕ × CMon))).getD i []
def RL (i : ℕ) : DMon × List ℤ := ([(⟨1, []⟩, []), (⟨0, []⟩, [2, 32, 1, 32])] : List (DMon × List ℤ)).getD i (⟨1, []⟩, [])
def IX : List ℕ := [0, 1]
def FL : List ℕ := [2]
noncomputable def GS (θ : ℝ) : List ℝ := [θ]
noncomputable def ρ (d : ℕ) (θ : ℝ) : Val := ⟨fun i => ([d] : List ℕ).getD i 1, fun j => ex ((GS θ).getD j 0)⟩
def S0 : List STm := [⟨CMon.ofQ 1, [], [], []⟩]
def S1 : List STm := [...]
theorem rhoOK ... : (ρ d θ).OK := ...
theorem flOK ... : FloorOK FL (ρ d θ) := ...
theorem envD ... : ∀ i ∈ IX, (E d θ).dimv i = (DT i).den (ρ d θ).dv := ...
theorem envP ... : ∀ ph j, ex ((E d θ).phase ph j) = phF (ρ d θ) (PT ph) j * rV (ρ d θ) (RL ph) j := ...
theorem scal0 ... : (s0) = exDen (ρ d θ) (fun _ => 0) S0 := ...
\end{lstlisting}
A phase vector in \lexpy!cp_rows! is written \lexlean!cpRow D ((k, m, c, d) : ℕ × ℕ × ℤ × ℕ)!. The
binders are the phase and dimension parameters; each dimension symbol $n$ gets a hypothesis
\lexlean!(h_n : f ≤ n)! with $f$ its floor (at least 1), each mixed-dimension constraint a
hypothesis \lexlean!hc<j>!.

\paragraph{Phase proofs.} \lexlean!envP! is proved per phase vector $p$ by splitting off its rule
(\lexpy!rule_pre!) and then:
\begin{itemize}
\item no rule: \lexlean!rw [show rV ρ (RL p) j = 1 by simp [rV, rEv, RL], mul_one]!; a marking
rule: \lexlean!rw [show RL p = (⟨0, []⟩, [k, m, c, d]) from rfl, rV_mark, mul_one]!; any other
rule through \lexlean!rule_ex!;
\item an empty row: \lexlean!simp [unfold, Φ, PT, phF, ex_zero]!;
\item a row matching the exporter's flat table of at most \lexpy!TABLE_CHUNK! entries
(\lexpy!tab_row!: \lexlean!ptOK! of the table's numerators over its denominator):
\lexlean!exact tab_ph ρ _ den _ (by norm_num) _ (by decide +kernel) j!;
\item a marking rule with a chunked table that matches: \lexlean!exact tab_phC ρ name 256 (by norm_num) den _ (by norm_num) _ (by decide +kernel) ok j!,
\lexlean!ok! being \lexlean!(ptOK_cpRow_of _ D q _ (by decide) (by decide +kernel) (by decide +kernel))!
when the table is \lexlean!cpRow D q! reduced modulo $M$, else \lexlean!(by decide +kernel)!;
\item otherwise a case split on each stored index with \lexlean!phC_den! (or \lexlean!phC_den_int!
when an integer multiple of an integer symbol dropped out), the remaining indices by
\lexlean!tab_zeroZ! or \lexlean!tab_zero!.
\end{itemize}
The scalar lemmas use the \lexlean!simp! set \lexpy!scalar_proof! names (\lexlean!exDen!,
\lexlean!CMon.den_mul!, \lexlean!ofQ_den!, \lexlean!exQ_den!, \lexlean!rtD_den!, \lexlean!dP_den!,
\lexlean!phC_den!, and \lexlean!phV!, \lexlean!ipDen!, \lexlean!imDen!, \lexlean!List.prod_cons! when the
model has a root of unity of symbolic order), then \lexlean!push_cast!, \lexlean!ring_nf! and
\lexlean!om_eq_ex!.

\subsubsection{Check declarations}\label{sec:lean-export-sym-checkdecls}

\lexpy!check_decls(name, doc, nets, s0, s1, s0t, s1t)! writes the proof of the check for
\lexpy!nets = [(tag, text, net), (tag, text, net)]!; the pair without boxes uses
\lexpy!chk_decls(A, B)! (name \lexlean!chk!, tags and texts \lexlean!A!, \lexlean!B!, scalars \lexlean!S0!,
\lexlean!S1!), a family lemma name \lexlean!C<i>! and tags \lexlean!C<i>x!, \lexlean!C<i>y!. Two networks
with the same text share one set of declarations.

\begin{algorithm}[\texttt{check\_decls}]\label{alg:lean-export-sym-checkdecls}
\begin{enumerate}
\item The \emph{whole} form: for \lexlean!chk! with a character certificate,
\lexlean!def FZN : ℕ!, \lexlean!def FZC : FzCert! and
\lexlean!theorem chk : symCheckF FZN FZC FL DT PT RL labs IX A B S0 S1 = true := by decide +kernel!;
for \lexlean!chk! with a range certificate, \lexlean!def RGC : RgCert! and \lexlean!symCheckR RGC ...!;
otherwise \lexlean!theorem name : symCheck FL DT PT RL labs IX A B s0 s1 = true := by decide +kernel!.
\item For \lexlean!chk! after a compact check (\lexpy!rx! nonempty): the \lexpy!rx_decls! form.
\item Otherwise the stepwise form: \lexpy!tr_steps! of each network (failure: the whole form);
the products must be admissible and $\mathit{rest} = \mathit{nf\_ex}$ of the difference empty,
or certified by \lexpy!fz_cert! (for \lexlean!chk! with \lexpy!fz!), or by \lexpy!rg_cert! (for \lexlean!chk!
with \lexpy!rg!); else the whole form.
\end{enumerate}
\end{algorithm}
The stepwise form, inside \lexlean!noncomputable section ... end!, per network tag \lexlean!A!
with $m$ nodes:
\begin{lstlisting}[style=lexlean]
def AH (i : ℕ) : List THint := (([[...], ...] : List (List THint))).getD i []
def An0 : Node := ...
def AD2 : List Node := []
def AD1 : List Node := An1 :: AD2
def AD0 : List Node := An0 :: AD1
def AM0 : ℕ := 0x0
def AF0 : ArchytasZX.Sym.Fac := (⟨nat_lit 2, []⟩, [([nat_lit 1], Int.ofNat (nat_lit 1))])
def AE0 : List STm := ([⟨..., [...], [AF0], []⟩] : List STm)
theorem AB2 : trNodesM (nfH FL AH) DT PT RL labs A.int AM2 2 AD2 = some AE2 := by decide +kernel
set_option maxHeartbeats 0 in
theorem AT0 : trBlockM (nfH FL AH) DT PT RL labs A.int AE2 AM0 0 [An0, An1] = some AE0 := by decide +kernel
theorem AB0 : trNodesM (nfH FL AH) DT PT RL labs A.int AM0 0 AD0 = some AE0 :=
  (trNodesM_block (nfH FL AH) DT PT RL labs A.int AE2 [An0, An1] AD2 AM0 AM2 0 (by decide +kernel) AB2).trans AT0
theorem AR : trNetG (nfH FL AH) DT PT RL labs A = some (exSums (labSums DT labs (A.int.filter fun l => !(labelsOf A.nodes).contains l)) AE0) := ...
def AML : ℕ := 0x6
def AX : List STm := exSums (labSums DT labs []) AE0
theorem AXE : ... = AX := by
  rw [filter_free A.int (labelsOf A.nodes) AML (by decide +kernel)]
  decide +kernel
theorem chkcA : netChk FL DT PT RL labs IX A = true := by decide +kernel
\end{lstlisting}
\lexlean!AH! holds \lexpy!tr_hints!; \lexlean!An<k>! and \lexlean!AD<k>! the nodes and their suffixes. The
block size is $\min(\mathit{BLOCK\_MAX}, \max(1, \lfloor \mathit{BLOCK\_TERMS} / \max_k
|\mathit{steps}_k|\rfloor))$; the marks are its multiples below $m$ and $m$; at each mark $k$,
\lexlean!AM<k>! is the mask of the labels of the nodes before $k$ and \lexlean!AE<k>! the stated
expression of \lexpy!tr_steps! with each distinct factor named \lexlean!AF<j>!. With every label at
dimension index 0 and $0\in ix$, the network check is \lexlean!netChkF FL DT PT RL labs IX AMK A!
with \lexlean!AMK! the mask of the boundary and internal labels, used through \lexlean!netChkF_ok!.
Then \lexlean!chkZH! holds the term hints of the difference, and the parts
\begin{lstlisting}[style=lexlean]
theorem chkp : A.bnd.isPerm B.bnd = true := by decide +kernel
theorem chkk : (exConst S0 && exConst S1) = true := by decide +kernel
theorem chkm : (exMulOK S0 AX && exMulOK (exNeg S1) BX) = true := by decide +kernel
theorem chkz : ((nfH FL chkZH 0) (exMul S0 AX ++ exMul (exNeg S1) BX)).isEmpty = true := by decide +kernel
\end{lstlisting}
the doc comment on the last. A character certificate adds \lexlean!chkFN!, \lexlean!chkFC! and reads
the residual through \lexlean!fzNf chkFN chkFC (nfH FL chkZH 0)! (proof \lexlean!fzNf_den!); a range
certificate adds \lexlean!chkRC!, uses \lexlean!nfHP! for every normalizer, reads the residual through
\lexlean!rgNf chkRC (nfHP FL chkZH 0)! (proof \lexlean!rgNf_den!) and adds
\lexlean!theorem chkr : labOK labs IX DT A.bnd chkRC.2.1 = true!. The stepwise form records the
arguments of its \lexlean!equal_sym! lemma: the two normalizers and their \lexlean!_den! proofs, the
\lexlean!R!/\lexlean!XE! translations, \lexlean!p!, both network checks, \lexlean!k!, \lexlean!m!, the
residual normalizer with its proof and \lexlean!z! (and \lexlean!r!).

\paragraph{The compact form \texttt{rx\_decls}.} Inside \lexlean!noncomputable section ... end!:
per network, \lexlean!AX! is the literal of \lexpy!tr_net_k! (when folded) or \lexpy!tr_net_id!
(\lexcode!a compact network whose products capture variables! when either fails), proved
\lexlean!theorem AR : trNetK DT PT RL labs A = some AX := by rw [← trNetKM_eq]; decide +kernel!
when folded, else by \lexpy!id_blocks!; and the network check as above. The products must be
admissible (\lexcode!a compact export whose products capture variables!); the certificate is
\lexpy!rx_cert_f! (folded) or \lexpy!rx_cert! of the difference (\lexcode!the compact residual does not cancel!
when none). The route $cx$ is \lexlean!X! (\lexpy!rx_cert! without a fourth entry), \lexlean!C!
(\lexpy!rx_cert_c!) or \lexlean!F!; C and F with cuts are \emph{chunked}. The declarations are
\lexlean!chkRR! (the ranges); per residual term $i$ and round $k$: the terms \lexlean!chkT<i>_<k>!,
the round hint \lexlean!chkH<i>_<k> : RHint! (or, chunked, \lexpy!chunk_decls!), and
\begin{lstlisting}[style=lexlean]
set_option maxHeartbeats 0 in
theorem chkS0_0 : normTmX chkRR FL [chkH0_0] chkT0_0 = some chkT0_1 := by decide +kernel
theorem chkN0 : normTmX chkRR FL [chkH0_0, chkH0_1] chkT0_0 = none := by rw [normTmX_step _ _ _ _ _ _ chkS0_0, normTmX_stop _ _ _ _ _ chkS0_1]; try rfl
def chkXH : List (List RHint) := [[chkH0_0, chkH0_1], ...]
theorem chkd : (exMul S0 AX ++ exMul (exNeg S1) BX) = [chkT0_0, ...] := by decide +kernel
theorem chkz : (rxNf chkRR FL chkXH (exMul S0 AX ++ exMul (exNeg S1) BX)).isEmpty = true := by
  rw [chkd]
  unfold rxNf chkXH
  rw [nfGoX_none _ _ _ _ _ _ _ chkN0, ...]
  decide +kernel
\end{lstlisting}
(\lexlean!normTmC!/\lexlean!rxNfC!/\lexlean!nfGoC! and \lexlean!normTmF!/\lexlean!rxNfF!/\lexlean!nfGoF! on the
other routes; a chunked round's \lexlean!S! theorem rewrites by \lexlean!normTm<cx>_one ... chkE<i>_<k>!
first), with the parts \lexlean!chkp!, \lexlean!chkk!, \lexlean!chkm! and
\lexlean!chkr : labOK labs IX DT A.bnd chkRR!. The lemma is \lexlean!equal_symSK! (folded) or
\lexlean!equal_symSR!, with identity normalizers in the unfolded case and the residual proved
through \lexlean!rxNf_den! (or its route's variant).

\paragraph{\texttt{id\_blocks} and \texttt{chunk\_decls}.} \lexpy!id_blocks! proves
\lexlean!AR : trNetG (fun _ e => e) DT PT RL labs A = some AX! like the stepwise form with the
identity normalizer, blocks of \lexpy!TR_BLOCK! nodes and the unnormalized suffix products as
\lexlean!AE<k>!, closing with \lexlean!trNetG_of!, \lexlean!trNodesM_eq!, \lexlean!filter_free! and
\lexlean!decide +kernel!. \lexpy!chunk_decls(name, i, k, c, es, cut, cx)! cuts round $k$ of term $i$
into chunks of \lexpy!RX_STEPS! hints \lexlean!chkC<i>_<k>_<c>! (one empty chunk for no hints), the
terms after each chunk \lexlean!chkU<i>_<k>_<c>! (from \lexpy!cut!), one theorem
\lexlean!chkK<i>_<k>_<c> : elimAll<cx>S chkRR chkC.. t = some u := by decide +kernel! per chunk,
the round hint \lexlean!chkH<i>_<k> : RHint := (nat_lit c, C0 ++ (C1 ++ ...))! and
\lexlean!chkE<i>_<k> : elimAll<cx> chkRR (...) chkT<i>_<k> = last! by \lexlean!elimAll<cx>_append! and
\lexlean!elimAll<cx>_of_S!.

\paragraph{The theorem.} \lexpy!sym(envt, ks, n0, n1, a0, a1, chk, l0, l1)! writes the final
application: with recorded stepwise or compact arguments, the recorded lemma
(\lexlean!equal_symS! by default, \lexlean!equal_symSR!, \lexlean!equal_symSK!) on them; otherwise
\lexlean!equal_symF ... FZN FZC!, \lexlean!equal_symR ... RGC! or \lexlean!equal_sym!, ending with the
check \lexlean!chk!:
\begin{lstlisting}[style=lexlean]
theorem main (d : ℕ) (θ : ℝ) (h_d : 2 ≤ d) :
    Equal (E d θ) s0 A s1 B :=
  equal_sym (E d θ) (ρ d θ) (rhoOK d θ h_d) FL (flOK d θ h_d) DT PT RL labs IX
    (envD d θ h_d) (envP d θ h_d) rfl A B S0 S1 _ _ (scal0 d θ h_d) (scal1 d θ h_d) chk
\end{lstlisting}
With a common factor the proof is \lexlean!Equal.congr_s (Equal.smul c proof) t0 t1!. For a family
the environment is \lexlean!(F args).env (ksOf [...])!, the theorem quantifies over the bare
multiplicities, the declarations comprise the scalar tables, the auxiliary networks and the
family lemmas (in reverse creation order), and the header notes the induction
(\lexcode!by induction on the multiplicity n, one copy peeled per step!, or
\lexcode!at fixed multiplicities, every copy peeled! when no multiplicity stays free), the floors
(\lexcode!assumes the dimension floors n ≥ 2!) and the constraints.

% ---------------------------------------------------------------------------
\section{The \texorpdfstring{$\mathbb{F}_2$}{F2} Mirror}\label{sec:lean-export-f2}

The module \module{archytaszx.semantics.lean_f2} mirrors \file{ArchytasZX/SymF2.lean} and
\file{ArchytasZX/SymF2Bridge.lean} on Python integers. The symbolic exporter uses it for
qubit networks (\lexpy!_SymExport.q_plan!): when it applies, the equality is proved by the
$\mathbb{F}_2$ checker \lexlean!qCheck! instead of the general normalizer. Each function computes
what its Lean namesake computes, in the same order, so the literal states written into the Lean
file agree with the ones Lean computes.

\begin{definition}[Forms and terms over $\mathbb{F}_2$]\label{def:lean-export-f2}
An affine form over $\mathbb{F}_2$ is a \emph{mask}: bit 0 is the constant, bit $v+1$ the
variable $v$ (\lexpy!vb(v) = 1 << (v + 1)!). A term \lexpy!(q, k, dels, phs)! is
\[
  q\,\sqrt2^{\,k}\prod_{d\in\mathit{dels}}[\,d = 0\,]\prod_{(u,w)\in\mathit{phs}}(-1)^{uw},
\]
$q\in\Z$, $k\in\Z$, each delta and each phase factor a form. A state \lexpy!(S, t)! holds the mask
$S$ of the labels summed so far (bit $\ell$ for label $\ell$) and the term, \lexpy!None! when the
term is zero.
\end{definition}

\paragraph{Constants.} \lexpy!CHAIN = 40!: at most 40 steps per \lexlean!qchain! block.
\lexpy!BLOCK_WORK = 6000!: at most 6000 forms visited per block past its first step.

\paragraph{Elementary maps.} \lexpy!tb(m, i)! tests bit $i$. \lexpy!sub(x, v, m)! substitutes the form
$v$ for the variable $x$ in $m$ (XOR of $v$ and $x$'s bit when $x$ occurs). \lexpy!clr(x, m)! clears
$x$.

\paragraph{Node translation (\texttt{tr\_q}).}
\begin{itemize}
\item Z without ports: $(1, 2, [\,], [\,])$, the value 2.
\item Z with ports $\ell_0, \ell_1, \ldots$: $(1, 0, [\ell_i + \ell_0 : i\ge1], [\,])$.
\item X with ports: $(1, 2-|\mathit{ports}|, [\sum_i \ell_i], [\,])$.
\item F with one output $o$ and one input $i$: $(1, -1, [\,], [(o, i)])$.
\item anything else: \lexpy!None!.
\end{itemize}

\begin{algorithm}[\texttt{sum\_x(x, t)}: summing a label]\label{alg:lean-export-f2-sum}
\begin{enumerate}
\item If some delta $d$ reads $x$ (the first one): $v = d + x$; substitute $x\mapsto v$ in the
other deltas, returning \lexpy!None! when one becomes the constant 1 and dropping those that become
0 (the kept list is reversed); substitute in both forms of every phase factor.
\item Otherwise, the character sum: with $u', w'$ the forms with $x$ cleared,
$A = \sum_{(u,w)} \big([x\in u][x\in w] + [x\in u]\,w' + [x\in w]\,u'\big)$; the phase factors
become $(u', w')$, dropping those with a zero form. $A = 1$ gives \lexpy!None!; otherwise the
result has $k+2$, and $A$ prepended to the deltas when $A\ne 0$.
\end{enumerate}
\end{algorithm}

\paragraph{Runs.} \lexpy!hints(nodes, internal)! lists, per node, the internal labels whose last
occurrence is at that node (in internal-label order). \lexpy!sum_ls(mask, labs, s)! sums each label,
returning \lexpy!None! when a label is not in the mask or already summed. \lexpy!step_q(mask, node, h, s)! requires that no port of the node is summed, multiplies its translation in
($q$ multiplied, $k$ added, lists concatenated with the node's first) and sums the hints $h$.
\lexpy!mask_of(labs)! sets the labels' bits.

\begin{algorithm}[\texttt{run\_q} and \texttt{plan}]\label{alg:lean-export-f2-run}
\lexpy!run_q(nodes, internal, hs)! starts at $(0, (1, 0, [\,], [\,]))$. Per node it charges the
work $w = 1 + (|\mathit{dels}|+|\mathit{phs}|)\cdot\max(1, |h|)$ (1 for a zero term); it closes
the current block before the node when the block holds \lexpy!CHAIN! steps or the work would exceed
\lexpy!BLOCK_WORK!. A failed step returns \lexpy!False!. After the last node the remaining internal
labels are summed (\lexlean!finQ!). The result is (final term, block sizes, total work).
\lexpy!plan(nodes, internal)! runs the nodes in listed and in reversed order and returns
\lexpy!(reversed, hints, result)! for the order with less total work (ties keep the listed order);
\lexpy!None! when neither passes.
\end{algorithm}

\paragraph{Comparison.} \lexpy!canon_q(t)! puts the deltas into reduced echelon form (pivot = highest
set bit; a delta reducing to 1 gives \lexpy!None!, to 0 is dropped; rows sorted decreasingly),
reduces both forms of every phase factor by the rows, orders each pair, drops factors with a
zero form, counts factors $(1,1)$ as signs, and sorts the rest. \lexpy!scal_eq(a0, a1, e)! decides
$a_0\sqrt2^{\,e} = a_1$ for even $e$ (false for odd $e$). \lexpy!cmp_q(s0, s1, t0, t1)! with scalars
$(a, b, K)$ meaning $\tfrac{a}{b}\sqrt2^{K}$: false when a denominator is 0; both zero terms
compare equal; otherwise the canonical deltas and phases must be equal and
\lexpy!scal_eq(a0 q0 b1, a1 q1 b0, (k0 + k0') - (k1 + k1'))!. \lexpy!scalar_of(model)! returns
$(a, b, K)$ when the model is a product of rational powers and integer powers of $\sqrt2$ and
$1/\sqrt2$, else \lexpy!None!.

\paragraph{When the $\mathbb{F}_2$ plan applies (\texttt{q\_plan}).} No multiplicities, constraints,
parameters or common factor; no opaque scalars; the only dimension is the constant 2
(\lexlean!DT = [⟨2, []⟩]!), no stored phase entry and no rule; every label has dimension index 0;
both networks convert; sorted boundaries agree; every node has dimension index 0 and a \lexpy!tr_q!
translation; both scalars have the form $\tfrac ab\sqrt2^{K}$; both plans exist; and \lexpy!cmp_q!
holds on the final terms.

\paragraph{Emitted proof (\texttt{q\_source}).} Under \lexlean!open ArchytasZX.SymF2!, per side
$t\in\{0,1\}$: \lexlean!H<t>! (the hints, \lexpy!lean_hints!), \lexlean!R<t> := (ordQ rev N<t>).zip H<t>!,
the summed-label mask \lexlean!IM<t>! in hexadecimal with \lexlean!IM<t>_eq : maskL N<t>.int = IM<t>!
by \lexlean!decide +kernel!, the block sizes \lexlean!KS<t>!, and the command
\lexlean!qchain q<t> IM<t> R<t> KS<t>! (under \lexcode!set_option maxHeartbeats 0 in!). Then
\lexlean!hph! (every phase is zero), \lexlean!hL! (every label has dimension index 0), \lexlean!qc : qCheck N0 N1 H0 H1 q0_fin q1_fin a0 b0 k0 a1 b1 k1 = true! by \lexlean!decide +kernel!, and
\begin{lstlisting}[style=lexlean]
theorem main :
    Equal E s0 N0 s1 N1 :=
  Equal.congr_s (equal_q E rfl hph hL N0 N1 false true H0 H1
    q0_fin q1_fin (1 : ℤ) 1 (0 : ℤ) (1 : ℤ) 1 (0 : ℤ) (by rw [IM0_eq]; exact q0) (by rw [IM1_eq]; exact q1) qc)
    (show (((1 : ℤ) : ℂ) / ((1 : ℕ) : ℂ) * rt 2 ^ (0 : ℤ)) = s0 by qscal) (show ... = s1 by qscal)
\end{lstlisting}
The diagram literals are bridged by packed wires: \lexpy!pkey(width, node, out, pos)! is
$(2\,\mathit{node} + [\lnot\mathit{out}])\cdot\mathit{width} + \mathit{pos}$, and
\lexpy!pack_w(width, bits, wires)! folds the wires from the last, each as
$\mathit{key}(a) + 2^{\mathit{bits}}\mathit{key}(b) + 2^{2\,\mathit{bits}}\mathit{acc}$, with
\lexcode!width! the largest node arity (at least 1) and \lexcode!bits! the largest key bit length. The
file then states \lexlean!WB<n> : ℕ := 0x...!, \lexlean!D<n>_bq : bridgeQ width bits off top WB<n> D<n> N<n> = true! by \lexlean!decide +kernel!, and derives \lexlean!D<n>_net! and \lexlean!D<n>_chk! from
\lexlean!bridge_ok!. The label table \lexlean!labs! is replaced by
\lexlean!(List.range top).map fun l => (l, 0)! with \lexcode!top! one more than the largest label.

% ---------------------------------------------------------------------------
\section{Qudit Digit Networks}\label{sec:lean-export-qdigits}

The module \module{archytaszx.semantics.lean_qdigits} proves \lexlean!Equal! for networks over
the powers of one odd prime $p$ with at least one port of dimension $p^{k}$, $k>1$, over
\file{ArchytasZX/PathSumQuditDigits.lean}. A label of dimension $p^{k}$ is $k$ digit
variables in $\Z_p$, low digit first; a label of dimension $p$ is its own digit variable. Z
spiders on $p^{k}$, dimension binders and splitters become digit constraints; Z, X and F on
dimension $p$ become the fragments of the qudit network export
(\cref{sec:lean-export-pathsum}).

\paragraph{Constant.} \lexpy!LABEL_OFFSET = 1 << 19!: the first label of a summed wire of more
than one digit. The left side uses \lexpy!LABEL_OFFSET!, the right \lexpy!LABEL_OFFSET + (1 << 17)!.

\begin{algorithm}[\texttt{\_side(diagram, phases, sc, label0)}]\label{alg:lean-export-qdigits-side}
\begin{enumerate}
\item The raw sum must be a qudit sum (\lexcode!the qudit digit network export covers odd prime path sums only!) without atoms (\lexcode!the qudit digit network export has no atoms!).
\item Boundary ports (outputs, then inputs) take as label their first digit variable, numbering
digits consecutively from 0; each digit maps to its class variable.
\item A summed wire's digits are the class variables when those are unused path variables, else
fresh variables from \lexpy!QPIN_OFFSET! $=2^{18}$. A one-digit wire is labeled by its digit, a
$k$-digit wire by $\mathit{label0}+w$. A wire on a boundary port raises \lexcode!a wire on a boundary port!.
\item Only Z, X, F, B and S nodes (\lexcode!node <id> is a <name>!), none weighted
(\lexcode!node <id> carries weights!), none legless (\lexcode!node <id> has no legs!).
\item Hyperedges: a Z spider joins per digit $j$ the $j$-th digits of its labels (one digit count,
else \lexcode!node <id> joins legs of different dimensions!); B and S join whole and parts as in
\cref{alg:lean-export-net-side}; an X or F with a port of more than one digit raises
\lexcode!node <id> is a X on dimension p^k, k > 1!. Anchors by breadth-first search as before
(\lexcode!a digit constraint not reached from its class's variable!). X spiders take the last
$|X|$ sorted path variables.
\item Fragments: B/S \lexlean!bFragQ p digT whole a b! with \lexlean!b_frag_q! or \lexlean!s_frag_q! and
the digit constraints; F \lexlean!fFragQ p out in! (proof \lexlean!f_frag_q! lifted by
\lexlean!lift_q!; a self-loop raises \lexcode!node <id> is an F on a self-loop!); Z with $k>1$
\lexlean!zFragQD p digT anchors legs poly! with \lexlean!z_frag_qd! and constraints from each port's
digits to the anchor digits; Z with $k=1$ \lexlean!zFragQ!; X \lexlean!xFragQ! (both lifted). The
phase tables come from \lexpy!_qudit_table! and \lexpy!_turns_table! over $p^{k}$ values, entries on
the $1/p$ grid (\lexcode!node <id> has a phase <t> off the 1/<p> grid!).
\item Pins: for each reached digit $x\ge$ \lexpy!QPIN_OFFSET!, deepest first, a constraint setting
$x$ to its parent (\lexcode!no constraint sets the digit <x> to <a>!) gives \lexlean!.perm [x, path]!,
\lexlean!.pin a!. Then \lexlean!.norm!, the path check (\lexcode!the pinned network keeps other path variables!), and the reduction steps of the pinned sum. The scalar must model exactly.
\end{enumerate}
\end{algorithm}

\begin{signature}
\lexpy!export_qudit_digits_net(left: Diagram, right: Diagram, name: str = "main") -> str!
\end{signature}
After building both sides it requires a multi-digit label (\lexcode!no leg of dimension p^k, k > 1!), one prime and boundary (\lexcode!the two path sums have different bases or boundaries!), equal
boundary labels (\lexcode!the two networks have different boundary labels!), one digit count per
label (\lexcode!the label <l> has two dimensions!), equal boundary digits (\lexcode!the two networks have different boundary digits!), one canonical form and one elimination order. The constants'
identity uses $n=12$ for $p=3$ and $n=p$ otherwise, exactly as \lexpy!qnet_body!. The file, under
\lexlean!import ArchytasZX.PathSumQuditDigits!, namespace
\lexlean!ArchytasZX.PathSumQuditDigitsProof!, \lexpy!QNET_OPTIONS! and the prime instance, has the
environment \lexpy!net_env(phases, labs, "p ^ k", p)!, per side \lexlean!ownT : PSem.Own!,
\lexlean!labsT!, \lexlean!digT!, \lexlean!ownTOK!, \lexlean!NT!, \lexlean!psT : List (Node × QFrag)!,
\lexlean!netT : QNetPSD E NT p ownT labsT digT psT! and \lexlean!cT!, the step lists, the \lexlean!lg!,
\lexlean!qg! and \lexlean!fac! theorems, \lexlean!consts!, and
\begin{lstlisting}[style=lexlean]
theorem main : Equal E (sL) NL (sR) NR := by
  apply equal_qpsd2 _ _ netL netR [0, 1] (by decide) (by decide) rfl
  intro τ
  rw [← QSem.val_smul, ← QSem.val_smul,
    QSem.run_sound stepsL (qnetDataD digL NL (psL.map Prod.snd)) _ τ (by decide),
    QSem.run_sound stepsR (qnetDataD digR NR (psR.map Prod.snd)) _ τ (by decide)]
  apply QSem.final_sound _ _ _ _ τ kL kR (by decide) (by decide) (by decide)
  rw [facL, facR, cL, cR]
  all_goals exact consts
\end{lstlisting}
The list after \lexlean!netR! holds the boundary digit variables. Unlike the other path-sum
exports, the runs are checked by one \lexlean!decide! each, without \lexlean!qrun_chain! checkpoints.

% ---------------------------------------------------------------------------
\section{Refutations}\label{sec:lean-export-refute}

A refutation states that two diagrams are \emph{not} equal for every value of their
parameters and proves it at one point. The module
\module{archytaszx.semantics.lean_refute} writes exact refutations: at the point each network's
entry is, by the Lean lemma \lexlean!eval_W!, a product of $(\sqrt d)^{-1}$ factors and
coefficient denominators times a formal sum in $\zeta = \omega_M$ (\lexlean!om M!);
\lexlean!eval_res! gives that sum modulo a prime $p \equiv 1 \pmod M$ at an element $w$ of order
$M$ by a checked contraction of residue tables, and \lexlean!refute_res! derives that the scaled
entries differ from their differing images in $\Zd{p}$.
\module{archytaszx.semantics.lean_numeric} writes refutations by ball (midpoint--radius
interval) arithmetic~\cite{vanderhoeven2009ball} when the values at the point leave
$\Q(\zeta_M)$. Both build the Lean networks through the network exporter state \lexpy!_Export! of
\module{archytaszx.semantics.lean_net} (\cref{ch:lean-export-net}), both raise
\lexpy!NetUnsupported!, and both re-check their evidence through
\module{archytaszx.semantics.refute_check} before returning a source. The public names of
\module{archytaszx.semantics.lean_refute} are \lexpy!export_refutation! and
\lexpy!export_interface_refutation!.

\subsection{Constants}\label{sec:lean-export-refute-constants}

\begin{center}
\begin{tabular}{@{}p{0.28\textwidth}p{0.17\textwidth}p{0.47\textwidth}@{}}
\toprule
Name & Value & Meaning \\
\midrule
\lexpy!MAX_ASSIGNMENTS! & 4096 & boundary assignments tried with whole-network residue plans, and by the numeric refuter \\
\lexpy!MAX_POINTS! & 16 & default sample points tried after the given point \\
\lexpy!BRUTE_LIMIT! & 256 & largest product of summed-label dimensions the group evaluator sums directly \\
\lexpy!MAX_TERMS! & 200\,000 & bound on node dimensions, table sizes and summation work \\
\lexpy!MAX_CYCLOTOMIC_DEGREE! & 512 & largest $\varphi(M)$ (or $\varphi(m)$ of a subfield) \\
\lexpy!MAX_SPARSE! & 64 & most entries of a sparse phase vector \\
\lexpy!MAX_PRIMES! & 4 & primes $p \equiv 1 \pmod M$ tried \\
\lexpy!MAX_PINNED! & 16 & boundary assignments tried with pinned residue plans \\
\lexpy!SYM_FIRST! & 10\,000 & widest node table above which the substitution contraction is tried first \\
\lexpy!ELIM_CHUNK! & 2\,000 & most entry sums per chunk theorem of an elimination \\
\lexpy!MAX_RES_ENTRIES! & 2\,000\,000 & most residue-table entries of one plan \\
\lexpy!RES_PRIME_FLOOR! & $2^{20} = 1\,048\,576$ & every prime tried exceeds it \\
\lexpy!LEAF_X_MIN! & $2^{14} = 16\,384$ & least table size times dimension for an X leaf to be checked by its row \\
\lexpy!ELIM_DESCENT! & \lexpy!True! & eliminations are proved by \lexlean!elimT_ok! when cheap \\
\lexpy!DESC_MAX! & 300\,000 & most cursor steps of an \lexlean!elimT_ok! descent \\
\lexpy!PT_DENSE_MIN! & 16 & least key count of a phase vector written as a dense \lexlean!Pt! list \\
\lexpy!_Refuter.MAX_TABLE! & 200\,000 & largest trial table of a residue elimination (\lexpy!MAX_TERMS!) \\
\lexpy!K! & \lexcode!(by decide +kernel)! & the kernel-decision proof text \\
\lexpy!DIGITS! & 40 & decimal digits of ball midpoints and radii \\
\lexpy!TAYLOR! & 10 & Taylor terms of $\cos$ and $\sin$ \\
\lexpy!SQRT_DIGITS! & 45 & decimal digits of square-root balls \\
\bottomrule
\end{tabular}
\end{center}
\lexcode!mpmath.mp.dps! is set to 60 for the numeric evaluation (\tool{mpmath}{https://mpmath.org}).
\lexpy!TABLE_CHUNK! (256) of \module{archytaszx.semantics.lean_net} bounds the phase tables written
inline.

\subsection{Dispatch \texttt{export\_refutation}}\label{sec:lean-export-refute-dispatch}

\begin{signature}
\lexpy!export_refutation(left, right, point, name="main", describe=None, index=None) -> str!
\end{signature}
\lexcode!index!, a boundary value per slot, is the first assignment tried at \lexcode!point!.
\begin{algorithm}[Refutation dispatch]\label{alg:lean-export-refute-dispatch}
\begin{enumerate}
\item If a \bbox of either diagram is extended, or a node phase reads a \bbox multiplicity
(\lexpy!_mult_phases!), refute the families instantiated at a point
(\lexpy!_instantiated_refutation!, passing \lexcode!index!).
\item If the interfaces differ (\lexpy!_interfaces_differ!: different slot counts, or a slot whose
dimension or port count differs symbolically) and the diagrams have \bboxes or a slot dimension
differs, try the interface refutation; on failure without \bboxes re-raise, with \bboxes fall
back to the instantiated refutation.
\item Candidate points: the given point, then the default sample grid
(\lexpy!sample_grid! of \module{archytaszx.semantics.decide}; an exception there leaves the grid
empty) ordered by the sum of absolute values (\lexpy!_by_size!; a non-numeric value counts 1,
ties keep their order), first \lexpy!MAX_POINTS!, duplicates removed in order.
\item For each candidate: build a \lexpy!_Refuter! (its \lexpy!hint! is \lexcode!index! at the first
candidate), run it (\cref{alg:lean-export-refute-run}) and return
\lexpy!checked_residue(src, left, right, candidate, describe, refuter)!
(\cref{sec:lean-export-refute-recheck}). If the run or the recheck raises
\lexpy!ScalarUnsupported! or \lexpy!NetUnsupported! and \lexpy!refuter.wide()! holds, build a second
refuter with the same hint and \lexpy!sparse_all = True! and return its checked source; a
\lexpy!RecheckFailed! of that second try is remembered, any other failure of it is dropped, and the
first try's exception is re-raised. A raised exception is remembered as the first failure (a
\lexpy!ScalarUnsupported! becomes a \lexpy!NetUnsupported! with its message); a \lexpy!RecheckFailed! is
also remembered as the first recheck failure.
\item For each candidate, return \lexpy!export_numeric_refutation! when it succeeds; a
\lexpy!RecheckFailed! is remembered as above, any other \lexpy!NetUnsupported! is dropped.
\item Raise the first recheck failure, else the first failure.
\end{enumerate}
\end{algorithm}
\lexpy!wide()! holds when the order $M$ is chosen and either some sparse phase vector is recorded
or $\varphi(M) > \mathtt{MAX\_CYCLOTOMIC\_DEGREE}$; before $M$ is chosen it holds when the order
choice failed on a phased node wider than \lexpy!MAX_TERMS! (\lexpy!wide_node!).

\paragraph{Interface refutations.} \lexpy!_legs(d)! lists each boundary slot's dimension and port
count (the product of the multiplicities of the \bboxes containing the slot), outputs then
inputs. \lexpy!export_interface_refutation(left, right, point, name="main")! tries the given point
and the sample grid (first \lexpy!MAX_POINTS!); at a point giving every symbol a value it flattens
both lists to dimension lists (dimension positive, count nonnegative, both integers) and, when
they differ, emits
\begin{lstlisting}[style=lexlean]
import ArchytasZX
namespace ArchytasZX.Export
theorem main :
    ¬ ∀ (n : ℕ) (h_n : 0 < n), (List.replicate n 2) = (List.replicate 2 n) := by
  intro h; exact absurd (h 3 (by norm_num)) (by decide)
end ArchytasZX.Export
\end{lstlisting}
(\lexcode!decide! alone without symbols). Symbol names are sanitized to \lexcode![A-Za-z0-9_]!, with
\lexcode!v_! prefixed to a leading digit; a dimension or count that is not a sum, product or power
of naturals and symbols raises \lexcode!<e> is not a natural number expression!. No point separating
the lists raises \lexcode!the leg dimensions agree at every sample point!.

\paragraph{Instantiated families.} For each candidate (the point, then the first
\lexpy!MAX_POINTS! of the size-ordered grid) giving every multiplicity a value, each diagram is
instantiated (\lexpy!instantiate_symbol! for each multiplicity its \bboxes read, substitution of the
multiplicities its scalar, phases or port dimensions still carry, \lexpy!expand_concrete_boxes!), a
closed scalar holding $\exp$ is replaced by its cosine form when that is equal and free of
$\exp$ and $\cos$ (\lexpy!_close_scalar!), a pair with \bboxes left is skipped, and
\lexpy!export_refutation! is called on the pair with the remaining point and, at the given point
only, with \lexcode!index!. The header gains the line \lexcode!the families instantiated at n = 3!.
Failure everywhere raises the first error or \lexcode!no sample point instantiates both families!.

\subsection{The exact refuter}\label{sec:lean-export-refute-exact}

\begin{algorithm}[\texttt{\_Refuter.run}]\label{alg:lean-export-refute-run}
\begin{enumerate}
\item Construction labels both networks (the right one on fresh labels when the boundary sizes
differ; within a family that raises \lexcode!the boundaries differ in size within a family!),
registers them and models both scalars; an opaque scalar raises
\lexcode!a diagram scalar has no exact Lean form!.
\item Different boundary sizes: emit \lexlean!refute_bnd (by decide) (h args)! with every
dimension parameter at 1, every other parameter at 0 and every multiplicity at 0; then
\lexlean!main_diagrams : ∀ binders, ¬ DEqual E s₀ D₀ s₁ D₁! over Lean literals of both
diagrams (the right one labelled again from 0), proved by
\lexlean!exact absurd h.1 (by decide)! on the boundary dimension lists; the header is
\lexcode!refutation: the boundaries differ in size!.
\item \emph{Environment} (\lexpy!setup_env!): a weighted node (kind \lexcode!ZW!, \lexcode!XW!, or a
bundle of one) raises \lexcode!a coefficient spider has no cyclotomic entry! when \lexpy!weights_ok! is
false (the numeric refuter) and \lexcode!a coefficient spider in a family! in a family; dimensions
and multiplicities take their values at the point (each must be a natural number, else
\lexcode!<e> = <v> is not a natural number!); a bundled dimension is its root dimension to the
product of its chain's multiplicities.
\item \emph{Order} $M$ (\lexpy!choose_m!, \cref{alg:lean-export-refute-order}).
\item Different boundary label sets raise \lexcode!the diagrams have different boundaries!.
\item \emph{Scale factors}: per network the quadruples $(d, c, q, k)$ of \lexpy!sqw_list!
(\cref{sec:lean-export-refute-weighted}) and $K_W = \prod (\sqrt d)^{-c} q^{-k}$ (\lexpy!_kw!).
With \lexpy!sparse_all!, $\lambda$ is the root factor of the pair $(s_0K_{W,0}, s_1K_{W,1})$
(\lexpy!_root_factor!); otherwise $\lambda = 1$.
\item $f_j = \mathrm{poly}(\lambda s_jK_{W,j})$ (\cref{sec:lean-export-refute-arith}); $\ell$ is
the lcm of the denominators of the coefficients of $f_0$ and $f_1$; $v_j$ is the coefficient
list of $\ell f_j$, constant first, of length $\deg f_j + 1$ (empty for $f_j = 0$;
\lexpy!_coeff_list!).
\item For each of the first \lexpy!MAX_PRIMES! pairs $(p, w)$ of \lexpy!_primes(M)!: set
$\hat f_j = \sum_i v_{j,i} w^i \bmod p$ and search a boundary assignment
(\lexpy!find_sigma!, \cref{alg:lean-export-refute-sigma}); stop at the first success. No success
raises \lexcode!the networks agree mod p at every assignment checked!.
\item Record the assignment $\sigma$ and the two residues (\lexpy!sigma!, \lexpy!residues!), clear
the certificate tables (\lexpy!init_cert!), and prove for each side
$\ell\lambda s_jK_{W,j} - \sum_i v_{j,i}\zeta^i = 0$ by \lexpy!combination!
(\cref{alg:lean-export-refute-combination}).
\item Write the source (\lexpy!source_res!, \cref{sec:lean-export-refute-source}).
\end{enumerate}
\end{algorithm}

\begin{definition}[Root factor and rational values]\label{def:lean-export-refute-root}
\lexpy!_rational(e)! is the rational value of $e$ when $e$ holds no $\omega_m$, $\exp$ or $i$ atom
and, after $\sqrt x \mapsto \sqrt x$, $x^{-1} \mapsto 1/x$, numerals to their values and
simplification, is rational; else \lexpy!None!. \lexpy!_root_factor(es)! takes the first $e$ of
\lexcode!es! that is not rational: when $e^2$ is not a positive rational $a/b$ it returns 1;
otherwise it returns $\sqrt x$ for the square-free part $x$ of $ab$ when $\sqrt x\, e$ is
rational for every $e$ of \lexcode!es!, and 1 otherwise. With every $e$ rational it returns 1.
\end{definition}

\subsubsection{The order of \texorpdfstring{$\zeta$}{zeta}}

\begin{algorithm}[\texttt{choose\_m}]\label{alg:lean-export-refute-order}
The needed orders start as $\{1\}$. For every node of both networks except binders B and
splitters S:
\begin{enumerate}
\item $d$ is the node dimension; $d = 0$ raises \lexcode!a node of dimension 0!; add $d$. A node
that reads no phase (\lexpy!_phased!: kinds Z, X, ZB, XB, ZW, XW and bundles of them) stops here.
\item With \lexpy!sparse_all!, a node that is not weighted and whose phase vector is not zero
(\lexpy!zero_vector!) takes its sparse vector (\lexpy!sparse_of!: \lexpy!None! when the vector has an
index rule, the export is a family, the vector has a phase table or digit polynomial, or it has
more than \lexpy!MAX_SPARSE! entries; else index $\mapsto$ phase at the point); when there is one,
record it and add the denominators of its entries.
\item A digit-polynomial phase vector (\lexpy!phase_polys!, outside a family, not weighted): record
it and add its denominator.
\item $d > \mathtt{MAX\_TERMS}$ sets \lexpy!wide_node! and raises
\lexcode!a node of dimension <d>, above 200000!.
\item A weighted node records the coefficient keys $(ph, j)$, $j < d$ (\lexpy!ckeys!).
\item A zero phase vector is recorded in \lexpy!zs!.
\item A phase table $(nums, den)$ (\lexpy!phase_tables!, outside a family) is recorded in \lexpy!tabs!
and $den$ is added.
\item Otherwise each key $(ph, j)$, $j < d$, is recorded once (\lexpy!keys!) with its phase at the
point in turns (\lexpy!phase!: the entry plus the index rule at $j \ne 0$; a value that is not
rational raises \lexcode!the phase <v> is not a rational number of turns!), and its denominator is
added.
\end{enumerate}
Then the atoms of both scalars at the point, of every coefficient entry and of $\sqrt d$ for
every node with a nonzero square-root count add their orders (\lexpy!atom_needs!): 4 for $i$, $m$
for $\omega_m$, and for each $\sqrt x$ with square-free part $s$, per prime $q \mid s$: 8 for
$q = 2$, 12 for $q = 3$, $q$ for $q \equiv 1 \pmod 4$, $4q$ otherwise; an \lexpy!Ex! raises
\lexcode!the exponential <a> at the counterexample!; a symbol without a value raises
\lexcode!the counterexample gives no value for <s>!. $M$ is the lcm. The phase table maps each
recorded key to $Mt \bmod M$; the coefficient tables follow
(\cref{sec:lean-export-refute-weighted}).
\end{algorithm}
The power of $\zeta$ of entry $k$ of phase vector $ph$ (\lexpy!P!) is read, in this order, from the
digit polynomial ($\mathrm{value}(k)\cdot M/den \bmod M$), the phase table
($nums_k \cdot M/den \bmod M$, 0 past the table), the sparse vector ($tM \bmod M$, 0 for an
absent index), and the key table (0 for an absent key); the atom part is always empty.

\subsubsection{Weighted spiders}\label{sec:lean-export-refute-weighted}

A weighted spider (\lexcode!ZW!, \lexcode!XW!) carries a coefficient vector. \lexpy!coef_at(ph, k)! is 1
at $k = 0$; otherwise it reads the coefficient models of the vector's body
(\lexpy!weight_tab!, \lexpy!coef_models! of the export), a missing index giving 1, at the point; a
vector without models raises \lexcode!a coefficient vector without concrete entries!. With
$c_{ph,k} = \mathrm{poly}(\mathrm{coef\_at}(ph,k))$, the denominator $q_{ph}$ (\lexpy!qtab!) is the
lcm of the denominators of every coefficient of every $c_{ph,k}$, and \lexpy!ctab[(ph, k)]! is
$q_{ph}c_{ph,k}$ as an integer vector over powers of $\zeta$. The tensor of \lexcode!ZW! is
$\sum_{k<d}[\text{all ports} = k]\,\mathrm{ctab}(ph,k)$; that of \lexcode!XW! is
$\sum_{k<d}\mathrm{ctab}(ph,k)\,\zeta^{(M/d)(k\sum o + (d-1)k\sum i)}$.

\lexpy!qc(kind)! (mirror of \lexlean!qcW!) is 1 for \lexcode!ZW! and \lexcode!XW!, the multiplicity times the
inner count for a bundle, 0 otherwise. \lexpy!sqw_list(net)! (mirror of \lexlean!sqLW!) gives per node
$(d, c, q, k)$: $d$ and $c$ from \lexpy!sq_list! (\cref{sec:lean-export-refute-groups}),
$q = q_{ph}$ (1 when absent), and $k = 0$ for B and S, else $\mathrm{qc}$.

\subsubsection{Exact arithmetic in \texorpdfstring{$\Q(\zeta)$}{Q(zeta)}}\label{sec:lean-export-refute-arith}

\lexpy!poly(e)! returns the rational value of $e$ as a constant when \lexpy!_rational! finds one;
otherwise \lexpy!reduced(poly_expr(e))!. A \emph{big field} (\lexpy!big_field!) is a
\lexpy!sparse_all! run with $\varphi(M) > \mathtt{MAX\_CYCLOTOMIC\_DEGREE}$; there the
\lexpy!modulus! is $\zeta^M - 1$, otherwise it is $\Phi_M$ (\lexpy!phi!, which raises
\lexcode!the cyclotomic field of order <M> has degree above 512! when $\varphi(M) > 512$).
\begin{itemize}
\item \lexpy!poly_expr! is \lexpy!dense_expr! (sympy polynomials reduced by the modulus after each
sum and product) outside a big field, and \lexpy!_poly_of(sparse_expr(e))! in one.
\item \lexpy!sparse_expr(e)! maps exponents in $[0, M)$ to nonzero \lexpy!Fraction! coefficients:
rationals are constants; $\zeta$ and $\omega_M$ are $\{1 \bmod M: 1\}$; sums add; products
multiply by \lexpy!_cyclic_mul! (the product modulo $\zeta^M - 1$, zero coefficients dropped);
integer powers by square and multiply, a negative power inverting the base first; numerals are
constants; an inverse of a nonzero value is \lexpy!inverse!; any other atom is rewritten by
\lexpy!substitute! (without emitting facts), and an atom it cannot rewrite raises
\lexcode!no exact value for <e>!.
\item \lexpy!dense_expr! follows the same rules on sympy polynomials, powers by \lexpy!power!.
\item \lexpy!subfield(f)! returns $(k, m, g)$ with $k = \gcd(M, \text{exponents of } f)$,
$m = M/k$ and $g(y) = \sum c\,y^{e/k} \bmod \Phi_m(y)$, so that $f(\zeta) = g(\zeta^k)$;
$\varphi(m) > 512$ raises \lexcode!the cyclotomic field of order <m> has degree above 512!.
\lexpy!lift(k, g)! is $g(\zeta^k)$.
\item \lexpy!reduced(f)! is $f \bmod \Phi_M$, in a big field $\mathrm{lift}(\mathrm{subfield}(f))$.
\item \lexpy!inverse(f)! is $f^{-1} \bmod \Phi_M$ by \lexpy!sp.invert!, in a big field the inverse of
$g$ modulo $\Phi_m$ lifted by $\zeta^k$.
\item \lexpy!power(b, n)! is $b^n$ modulo the modulus, in a big field by square and multiply with
\lexpy!_cyclic_mul!.
\end{itemize}
\lexpy!substitute! rewrites atoms: $i = \zeta^{M/4}$ (\lexlean!I_sub!); $\omega_m = \zeta^{M/m}$ for
$m \ne M$ (\lexlean!om_sub!); $\sqrt1 = 1$ (\lexlean!rt_one!); $\sqrt{k^2s} = k\sqrt s$
(\lexlean!rt_split!); $\sqrt2 = \zeta^{M/8}+\zeta^{7M/8}$ (\lexlean!rt_sub2!);
$\sqrt3 = \zeta^{M/12}+\zeta^{11M/12}$ (\lexlean!rt_sub3!); $\sqrt{s} = \sqrt q\sqrt{s/q}$ for
composite $s$ and its least prime $q$ (\lexlean!rt_prod!); $\sqrt p$ for an odd prime by the Gauss
sum~\cite{berndt1998gauss} $h = \sum_{a<p} \left(\frac ap\right)\zeta^{aM/p}$, times
$\zeta^{M/4}$ when $p\equiv 3\pmod 4$, negated when its numeric real part is negative, proved by
\lexlean!sqrt_of_sq! with $h^2 = p$ and a ball showing $\mathrm{Re}\,h > 0$ (\lexpy!ball_tree!;
failure raises \lexcode!the Gauss sum for √<p> is not separated from zero!).

\begin{algorithm}[\texttt{combination(target, hyps)}]\label{alg:lean-export-refute-combination}
Returns \lexcode!linear_combination ...! proving $\mathit{target}=0$.
\begin{enumerate}
\item Subtract the hypotheses (each \lexcode!lhs - rhs! with a coefficient).
\item Up to 64 rounds: reduce every $\sqrt x$ and $(\sqrt x)^{-1}$ atom to degree below 2 by
$a^2 = x$ or $a^2 = 1/x$ (facts \lexlean!hsq<k>!); take the least remaining atom by name,
declare its fact \lexlean!hf<k> : atom = rhs! (recursively for atoms of the right side; an inverse
by \lexlean!inv_eq_of_mul_eq_one_right!), substitute, and record the quotient. A remainder raises
\lexcode!an atom substitution left a remainder!; 64 rounds raise
\lexcode!atom substitution did not terminate!.
\item Reduce declared inverse pairs $x\,x^{-1}=1$. When a remainder is left, declare the
relations (\lexpy!relations!) once: \lexlean!hzM : om M ^ M = 1! and, for each prime $p \mid M$,
\lexlean!hg<p>! ($\sum_{j<p}\zeta^{jM/p}=0$, by \lexlean!om_geom!), with the B\'ezout coefficients
writing $\Phi_M$ over them (an extended-gcd chain; a normalized gcd other than $\Phi_M$ raises
\lexcode!the relations of ζ do not generate Φ_M!); a big field skips the B\'ezout
coefficients.
\item Divide by $\zeta^M - 1$ (\lexpy!_div_cyclic!: a term $c\zeta^{aM+b}$, $0 \le b < M$, gives
the quotient $c\zeta^b\sum_{j<a}\zeta^{jM}$ and the remainder $c\zeta^b$; sympy division when an
exponent is not a nonnegative integer or a coefficient holds $\zeta$) and record the quotient
against \lexlean!hzM!.
\item Outside a big field divide the rest by $\Phi_M$ (a remainder raises
\lexcode!the scalar identity does not hold in ℚ(ζ)!) and write the quotient over the
relations by the B\'ezout coefficients. In a big field use \lexpy!sub_terms!
(\cref{alg:lean-export-refute-subterms}).
\end{enumerate}
\end{algorithm}

\begin{algorithm}[\texttt{sub\_terms(rest)}]\label{alg:lean-export-refute-subterms}
\begin{enumerate}
\item $(k, m, g) = \mathrm{subfield}(\mathit{rest})$; $g \ne 0$ raises
\lexcode!the scalar identity does not hold in ℚ(ζ)!.
\item $R(y) = \sum c\,y^{e/k}$ over the terms $c\zeta^e$ of \lexcode!rest!; the quotient is
$R / \Phi_m$.
\item The relations are $y^m - 1$ (named \lexlean!hzM!) and, for each prime $p \mid m$,
$\sum_{j<p} y^{jm/p}$ (named \lexlean!hg<p>!); at $y = \zeta^k$ they are the declared
\lexlean!hzM! and \lexlean!hg<p>!. An extended-gcd chain gives B\'ezout coefficients; a normalized
gcd other than $\Phi_m$ raises \lexcode!the relations of ζ do not generate Φ_m!.
\item Each nonzero quotient times coefficient over the leading coefficient, at $y = \zeta^k$,
is a term against its relation.
\end{enumerate}
\end{algorithm}

\subsubsection{Residues modulo a prime}\label{sec:lean-export-refute-residues}

\lexpy!_primes(m)! yields, for $t = \lfloor \mathtt{RES\_PRIME\_FLOOR}/m\rfloor + 1, t+1, \ldots$,
every prime $p = mt + 1$ with $w = g^{(p-1)/m} \bmod p$, $g$ the least primitive root
(\lexpy!sp.primitive_root!); $w$ has order exactly $m$. For the example $M = 8$ the first pair is
$(1\,048\,601, 386\,772)$. \lexpy!res(v)! (mirror of \lexlean!tmN!) is
$\sum c\,w^{e \bmod M} \bmod p$ over the terms of an exact value (atom parts ignored).
\lexpy!node_res(n, σ)! is \lexpy!res(node(n, σ))!, memoized per
$(p, w, \text{kind}, \text{dimension index}, \text{phase index}, \text{output values},
\text{input values})$; \lexpy!share_node_res(other)! adopts another refuter's memo when both have
equal \lexpy!M!, \lexpy!ptab!, \lexpy!tabs!, \lexpy!polys!, \lexpy!sparse!, \lexpy!ctab!, \lexpy!qtab!,
\lexpy!dimv_tab! and \lexpy!ks!.

\begin{algorithm}[\texttt{find\_sigma($\hat f$)}]\label{alg:lean-export-refute-sigma}
The candidates are the hint (\lexpy!hinted!) followed by the odometer order of the boundary labels'
ranges (\lexpy!_odometer!: \lexpy!itertools.product! order, last label fastest, nothing for an empty
range).
\begin{enumerate}
\item For the first \lexpy!MAX_PINNED! candidates $\sigma$: contract both networks at $\sigma$
(\lexpy!net_plan!) to residues $r_0, r_1$; return $\sigma$ and the plans at the first with
$\hat f_0 r_0 - \hat f_1 r_1 \not\equiv 0 \pmod p$.
\item Build both networks' residue plans with no pinned label (\lexpy!res_plan(net, {})!); a
\lexpy!NetUnsupported! returns \lexpy!None!.
\item For the first \lexpy!MAX_ASSIGNMENTS! candidates: read both residues from the final factors
(\lexpy!res_final!); at the first that differs as above, return $\sigma$ with both networks'
\lexpy!net_plan! at $\sigma$.
\item Return \lexpy!None!.
\end{enumerate}
\end{algorithm}

\paragraph{Hints.} \lexpy!hinted(bnd)! returns no candidate when there is no hint or the boundary
labels are not exactly $0, \ldots, |\mathit{bnd}|-1$. The hint lists a value per slot in label
order. When its length differs from the label count, \lexpy!bundled_hint! regroups it: each label
in ascending order takes $\prod$ (multiplicities of its chain) entries (one for an unbundled
label or outside a family) and becomes
$\mathrm{value}(d, c, ys) = \sum_{i < \mathrm{mult}(c_0)} \mathrm{value}(d, c_{1:},
ys_{[i\cdot\mathit{per},(i+1)\mathit{per})})\,b^i$ with $\mathit{per}$ the product of the
multiplicities of $c_{1:}$, $b = d^{\mathit{per}}$ and $d$ the root dimension (a value of an
empty chain is its single entry); a count mismatch gives no candidate. A value outside its range
gives no candidate; otherwise the single candidate places the value of the $s$-th label in
ascending order at that label's position in \lexcode!bnd!.

\paragraph{Contraction routes.} \lexpy!net_plan(net, σ)! returns \lexcode!("res", plan, r)! or
\lexcode!("sym", plan, r)!. When the widest node table (product of the dimensions of a node's
distinct labels outside $\sigma$) exceeds \lexpy!SYM_FIRST!, it tries the substitution contraction
\lexpy!sym_plan(self, net, σ)! of \module{archytaszx.semantics.lean_subst} first, falling through
on \lexpy!NetUnsupported!; then the residue plan \lexpy!res_plan(net, σ)! with
$r = \mathrm{res\_final}$; on \lexpy!NetUnsupported! the substitution contraction, whose failure
propagates.

\begin{algorithm}[\texttt{res\_plan(net, at)}]\label{alg:lean-export-refute-resplan}
Mirror of \lexlean!leafR! and \lexlean!elimR!.
\begin{enumerate}
\item Reset the entry count. Each node's leaf has scope its distinct labels (outputs then
inputs) outside \lexcode!at! and table $\tau \mapsto \mathrm{node\_res}(n, \tau \cup at)$, indexed
with the first label least significant (\lexpy!res_table!). Leaves are memoized per
$(p, w, M, \text{node}, \text{pinned values})$; a memoized leaf still counts its size.
\item Eliminate the distinct summed labels (\lexpy!res_eliminate!).
\item Reorder every scope (\lexpy!res_sorted!).
\end{enumerate}
\lexpy!res_count(size)! adds \lexcode!size! to the entry count and raises
\lexcode!a residue table is too large at the counterexample! when $\mathit{size} >
\mathtt{MAX\_TERMS}$ or the count exceeds \lexpy!MAX_RES_ENTRIES!.
\end{algorithm}

\begin{algorithm}[\texttt{res\_eliminate(leaves, labels)}]\label{alg:lean-export-refute-elim}
\begin{enumerate}
\item The \emph{trial table} of a label $l$ has scope the distinct labels other than $l$ of the
factors holding $l$, in factor order, and entries
$\sum_{j < \dim l}\prod_{f \ni l} f(\tau, l{=}j) \bmod p$ (\lexpy!res_sum!, by broadcasting
NumPy~\cite{harris2020numpy} arrays, \lexcode!int64! for $p < 2^{31}$ and Python integers otherwise,
memoized per scope, label and factor tables); it is absent when its scope's size exceeds
\lexpy!MAX_TABLE!. The cost of $l$ is $(1, 0)$ without a trial table, else $(0, \text{its nonzero
entry count})$.
\item While labels remain: take the label of least (cost, label); its new factor is the trial
table (now counted by \lexpy!res_count!) or, without one, \lexpy!res_sum! (counted). A new factor with a
nonempty scope and all entries equal becomes the constant factor $((), [v])$, its full form kept
in \lexpy!_untrim!. Remove the factors holding $l$, append the new factor, drop the trial tables of
the new scope's labels and recompute their costs, and record the step $(l, \text{factors})$, the
new factor last.
\item The final factors are those left.
\end{enumerate}
\end{algorithm}
\lexpy!res_sorted! orders each factor's scope by the key
$(-\mathrm{step}(x), x)$, $\mathrm{step}(x)$ the index of the step eliminating $x$ (the step
count for a label never eliminated), so labels never eliminated come first, then the latest
eliminated; tables are transposed to match and the full forms in \lexpy!_untrim! are reordered too.
\lexpy!res_final(final, σ)! (mirror of \lexlean!rfold!) multiplies the factors' entries at $\sigma$
from the last factor to the first modulo $p$; \lexpy!res_at! (mirror of \lexlean!RFac.at!) indexes a
table by mixed radix, first label least significant.

\subsubsection{Substitution plans (\texttt{lean\_subst})}\label{sec:lean-export-subst}

\module{archytaszx.semantics.lean_subst} plans and emits the symbolic eliminations of \file{RefuteSym} (\cref{ch:lean-core}) for one network at one boundary assignment $\sigma$; \module{archytaszx.semantics.lean_refute} and \module{archytaszx.semantics.refute_check} call it. Label expressions are tuples \lexpy{('v', x)}, \lexpy{('c', k)}, \lexpy{('add', a, b)}, \lexpy{('mul', k, a)}, \lexpy{('div', a, k)}, \lexpy{('mod', a, k)}, \lexpy{('neg', a, n)}; \lexpy{ev}, \lexpy{vars_}, \lexpy{subst}, \lexpy{ub}, \lexpy{lnorm}, and the smart constructors mirror \lexlean{LX.eval}, \lexlean{LX.vars}, \lexlean{LX.subst}, \lexlean{LX.ub}, \lexlean{LX.norm}, \lexlean{LX.sadd}, \lexlean{LX.smod}, \lexlean{LX.sdiv} (with division and remainder by $0$ giving $0$ and the value itself), and \lexpy{text} prints them as \lexlean{LX} terms. A \lexpy{State} holds bindings, node indices, constraints, and a constant (\lexpy{texpr} reads the first binding); a \lexpy{Plan} holds the states, the free labels after each step, the steps, the leaves, the groups and scopes of the final contraction, its plan, and its value.

\begin{center}\small
\begin{tabular}{@{}lll@{}}
\toprule
Constant & Value & Use \\
\midrule
\lexpy{MAX_SYM_STEPS} & $4000$ & most substitution steps \\
\lexpy{MAX_XSTEP} & $4096$ & largest X-spider dimension removed by an X step \\
\lexpy{RUN_STEPS} & $64$ & most steps of one \lexlean{sym_run} segment \\
\lexpy{LEAF_WALK_MIN} & $2^{18}$ & leaf checks past this size use \lexlean{leafS_ofW} \\
\bottomrule
\end{tabular}
\end{center}

\begin{algorithm}[\texttt{sym\_plan(refuter, net, sigma)}]\label{alg:lean-export-sym-plan}
With $d$ the refuter's label dimensions, $p$ its prime, $M$ its order, the free labels the network's summed labels in order, and \emph{pinned} the refuter's \lexpy{sparse_all} flag:
\begin{enumerate}
\item When pinned: remove each X spider of dimension $0<d_n\le$ \lexpy{MAX_XSTEP} with $d_n\mid M$ and all phase exponents $0$, adding the constraint $(d_n,\mathrm{xexpr})$ and multiplying the constant by $d_n$ (step \lexpy{("X", i, rest)}). When pinned, every later binding is normalized with the boundary labels read from $\sigma$.
\item Remove each \lexcode{Bun m X} node with one output, no inputs, and the same bounds, with modulus $d_n^{\mu(m)}$ (step \lexpy{('XB', i, rest)}).
\item Repeat while free labels remain and fewer than \lexpy{MAX_SYM_STEPS} steps exist, taking the first of: a delta step whose solution reads no free label; a delta step from a constraint; a geometric step; any delta step. A delta step's candidates are, per remaining node in order, every ordered pair of distinct legs of a Z spider (or \lexcode{Bun m Z}) whose first leg's expression is an unbound free label (possibly under \lexcode{neg}), and the three positions of a binder or splitter; then every constraint. A candidate on label $l$ is accepted when $l$ is free and unbound, the solution (as in \lexlean{solve}, with bare constraints allowed when pinned or after an \lexpy{XB} step) does not read $l$, reads only free or boundary labels, and has upper bound below $d(l)$. A geometric step takes the first free unbound label read by exactly two nodes, by no constraint, with $d(l)>0$, $d(l)\mid M$, both nodes of dimension $d(l)$, and both $F$ nodes with \lexpy{_fleg}. A delta step binds $l$ and substitutes everywhere; a geometric step removes the two nodes and adds the constraint $(d(l),X+Z)$, multiplying the constant by $d(l)$ mod $p$.
\item Finish (\lexpy{_finish}): each remaining node is one group with scope the labels its normalized leg expressions read; its residue table (raising \lexpy{NetUnsupported} past \lexpy{MAX_TABLE} entries) is the product of the refuter's node residues; each constraint's table is its divisibility indicator; the constant is the last leaf. Tables are reduced by dropping every coordinate they do not depend on (\lexpy{_reduce}); the refuter's eliminator orders and eliminates the remaining free labels; its final value is the plan's value.
\end{enumerate}
\end{algorithm}

\begin{algorithm}[\texttt{sym\_decls(refuter, tag, net, plan, env, unbnd)}]\label{alg:lean-export-sym-decls}
Emits, with prefix \lexpy{tag}: one \lexlean{Node} per network node (\lexpy{tag}\lexcode{k}$i$); the states and free lists at the segment ends (\lexcode{s}$k$, \lexcode{R}$k$); \lexcode{hU} (the boundary assignment in range, by \lexcode{decide +kernel}); for each segment a theorem \lexcode{t}$k$: a \lexlean{sym_run} over up to \lexpy{RUN_STEPS} non-geometric steps when pinned, or else one \lexlean{sym_dstep}, \lexlean{sym_xbstep}, or \lexlean{sym_gstep}; the residue declarations of the final contraction; the binding trie \lexcode{xbt} with \lexlean{toBT} checked; one \lexlean{leafS} theorem per group (by \lexlean{leafS_ofT}, or \lexlean{leafS_ofW} when the group's table equals its scope and $\prod d\cdot(\max d-1)/2$ exceeds \lexpy{LEAF_WALK_MIN}); one \lexlean{leafQ} theorem per constraint; and \lexcode{fin} by \lexlean{sym_final} (\lexlean{sym_finalW} for $M>2^{64}$). Every theorem carries \lexcode{set_option maxHeartbeats 0}. It returns the declarations and the chain \lexlean{sym_start} \lexcode{|>.trans} \lexcode{t}$k$ \ldots{} \lexcode{|>.trans (fin hw1)}. \lexpy{wide}, \lexpy{bound}, and \lexpy{prim_head} choose the \lexcode{W} variants, the $M\le 2^{64}$ proof, and \lexlean{prim_of} or \lexlean{prim_ofB} by the size of $M$.
\end{algorithm}

\subsubsection{The emitted refutation}\label{sec:lean-export-refute-source}

\lexpy!source_res! writes, before the theorem:
\begin{itemize}
\item \lexlean!Pt : ℕ → ℕ → ℕ × List ℤ! by cases: \lexlean!| ph, _ => (0, [])! per zero vector;
\lexlean!| ph, k => (((T * (M/den)) % M).toNat, [])! per phase table, $T$ entry $k$ of the table
(\lexlean!(([...]) : List ℤ).getD k 0!, or the chunked definitions of \lexpy!_Export.table_get! past
\lexpy!TABLE_CHUNK! entries; \lexpy!table_get!); the same with the digit polynomial's Lean text per
digit-polynomial vector; an \lexcode!if k = <k> then (<e>, []) else ...! chain ending in
\lexlean!(0, [])! per sparse vector (\lexpy!sparse_cases!); \lexlean!| ph, k => (([e0, e1, ...] : List ℕ).getD k 0, [])!
per \emph{dense} vector (\lexpy!dense_keys!: keys exactly $0..n-1$ with $n \ge \mathtt{PT\_DENSE\_MIN}$);
and a final lookup case over the remaining keys;
\item \lexlean!Ct! (each coefficient key's \lexpy!ctab! vector as \lexlean!Terms!), \lexlean!Qt! (each
coefficient denominator, default 1) and \lexlean!σ0! (the separating assignment, default 0);
\item per network (once when both sides share a network name) the contraction declarations of
its route, below.
\end{itemize}
The theorem \lexcode!main! (with \lexcode!set_option maxHeartbeats 0!) states the negated equality
(\lexpy!statement!) and proves it as follows.
\begin{lstlisting}[style=lexlean]
theorem main :
    ¬ ∀ (d : ℕ) (h_d : 0 < d), Equal (E d) (rt d) N0 1 N0 := by
  intro h
  have he := h 2 (by norm_num)
  have hp : Fact (Nat.Prime 1048601) := ⟨by norm_num⟩
  have hw := prim_of 1048601 386772 8 (by norm_num) (by norm_num) [(2, 3)] (by decide +kernel) (by norm_num) (by decide +kernel) (by decide +kernel)
  have hw1 : ((386772 : ℕ) : ZMod 1048601) ^ 8 = 1 := hw.pow_eq_one
  have hP : PhOK (E 2) 8 [] Pt [] := trivial
  have hzM : om 8 ^ 8 = 1 := om_pow_self 8 (by norm_num)
  have hg2 : (1 : ℂ) + om 8 ^ 4 = 0 := by simpa [Finset.sum_range_succ] using om_geom 8 2 (by norm_num) (by norm_num) (by norm_num)
  have hsq0 : (rt 2) ^ 2 = (2 : ℂ) := by have e := rt_sq' 2; linear_combination e
  have hf1 : (rt 2) = (((om 8) ^ 7) + (om 8)) := by have e := rt_sub2 8 1 (by norm_num) (by norm_num); linear_combination e
  have hC : CoefOK (E 2) 8 [] Ct Qt [] := trivial
  have hF : ∀ i ∈ ([0] : List ℕ), ∀ k, ex ((E 2).phase i k) = om 8 ^ (Pt i k).1 * mono (([] : List ℝ).map ex) (Pt i k).2 := by
    intro i hi k
    fin_cases hi
    · simp [E, Φ, Pt, ex_zero, mono_nil]
  have e0 := eval_W (E 2) 8 [] Pt Ct Qt [] [] [0] N0 σ0 (by norm_num) (by decide +kernel) hP hC hF
  have r0 := eval_res (E 2).shape 8 Pt Ct 1048601 386772 (by norm_num) hw1 (by norm_num) (by norm_num) N0 σ0 LN0 (by decide +kernel) HN00 STN0 FNN0 (by decide +kernel) CHN0 (by decide +kernel) (by decide +kernel) (by decide +kernel)
  refine refute_res σ0 (inRange_of _ _ _ (by decide +kernel)) 8 1048601 386772 (by norm_num) (by norm_num) hw _ _ _ _ _ _ ⟨e0, r0⟩ ⟨e0, r0⟩ he (1 : ℂ) [0, 1, 0, -1] [1] ?_ ?_ (by decide +kernel)
  · rw [KW, show sqLW (E 2).shape Qt N0 = [(2, 0, 1, 0)] from by decide +kernel]
    simp only [List.map_cons, List.map_nil, List.prod_cons, List.prod_nil, cvz, pow_zero, pow_one, mul_one, Nat.cast_one, inv_one] <;> (try push_cast) <;> linear_combination (1) * hf1 + (((om 8) ^ 3)) * hg2
  · rw [KW, show sqLW (E 2).shape Qt N0 = [(2, 0, 1, 0)] from by decide +kernel]
    simp only [...] <;> (try push_cast) <;> linear_combination 0
\end{lstlisting}
(the refutation of $\sqrt d\cdot Z$ against $Z$ at $d = 2$). In order:
\begin{enumerate}
\item \lexlean!he! instantiates the statement at the point (parameter values by \lexpy!value_text!:
naturals bare, others \lexlean!(a : ty)! or \lexlean!((a : ty) / b)!; bare multiplicities as
naturals; a \lexlean!(by norm_num)! per natural parameter).
\item \lexlean!hp! proves $p$ prime (\lexpy!prime_proof!): \lexlean!by norm_num!, or with
\lexpy!sparse_all! \lexlean!prime_lucas p g ...! with the least primitive root $g$ and the
factorization of $p - 1$ (\lexpy!_fac_list!: \lexlean![(q, e), ...]!), \lexlean!prime_lucasB <bits> p g ...!
for $p \ge 2^{64}$.
\item \lexlean!hw! proves that $w$ has order $M$ by \lexlean!prim_of p w M ... <factorization of M> ...!
(\lexlean!prim_ofB <bits>! for $M \ge 2^{64}$; \lexpy!prim_head!); \lexlean!hw1! is $w^M = 1$ in
\lexlean!ZMod p!.
\item \lexlean!hP : PhOK env M [] Pt keys! with one \lexlean!ex_atoms M e (z) [] [] (by norm_num) (by norm_num [E, Φ])!
per key ($e$ the power, $z = (tM - e)/M$; the unfold list is \lexlean!Fam.env, F, Φ! in a family),
then the relation lines and the atom facts of the certificate.
\item \lexlean!hC : CoefOK env M [] Ct Qt ckeys!: per coefficient key the \lexlean!Qt! value by
\lexlean!rfl! and \lexlean!norm_num!, and \lexpy!coef_proof!: \lexlean!Qt ph * coef = tv (om M) [] (Ct ph j)!
by \lexpy!combination! of $q_{ph}\,\mathrm{coef\_at}(ph, j) - \sum c\,\zeta^e$.
\item \lexlean!hF! covers the vectors \lexpy!zs!, \lexpy!tabs!, \lexpy!polys! and \lexpy!sparse! by
\lexlean!fin_cases!: \lexlean!simp [unfold, Pt, ex_zero, mono_nil]! per zero vector,
\lexlean!ex_tab M den (M/den) ...! per table or digit polynomial, and per sparse vector a case split
on membership of $k$ in its keys with one \lexlean!ex_atoms! per entry (\lexpy!sparse_proof!); with
no such vector it is \lexlean!by intro i hi; simp at hi!.
\item Per network $i$: \lexlean!e<i> := eval_W env M [] Pt Ct Qt keys ckeys fs N σ0 ...! and
\lexlean!r<i>! from its route. The residue route uses \lexlean!eval_res! with the leaf check
\lexlean!HN0!, or \lexlean!eval_resG! with the leaf table \lexlean!GN! when \lexpy!x_tables! gives some X
leaf a row check; the substitution route uses the declarations and term of \lexpy!sym_decls! of
\module{archytaszx.semantics.lean_subst}. For $M > 2^{64}$ the lemmas \lexlean!eval_res!,
\lexlean!eval_resG! and \lexlean!refute_res! take their \lexlean!W! forms and lose the
\lexlean!M ≤ 2 ^ 64! argument (\lexpy!wide!, \lexpy!bound!).
\item \lexlean!refute_res! takes $\sigma_0$, $M$, $p$, $w$, \lexlean!hw!, both pairs
\lexlean!⟨e<i>, r<i>⟩!, \lexlean!he!, the multiplier $L$ (\lexpy!l_text!: \lexlean!(<lcd> : ℂ)!, or
\lexlean!((<lcd> : ℂ) * <λ>)! with $\lambda$ as Lean text, $L = \ell\lambda$), $v_0$ and $v_1$, leaves the two goals $L s_jK_{W,j} = \sum_i v_{j,i}\zeta^i$
and decides the residue inequality. Each goal is closed by rewriting \lexlean!KW! with the
\lexlean!sqLW! list (\lexpy!_quads!), \lexlean!simp only!, \lexlean!push_cast! and the side's
\lexcode!linear_combination!.
\end{enumerate}
After the theorem the negated diagram-level restatement follows
(\lexpy!diagram_theorem(..., negated=True)!, or \lexpy!family_diagram_theorem! in a family). The
header is \lexcode!refutation at <point>, ζ = om <M>, mod <p>! (point
\lexcode!no parameters! when empty), and the theorem's doc comment
\lexcode!The diagrams are not equal: they differ at <point>.!.

\paragraph{Residue contraction declarations.} \lexpy!res_decls(tag, plan, dim)! writes, for tag
\lexlean!N0! and \lexlean!dim! the text \lexlean!(E 2).shape.dim!:
\begin{lstlisting}[style=lexlean]
def N0a0 : RFac := ⟨([] : List ℕ), RT.leaf 1⟩
def N0F0 : List RFac := [N0a0]
def N0S0 : List (ℕ × RFac) := []
theorem N0C0 : chainR 1048601 (E 2).shape.dim N0F0 N0S0 N0F0 = true :=
  chainR_nil 1048601 (E 2).shape.dim N0F0
def LN0 : List RFac := N0F0
def STN0 : List (ℕ × RFac) := N0S0
def FNN0 : List RFac := N0F0
theorem CHN0 : chainR 1048601 (E 2).shape.dim LN0 STN0 FNN0 = true := N0C0
noncomputable def QN0 : Node × RFac → Bool := fun q => leafR (E 2).shape 8 Pt Ct 1048601 386772 N0.bnd σ0 q.1 q.2
theorem HN01 : ((N0.nodes.drop 1).zip (LN0.drop 1)).all QN0 = true := by decide +kernel
theorem HN00 : ((N0.nodes.drop 0).zip (LN0.drop 0)).all QN0 = true :=
  zipAll_drop _ _ _ 0 (by decide +kernel) HN01
\end{lstlisting}
Leaves are \lexlean!<tag>a<i>!, new factors \lexlean!<tag>b<k>! (a constant new factor is defined by
its full form and enters the next list as \lexlean!<tag>b<k>.trim!), the lists between steps
\lexlean!<tag>F<k>!, the step lists \lexlean!<tag>S<k>!, the chain \lexlean!<tag>C<k>! through
\lexlean!chainR_step p dim l F<k> F<k+1> new S<k+1> F<n> E<k> rfl C<k+1>!, and
\lexlean!L<tag>!, \lexlean!ST<tag>!, \lexlean!FN<tag>!, \lexlean!CH<tag>!. A factor is
\lexlean!⟨scope, RT...⟩! nested one level per label, first label outermost, leaves
\lexlean!RT.leaf v!; since \lexpy!ELIM_DESCENT! is true every factor is packed (\lexpy!_rfac!): a level
whose children are equal leaves is written as that leaf. \lexpy!_ptree! gives the same packing as
Python data (\lexcode!(v,)! a leaf, \lexcode!("N", kids)! a level).

\begin{algorithm}[Elimination theorems \texttt{elim\_decls}]\label{alg:lean-export-refute-elimdecls}
The theorem \lexlean!<tag>E<k> : elimR p dim l F<k> new = true! is proved:
\begin{enumerate}
\item by \lexlean!elimT_ok p dim l F<k> new (by decide +kernel)! when \lexpy!desc_cost! is at most
\lexpy!DESC_MAX!;
\item otherwise, with $k$ the shortest prefix of the new scope whose suffix size times
$\dim l$ is at most \lexpy!ELIM_CHUNK!: for $k = 0$ by one \lexlean!decide +kernel!; else one theorem
\lexlean!<name>c<i> : elimC p dim l F new pre rest vs = true := by decide +kernel! per value list
\lexcode!vs! of the prefix labels (first prefix label fastest), and \lexlean!elimR_chunks ... rfl (by decide +kernel) [vs, ...] (by decide +kernel) (by simp only [...])!
over them.
\end{enumerate}
\lexpy!desc_cost(scope, ins, l)! mirrors the cursor steps of \lexlean!descR! over the packed tables of
the factors \lexcode!ins!: a descent over \lexcode!scope! adds, per call, the factor count times 1
(labels left) or $1 + 2\dim l$ (none left), stops at the bound, returns at an empty scope or
when a factor's subtree is the leaf 0, and otherwise, for each value $v$ of the first label,
adds $v$ per factor whose scope starts with that label, enters child $v$ of each such factor
that is not a leaf, and recurses; it returns \lexpy!None! past \lexpy!DESC_MAX!.
\end{algorithm}

\paragraph{Leaf checks.} \lexpy!leaf_decls(tag, env, n, tables)! defines \lexlean!Q<tag>!, the leaf
check \lexlean!leafR! (or \lexlean!leafG! over \lexlean!Node × RFac × LeafH! with \lexlean!tables!), and
\lexlean!H<tag><k>!, its \lexlean!all! over the leaves from $k$ on: \lexlean!H<tag><n>! by
\lexlean!decide +kernel!, each earlier one by \lexlean!zipAll_drop! from one \lexlean!decide +kernel!.
\lexpy!x_tables(net, leaves)! gives per leaf a \lexlean!LeafH! for a node of kind exactly \lexcode!X!
with $M$ divisible by $d$ and table size times $d$ at least \lexpy!LEAF_X_MIN!: with $g_s$ the
residue of the X tensor at the single output $s < d$ and no inputs, \lexlean!.xd <packed> <b>!
(checked by \lexlean!leafXD!), $b = \max(1, \text{bit length of } \max g)$ and
$\mathit{packed} = \sum_s g_s 2^{bs}$, when the node's labels are distinct and form its scope;
otherwise \lexlean!.x (RT.node [RT.leaf g0, ...])! (\lexlean!leafXR!); \lexpy!None! for any other leaf.
\lexlean!G<tag> : List LeafH! lists them, \lexlean!.plain! for \lexpy!None!.

\subsubsection{Recheck}\label{sec:lean-export-refute-recheck}

\lexpy!checked_residue(source, left, right, point, describe, found)! of
\module{archytaszx.semantics.refute_check} returns \lexcode!source! when
\lexpy!recheck_residue! finds nothing wrong, and otherwise raises \lexpy!RecheckFailed! (a subclass of
\lexpy!NetUnsupported!) with \lexcode!residue recheck failed: <why>!. The recheck fails with
\begin{itemize}
\item \lexcode!both boundaries have <n> slots! for an arity refutation whose sizes agree;
\item \lexcode!the search recorded no boundary assignment!;
\item \lexcode!<w> does not have order <M> mod the prime <p>! unless $p$ is prime, $w^M \equiv 1$ and
$w^{M/q} \not\equiv 1$ for every prime $q \mid M$;
\item \lexcode!the order of ζ is <M'>, not <M>! when a fresh \lexpy!_Refuter! with the same
\lexpy!sparse_all! chooses another order;
\item \lexcode!the fresh evaluation failed (<message>)! when the fresh refuter fails with
the given message while it recomputes $K_W$, $\lambda$, $f_j$, $\ell$, $\hat f_j$ and both
residues at $\sigma$ (by tables, then by substitution; the reverse order when the widest node
table exceeds \lexpy!SYM_FIRST!; the fresh refuter shares the found refuter's node residues);
\item \lexcode!the residues <res> mod <p> at σ differ from the search's <list>!;
\item \lexcode!the scaled entries agree mod <p> at σ = <σ>!.
\end{itemize}

\paragraph{Recheck details.} \lexpy!_order_ok(p, w, m)! holds when $p$ is prime, $w^m\equiv1$, and $w^{m/q}\not\equiv1 \pmod p$ for every prime $q\mid m$. \lexpy!_net_residue(r, net, sigma)! recomputes a network's residue at $\sigma$ by two routes, substitution (\lexpy!sym_plan(r, net, sigma).value!, \cref{sec:lean-export-subst}) and residue tables (one table per node over its legs not in $\sigma$, eliminated by \lexpy!res_eliminate! and read by \lexpy!res_final!), trying substitution first when the widest node table has more than \lexpy!SYM_FIRST! entries and tables first otherwise, the second route on a \lexpy!NetUnsupported!. \lexpy!recheck_residue! builds a fresh \lexpy!_Refuter! with the recorded \lexpy!sparse_all!, recomputes $M$ (which must equal the recorded one), the scaled constants $f_0,f_1$ at $w$ (with the root factor \lexpy!lmul! when \lexpy!sparse_all!), and both residues, which must equal the search's when recorded; it fails when $f_0r_0\equiv f_1r_1\pmod p$.

\lexpy!recheck_numeric(left, right, point, describe, found)! reads the evidence $(\sigma,g_0,g_1,x,y,\rho)$ and fails with \lexcode!the ball contains 0! when $\max(|x|,|y|)\le\rho$; it rebuilds a \lexpy!_NumRefuter!, requires the same $M$ and atoms, recomputes the entry groups of both networks at $\sigma$ (by plan or by brute force, \lexpy!use_plan!), which must equal $g_0,g_1$, and evaluates the difference $v$ with the factors $\prod(1/\sqrt d)^c$ of both sides. It fails when $|v|<$ \lexpy!NUMERIC_FLOOR! $=10^{-25}$ (\lexcode!the difference is 0 numerically!) or when $|\Re v-x|+|\Im v-y|>\rho+$ \lexpy!NUMERIC_SLACK! $=10^{-30}$ (\lexcode!the ball does not hold the difference!).

\lexpy!recheck_interface(left, right)! lists each boundary port's (dimension, multiplicity) on both sides, inputs and outputs separately (\lexpy!_side!, multiplicity the product of the \bbox{} multiplicities over the port); with free symbols it takes the first \lexpy!MAX_POINTS! points of \lexpy!sample_grid! (an exception gives none), else the empty point. It returns \lexpy!None! (the interface refutation holds) at the first point where both sides' flattened dimension lists (\lexpy!_flat!: each dimension repeated by its multiplicity, \lexpy!None! for a non-integer or nonpositive value) exist and differ, and otherwise \lexcode!interface recheck failed: the leg dimensions agree at every sample point!.

\subsection{The group evaluator}\label{sec:lean-export-refute-groups}

The numeric refuter evaluates networks to groups (mirror of \lexlean!evalG!): one coefficient
vector of length $M$ per atom monomial, zero vectors dropped. A value is a dictionary from
(atom exponents, power of $\zeta$) to integer coefficients, reduced modulo $M$ in the exponent,
trailing zero exponents trimmed (\lexpy!_add!, \lexpy!_mul!). Node tensors at port values: Z gives
$\zeta^{P(k)}$ when all ports equal $k$; X gives
$\sum_{k<d}\zeta^{P(k) + (M/d)(k\sum o + (d-1)k\sum i)}$; F gives $\zeta^{(M/d)\,o\,i}$; T is 1
when $o=0$ or $o=i$; Ti is $1$ at $o=i$, $-1$ at $o=0<i$; W is 1 when $i=0$ and all outputs 0,
or $0<i<d$ and the nonzero outputs are exactly $[i]$; I is $[o=i<d]$; B and S are
$[a<s][b<t][k = at+b]$; ZW and XW as in \cref{sec:lean-export-refute-weighted}; bundled kinds
read their bundled ports digit by digit (\lexpy!bundled!, \lexpy!digs_c!, \lexpy!copy_legs!), and
\lexcode!Bun! multiplies its copies. \lexpy!sq_list! gives per node $(d, c)$, $c$ the square-root
count: the port count of X and XW (and of bundled X), 1 for F, the multiplicity times the inner
count for a bundle, 0 otherwise and for B and S.

A network evaluates (\lexpy!groups!) by direct summation over its summed labels
(\lexpy!groups_brute!; work above \lexpy!MAX_TERMS! raises
\lexcode!the networks are too large to evaluate at the counterexample!) or, when the product of
summed dimensions exceeds \lexpy!BRUTE_LIMIT! (\lexpy!use_plan!), by contraction (\lexpy!plan!, cached
per network): leaves are node tables over their labels outside the pinned assignment, then
repeatedly the summed label of least cost (product of the other labels' dimensions among the
factors holding it; ties by label) is eliminated into a new factor. Tables above
\lexpy!MAX_TERMS! entries raise \lexcode!a factor table is too large at the counterexample!,
\lexcode!the factor tables hold more than 200000 terms at the counterexample! or
\lexcode!an intermediate factor is too large!. \lexpy!plan_decls! writes factors \lexlean!<tag>a<i>!,
\lexlean!<tag>b<k>!, lists \lexlean!<tag>F<k>!, steps \lexlean!<tag>S<k>!, per step
\lexlean!<tag>E<k> : elimNew ... = true := by decide +kernel!, the chain \lexlean!<tag>C<k>! through
\lexlean!chainNew_step!, and \lexlean!L<tag>!, \lexlean!ST<tag>!, \lexlean!FN<tag>!, \lexlean!CH<tag>!; a factor
is \lexlean!⟨scope, FT.node [...]⟩! nested one level per label, first label outermost, leaves
\lexlean!FT.leaf terms!. \lexpy!eval_rewrites! gives per network \lexlean!eval_cyc! (small networks) or
\lexlean!eval_ve_at! and \lexlean!eval_ve_groups! (contracted networks) with the \lexlean!show! terms of
\lexlean!sqL! and of the groups.

\subsection{Ball arithmetic \texttt{export\_numeric\_refutation}}\label{sec:lean-export-numeric}

\begin{signature}
\lexpy!export_numeric_refutation(left, right, point, name="main", describe=None) -> str!
\end{signature}
The numeric refuter (\lexpy!_NumRefuter!, a subclass of \lexpy!_Refuter! constructed with opaque
scalars allowed) refutes at a point where the values leave $\Q(\zeta)$: irrational phases, square
roots of any natural number, $\pi$. A \lexpy!ScalarUnsupported! becomes \lexpy!NetUnsupported!.
Different boundary sizes raise \lexcode!the boundaries differ in size!. The class attribute
\lexpy!weights_ok! is false. The written source is re-checked by
\lexpy!checked_numeric(source, left, right, point, describe, refuter)! of
\module{archytaszx.semantics.refute_check}, which reads the evidence
\lexpy!(sigma, g0, g1, x, y, r)! the refuter records; a failed recheck raises \lexpy!RecheckFailed!
with \lexcode!numeric recheck failed: <why>!, and a \lexpy!ScalarUnsupported! there raises
\lexpy!NetUnsupported! with \lexcode!numeric recheck failed: <message>!.

\begin{algorithm}[Numeric refutation]\label{alg:lean-export-numeric}
\begin{enumerate}
\item Opaque scalar symbols are replaced by their expressions (\lexpy!model_at!); a symbol without
a value at the point raises \lexcode!the counterexample gives no value for <s>!.
\item \emph{Atoms} (\lexpy!plan_atoms!): every phase entry at the point and every \lexpy!Ex! argument
of the scalars, in turns, splits into a rational $q$ plus rational multiples of non-rational
terms $b$; the bases are sorted by text, each with the lcm $L$ of its denominators; atom $j$ is
$x_j = e^{2\pi i\,b_j/L_j}$ (symbols \lexcode!x_lean_j!, inverses \lexcode!xi_lean_j!).
\item $M$ is the lcm of 1, the node dimensions and the denominators of the rational parts. The
entry table maps (phase, value) to ($qM \bmod M$, atom exponents).
\item Different boundary label sets raise \lexcode!the diagrams have different boundaries!.
\item Search assignments (odometer order, at most \lexpy!MAX_ASSIGNMENTS!): evaluate the difference
$s_0k_0g_0 - s_1k_1g_1$ ($g$ the groups, \cref{sec:lean-export-refute-groups}; $k$ the
square-root factor of \lexpy!sq_list!) numerically (\lexpy!num!, mpmath with 60 digits) and keep the
largest; stop at the first with a component above $10^{-6}$. If the best is below $10^{-25}$,
raise \lexcode!the networks agree numerically at every assignment checked!.
\item Build a ball for the difference (\lexpy!tree!); if neither $|x|$ nor $|y|$ exceeds the radius,
raise \lexcode!the balls are too wide to separate the entry from zero!.
\end{enumerate}
\end{algorithm}

\paragraph{Balls.} \lexlean!NearR t m d! states $|t - m| \le d$ for a real $t$; \lexlean!Near t x y r!
states that the complex $t$ lies within $r$ of $x+iy$. Midpoints are rounded to 40 decimal
digits (\lexpy!_round!: nearest multiple of $10^{-40}$) and radii rounded up (\lexpy!_up!), each
rounding absorbed by \lexlean!nearR_round! or \lexlean!near_round!. Real balls (\lexpy!rball!):
\begin{itemize}
\item a rational: radius 0 (\lexlean!nearR_const!);
\item $\pi$: midpoint $3.14159265358979323846$, radius $10^{-20}$ (\lexlean!nearR_pi!);
\item $\sqrt p$: midpoint $\lfloor\sqrt{p}\,10^{45}\rfloor/10^{45}$, radius $2\cdot10^{-45}$
(\lexlean!nearR_sqrt!, four \lexlean!norm_num! side conditions);
\item $1/\sqrt p$ and other inverses: from the ball $(m,d)$ of the base with $m-d>0$, the interval
$[1/(m+d), 1/(m-d)]$ re-centered (\lexlean!nearR_inv!);
\item sums and products: $(m_1+m_2, d_1+d_2)$ and $(m_1m_2, |m_1|d_2+|m_2|d_1+d_1d_2)$, rounded
(\lexlean!nearR_add!, \lexlean!nearR_mul!).
\end{itemize}
Complex balls (\lexpy!tree!, \lexpy!atom_ball!): $i$ exactly (\lexlean!near_I!); rationals exactly
(\lexlean!near_rat!); real constants through \lexlean!near_ofReal!; sums and products componentwise,
the product radius $(|x|+|y|)r_2 + (|x_2|+|y_2|)r + rr_2$ (\lexlean!near_add!, \lexlean!near_mul!).

\begin{algorithm}[\texttt{ex\_ball(ctx, $\theta$)}: a ball for $e^{2\pi i\theta}$]\label{alg:lean-export-exball}
\begin{enumerate}
\item $z = \mathrm{round}(\theta)$; $\rho = \theta - z$; the argument is $a = 2\pi(\rho/8)$.
\item A real ball $(A, \varepsilon)$ for $a$ (\lexpy!angle_ball!: balls for $\rho$ and $\pi$, then
$2\pi$, $\rho/8$ and their product), or $(0,0)$ when $\rho=0$.
\item Taylor polynomials with $n=10$ terms and their tail bound,
\[
C = \sum_{k<n}\frac{(-1)^kA^{2k}}{(2k)!},\qquad
S = \sum_{k<n}\frac{(-1)^kA^{2k+1}}{(2k+1)!},\qquad
t = \frac{|A|^{2n}(2n+1)}{(2n)!\,2n},
\]
give the radius $2(\varepsilon + t)$ (\lexlean!near_exp!). The centers are checked by
\lexlean!norm_num [cosT, sinT, Finset.sum_range_succ, Nat.factorial]!.
\item Square three times (\lexlean!near_mul!) to reach $w^{8}$, and rewrite
$e^{2\pi i\theta} = w^{8}$ by \lexlean!ex_congr_int!, \lexlean!ex_pow_eight!, \lexlean!ex_eq_exp!.
\end{enumerate}
\end{algorithm}
Atoms $\omega_k$, $\zeta^{k}$, $x_j$, $x_j^{-1}$ (\lexlean!ex_inv!) and $e^{a}$ (\lexlean!exp_eq_ex!) all go
through \lexpy!ex_ball!.

The theorem body follows the group evaluator's form:
\lexlean!refine refute σ0 (inRange_of _ _ _ (by decide +kernel)) ?_ he!, \lexlean!intro h0!,
\lexlean!hP : PhOK env M GS Pt keys! with the atom list \lexlean!GS! and \lexlean!Pt! carrying atom
exponents (\lexlean!ex_atoms M p (z) gs ks!), the \lexlean!rw! of \lexpy!eval_rewrites! at \lexlean!h0!,
\lexlean!unfold sqK at h0!, the \lexlean!show! rewrites, a \lexlean!simp only! set extended by
\lexlean!mono!, \lexlean!zpow_neg!, \lexlean!zpow_one!, \lexlean!zpow_ofNat!, \lexlean!zpow_zero!, then
\lexlean!have hD : <difference> = 0 := by linear_combination h0!, the ball facts, and
\lexlean!exact ne_zero_of_near <ball> (Or.inl (by norm_num)) hD! (\lexlean!Or.inr! when the imaginary
part separates), then the \lexlean!¬ DEqual! restatement \lexlean!main_diagrams! over the engine's
diagrams as for the exact refutation (\lexpy!_Refuter.diagram_negation!). The header is
\lexcode!refutation by ball arithmetic at <point>!.
\lexpy!ball_tree(M, out, expr, prefix)! exposes the ball construction for a polynomial in
$\zeta = \omega_M$ (used for Gauss-sum signs), returning the facts, the text, the fact name and
the ball.

% ---------------------------------------------------------------------------
\section{Linked Diagrams: Peel, Slide, Regroup and Split}\label{sec:lean-export-linked}

Three modules export single proof-kernel steps over \bboxes (\cref{ch:kernel}) as Lean
theorems about \emph{linked diagrams} (\file{ArchytasZX/Linked.lean} and its siblings
\file{LinkedPeel.lean}, \file{LinkedSlide.lean}, \file{LinkedForm.lean},
\file{LinkedGhost.lean}, \file{LinkedSplice.lean}, \file{LinkedFirst.lean},
\file{LinkedDropX.lean}, \file{LinkedRead.lean}, \file{LinkedPair.lean},
\file{LinkedRegroup.lean}, \file{LinkedSplit.lean}; \cref{ch:lean-boxes}):
\module{archytaszx.semantics.lean_linked} (peel and identity slide),
\module{archytaszx.semantics.lean_regroup} (box-group regroup) and
\module{archytaszx.semantics.lean_split} (power split). Each export returns the Lean source or
\lexpy!None! when the step lies outside its form; it catches \lexpy!LinkedUnsupported! itself.
Every Lean side condition is first checked on a Python mirror of its Lean definition, field
by field (for the pair-box conditions \lexlean!PairFirst! and \lexlean!PairLast! of
\cref{sec:lean-export-pair}, \lexpy!pair_ok! checks \lexlean!multL!, \lexlean!multR! and
\lexlean!form!: each pair leg a leg of \lexlean!Lb! on a node of $B$ with leg chain
$[\mathit{mb}]$), and then in Lean by \lexlean!decide +kernel!. Every file
imports \lexlean!ArchytasZX!, opens namespace \lexlean!ArchytasZX.Export! and
\lexlean!open ArchytasZX ArchytasZX.LDiagram!.

The Lean digit sums the theorems compare are: \lexlean!evalL! (plain linked diagram),
\lexlean!evalG gs! (with ghosts \lexlean!gs!), \lexlean!evalX! and \lexlean!evalGX gs! (with box index
readings \lexlean!cph!, \lexlean!cdim!), \lexlean!evalN! and \lexlean!evalGXR gs gr! (with node dimension
readings \lexlean!cnd! and ghost read bounds \lexlean!gr!), and \lexlean!evalP ps! (with pair boxes
\lexlean!ps!).

\subsection{The linked-diagram mirror}

\begin{signature}
\lexpy!Port = (node, out, pos)!: a port as node id, output side flag and position.\\
\lexpy!LNode(id, kind, dim, ph, outs, ins)!: kind one of \lexcode!Z X F T Ti W B S!, dimension, phase
and per-port dimension indices.\\
\lexpy!LBox(id, mult, parent, nodes)!: a node-scoped \bbox with multiplicity index.\\
\lexpy!LLink(box, l1, l2, e1, e2)!: a chain link of a \bbox, its two ports and their partners.\\
\lexpy!Cph = (node, box, coef, mult)!: a copy-indexed phase; \lexpy!CDim = (port, box, coef, mult)!:
a copy-indexed port dimension.\\
\lexpy!LDiag(nodes, wires, outs, ins, boxes, links, cph=[], cdim=[], cnd=[])!: the mirror of the Lean
\lexlean!LDiagram! without aligned wires; \lexcode!cnd! holds node dimension readings (\lexpy!Cph! shape).\\
\lexpy!class LinkedUnsupported(Exception)!.
\end{signature}
A reading $(x, y, k, m)$ states that the entry of node (or port) $x$ is the index of \bbox $y$
read at the copy digit plus $k$ times the value of multiplicity index $m$; the engine writes
$k = m = 0$ for a plain reading. \lexpy!LDiag.walk(pred, fuel, start)! repeatedly takes the first box
meeting \lexcode!pred! whose parent is the previous box (starting at \lexcode!start!), at most
\lexcode!fuel! times. \lexpy!node_boxes(x)! is \lexpy!walk(x in nodes, |boxes|, None)!; \lexpy!chain(x)! is
their multiplicity indices, \lexpy!depth(y, x)! the position of box $y$ in that walk (its length
when absent), \lexpy!mult_of(y)! the multiplicity index of the first box $y$ (0 when absent),
\lexpy!legs()! every port of every node (outputs then inputs) and \lexpy!bnd()! outputs then inputs.
\lexpy!flip(D, bid)! mirrors \lexlean!flipLinks!: every link of box \lexcode!bid! becomes
\lexpy!LLink(box, l2, l1, e2, e1)!, the readings kept except \lexcode!cnd!. \lexpy!strip(D)! mirrors
\lexlean!strip!: $D$ without \lexcode!cph!, \lexcode!cdim! and \lexcode!cnd!. \lexpy!_find(entries, key)! mirrors
\lexlean!(entries.find? (·.1 = key)).map (·.2)!: the triple after the first entry of that key, else
\lexpy!None!. \lexpy!_Tables! numbers dimension texts, phase keys and expanded multiplicity texts in
order of first use (\lexpy!index(table, key)! inserts the next number), shared by both sides.

\subsection{Phase keys and box index readings}\label{sec:lean-export-readings}

\begin{itemize}
\item \lexpy!_phase_key(phase)!: \lexpy!repr(None)! for a \lexpy!PhaseVector! with \lexpy!is_zero!, else
\lexpy!repr(phase)!. The phaseless index \lexlean!zero! is therefore shared by phaseless nodes and zero
phase vectors.
\item \lexpy!_reads(text, sym)!: \lexcode!sym! is one of the names (regular expression
\lexcode![A-Za-z_]\w*!) of \lexcode!text!. \lexpy!_sym(text)!: \lexpy!sympify! with every name a plain
symbol. \lexpy!_index_names(d)!: the set of box indices of $d$. \lexpy!_equal(a, b)!:
\lexpy!simplify(a - b) == 0!.
\item \lexpy!_seed(ds, tables)!: for each diagram in order, each node in id order, number every port
dimension text reading no box index of any diagram in \lexcode!ds!, and the phase key of every node
whose phase is \lexpy!None! or has no free symbol among those indices; then number the empty
dimension text. Every plain index thus precedes every index of a reading text.
\item \lexpy!_value_key(tables, value, names)!: the first nonempty numbered dimension key reading
none of \lexcode!names! whose value equals \lexcode!value!, else \lexpy!str(value)!.
\item \lexpy!_end_dim(d, bp, text, tables)!: the dimension text of boundary link port \lexcode!bp! at
the copy the boundary meets. With no index read, \lexcode!text!. Otherwise for the first link of $d$
holding \lexcode!bp!: the read indices must be exactly that box's index (else
\lexcode!a link port reading another box index!); the index takes $\mu_b - 1$ when \lexcode!bp! is the
link's first port, $0$ when the second, and the result is \lexpy!_value_key! of the substituted
value. No link holding \lexcode!bp! raises \lexcode!an identity end reading a box index off a link!.
\end{itemize}

\lexpy!_readings(d)! returns \lexcode!(cph, cdim)!, every entry with $k = m = 0$. Per node (id order)
the \emph{first dimension} is that of its first port (the empty text when legless), for a binder
B with two inputs that of input 0; the \emph{second dimension} is input 1 of such a B or output 1
of a splitter S with two outputs. For each box index \lexcode!sym! (sorted) that the node's phase,
its second dimension or some port dimension reads:
\begin{enumerate}
\item the binders are the boxes holding the node, sorted by depth, whose index is \lexcode!sym!;
none, or a first dimension reading \lexcode!sym!, raises
\lexcode!a box index read outside the linked form!;
\item the innermost binder $y$ is the box of the reading;
\item a phase or second-dimension reading adds $(x, y, 0, 0)$ to \lexcode!cph!; a node already in
\lexcode!cph!, or one reading \lexcode!sym! in both phase and second dimension, raises
\lexcode!a node reading two box indices!;
\item each port reading \lexcode!sym! adds $(p, y, 0, 0)$ to \lexcode!cdim!; a port already in
\lexcode!cdim! raises \lexcode!a leg reading two box indices!.
\end{enumerate}

\begin{signature}
\lexpy!linked(d: Diagram, tables: _Tables, idents: Mapping[Port, int] | None = None, ghost: int | None = None, held: Mapping[Port, LBox] | None = None, readings: bool = False) -> LDiag!
\end{signature}
\begin{algorithm}[\texttt{linked}]\label{alg:lean-export-linked}
\begin{enumerate}
\item With \lexcode!readings!, \lexpy!cph, cdim = _readings(d)!; otherwise both empty.
\item Nodes in id order; a generator outside the kinds raises \lexcode!generator <name>!. The node
dimension is its first port's dimension index (the index of the empty text for a legless node);
the phase index is that of \lexpy!_phase_key(node.phase)!. A binder B with two inputs takes
\lexcode!(dim, ph)! from its two input dimension indices; a splitter S with two outputs from its two
output dimension indices.
\item Wires in \lexpy!sort_key! order; the boundary as ports.
\item For each boundary port $p$ in \lexcode!idents! (sorted), add a phaseless Z ($1\to1$, phase
index of \lexpy!repr(None)!) of id \lexcode!idents[p]!, wired to $p$, its other port taking $p$'s
boundary place; its dimension is that of $p$ (with \lexcode!readings!, the text of
\lexpy!_end_dim!). A port off the boundary raises \lexcode!an identity end off the boundary!.
\item With \lexcode!ghost!, every boundary port in slot $j$ that is a link port becomes
\lexcode!(ghost, True, j)!, or \lexcode!(x, True, j + 1)! when it is in \lexcode!held! with \lexcode!held[p]!
holding the node $x$; the new port is wired to the original port.
\item Every box of \lexcode!held! must hold one node, else
\lexcode!a held identity box of other than one node!; this is checked before the ghost step.
\item Each \bbox in id order must be a plain node box (no port scope, nonempty node scope, not
cyclic, no index unless \lexcode!readings!, no copy scalar, pair, tree, feed, sheets or alignment),
else \lexcode!box <id> is no plain node box!. Each link port must be on exactly one wire not used by
another link port (\lexcode!a link port without one wire of its own!); the link records the
partners, and those wires leave the wire list. The boxes of \lexcode!held! follow, by sorted port.
\item Every kept wire must join two nodes whose box-id walks are prefixes of one another, else
\lexcode!a wire between boxes neither of which holds the other!.
\end{enumerate}
\end{algorithm}
\lexpy!literal(D)! writes \lexlean!{ base := { nodes := [⟨id, .Z, dim, ph, [outs], [ins]⟩, ...], wires := [...], outs := [...], ins := [...], boxes := [⟨id, mult, none, [nodes], []⟩, ...] }, links := [⟨box, l1, l2, e1, e2⟩, ...] }!,
ports as \lexlean!⟨node, true, pos⟩!, a parent as \lexlean!some p!. With readings it appends
\lexlean!cph := [(x, y, k, m), ...]! and \lexlean!cdim := [(⟨..⟩, y, k, m), ...]!, and with
\lexcode!cnd! also \lexlean!cnd := [...]!.

\subsection{Identity drops and ghosts}

Each export writes each side with a phaseless Z($1\to1$) behind every boundary link port
(\lexlean!Lb!, \lexlean!Lk!) and proves that side's digit sum equal to the ghost form's (\lexpy!_chain!,
\lexpy!_chain_n!). Two step kinds remove the identities; each mirror runs on the diagram with the
identity $x$ moved last (\lexpy!_move_last!).
\begin{itemize}
\item \emph{drop} (\lexlean!DropOK!, 18 fields, \lexpy!drop_ok!): wire chains agree, node ids are
unique, $x$ is last, a Z with \lexlean!outs = ins = [dim]!, $\{u, v\}$ its two ports, $w$ is $(u, P)$
or $(P, u)$ with $P$ a port of another node, $u$ on no other wire, not on the boundary and on no
link, $v$ no link port, $w$ a wire, no \lexlean!cph! entry, and the port dimension of $P$ equals
$x$'s (\lexlean!pdim!). The Lean field \lexlean!ports! has no conjunct in \lexpy!drop_ok!: an
\lexpy!LBox! has no port scope. \lexpy!drop_id! removes $x$ and $w$ and
renames $v$ to $P$ in wires, boundary and links; Lean step \lexlean!evalL_dropAt!.
\item \emph{indexed drop} (\lexlean!DropOKX!, \lexlean!DropOK! on \lexlean!L.strip! plus 4 fields,
\lexpy!drop_x_ok!): \lexpy!drop_ok! on the stripped diagram, the dimension of $P$ equals $x$'s, $x$ in
no \lexcode!cph! entry and no \lexcode!cdim! entry, $P$ in no \lexcode!cdim! entry; Lean step
\lexlean!evalX_dropAt!, \lexlean!dropAtX!.
\item \emph{read drop} (\lexlean!DropN!, \lexlean!DropOK! plus 6 fields, \lexpy!drop_n_ok!): as
\lexpy!drop_x_ok! for the stripped \lexpy!drop_ok! and the dimension of $P$; $x$ in no \lexcode!cph!
entry; with $o$ the \lexcode!cnd! reading of $x$, the \lexcode!cdim! readings of both ports of $x$ and
of $P$ equal $o$; and when $o$ exists, $P$'s node and $x$ lie at one depth in $o$'s box. Lean step
\lexlean!evalN_dropAt!, \lexlean!dropAtN!.
\item \emph{ghost} (\lexlean!GhostOK!, 23 fields, \lexpy!ghost_ok!): wire chains agree, ids unique, $x$
last, a Z with \lexlean!outs = ins = [dim]!, $\{a, c\}$ its two ports, $x$ on no wire, the boundary
without repeats, slot $i$ holds $c$, $a$ and $g$ off the boundary, $g$ on no wire and distinct
from $a$ and $c$, $g$'s node has $x$'s chain, $c$ and $g$ on no link, $a$ no link port, every link
with an end at $a$ has the remaining chain (its chain without the link's box) equal to $x$'s
chain, and no ghost of \lexcode!gs! is at $a$, $c$ or $g$. \lexpy!ghost_drop! mirrors
\lexlean!ghostDropX!: removes $x$, renames $a$ and $c$ to $g$ in the boundary and the links, and
drops $x$'s \lexcode!cdim! and \lexcode!cnd! entries; Lean step \lexlean!evalG_ghostAt!.
\item \emph{indexed ghost} (\lexlean!GhostOKX!, \lexlean!GhostOK! plus 3 fields, \lexpy!ghost_x_ok!):
\lexpy!ghost_ok! and $x$ in no \lexcode!cph! or \lexcode!cdim! entry; Lean step \lexlean!evalGX_ghostAt!,
\lexlean!ghostAtX!.
\item \emph{read ghost} (\lexlean!GhostN!, \lexlean!GhostOK! plus 6 fields, \lexpy!ghost_n_ok!):
\lexpy!ghost_ok!, $x$ in no \lexcode!cph! entry, the \lexcode!cdim! readings of both ports of $x$ equal
its \lexcode!cnd! reading $o$, $g$'s node and $x$ at one depth in $o$'s box when $o$ exists, and no
read bound of \lexcode!gr! at $g$. \lexpy!_gr_add(D, x, g, gr)! mirrors \lexlean!grAdd!: \lexcode!gr! with
$(g, o)$ prepended when $x$ has a \lexcode!cnd! reading $o$. Lean step \lexlean!evalGXR_ghostAt!,
\lexlean!ghostAtN!.
\end{itemize}
The ghost of slot $i$ is $(x, \mathrm{true}, i+1)$ when $x$ lies in a box (a held identity), else
$(G, \mathrm{true}, i)$.

\begin{algorithm}[\texttt{\_reduce(D, ids, drops, G, xm=False, rm=False)}]\label{alg:lean-export-reduce}
\begin{enumerate}
\item With \lexcode!xm! and nonempty \lexcode!drops!, first run \lexpy!_reduce(D, ids, set(), G, True, rm)!;
return it when its ghosts are not \lexpy!None!.
\item For each node of \lexcode!drops! in node order: \lexpy!_drops! lists, over its two ports $u$ (output
first) and the wires at $u$ that are no self-loop, every $(w, u, v, P)$ meeting \lexpy!drop_ok!
(with \lexcode!xm! \lexpy!drop_x_ok!, with \lexcode!rm! \lexpy!drop_n_ok!); the first is applied, the node's
\lexcode!cdim! and \lexcode!cnd! entries removed.
\item For each remaining node of \lexcode!ids! in node order, moved last: for its port $c$ on the
boundary (output first), $a$ the other port, $i$ the slot of $c$, $g$ as above; when the ghost
condition holds (with \lexcode!rm! and remaining \lexcode!cnd! \lexpy!ghost_n_ok!, else with \lexcode!xm!
\lexpy!ghost_x_ok!, else \lexpy!ghost_ok!), record the step, update the read bounds by \lexpy!_gr_add!,
apply \lexpy!ghost_drop! and prepend the ghost $(g, i, \mathit{dim})$.
\item Return the diagram, the ghosts newest first (\lexpy!None! when a node of \lexcode!ids! remains),
the steps and the read bounds.
\end{enumerate}
\end{algorithm}

\lexpy!_chain(tag, start, end, gname, steps, gs, define=True, x=False)! writes per step
\lexlean!abbrev x<tag><j> : DNode!, the theorem \lexlean!ok<tag><j>! (its fields by
\lexlean!decide +kernel!; with \lexcode!x! the nested proof \lexlean!⟨⟨...⟩, ...⟩! of \lexlean!DropOKX! or
\lexlean!GhostOKX!), the ghosts \lexlean!abbrev gsL : Ghosts := [(⟨g⟩, i, dim), ...]! (with
\lexcode!define!), \lexlean!eq_<tag> : <last> = end! when there are steps, and
\begin{lstlisting}[style=lexlean]
theorem chain_b : ∀ (E : Env), (∀ k, E.phase zero k = 0) →
    ∀ (ks : ℕ → ℕ) (σ : ℕ → List ℕ → ℕ),
      Lb.evalL E ks σ = L.evalG gsL E ks σ :=
  fun E hE ks σ => ((evalL_dropAt okb1 E hE ks σ).trans ((evalG_nil ... E ks σ).symm.trans
    (evalG_ghostAt okb2 E hE ks σ)))).trans (by rw [eq_b])
\end{lstlisting}
With \lexcode!x! the names are \lexlean!evalX!, \lexlean!evalX_dropAt!, \lexlean!dropAtX!,
\lexlean!evalGX_nil!, \lexlean!evalGX_ghostAt!, \lexlean!ghostAtX! and the right side
\lexlean!L.evalGX gsL!. \lexlean!evalG_nil! (\lexlean!evalGX_nil!) switches to the ghost sum before the
first ghost step, or at the end when there is none.

\lexpy!_chain_n(tag, start, end, gname, rname, steps, gs, gr, D0)! writes the same chain over
\lexlean!evalN!, replaying the steps on $D_0$: drops through \lexlean!DropN!
(\lexlean!⟨⟨18 fields⟩, 6 fields⟩!) and \lexlean!evalN_dropAt!; ghost steps while \lexcode!cnd! readings
remain through \lexlean!GhostN ... gr ...! (the read bounds before the step) and
\lexlean!evalGXR_ghostAt!. At the first ghost step with no \lexcode!cnd! left (or at the end, when every
step is a drop and no \lexcode!cnd! is left) it inserts
\lexlean!(evalN_plain <at> (by decide +kernel) E ks σ)! and \lexlean!(evalGX_nil <at> E ks σ).symm!
and continues with \lexlean!GhostOKX! and \lexlean!evalGX_ghostAt!. When the chain ends in the read
mode it also writes \lexlean!abbrev grL : GReads := [...]!, the right side is
\lexlean!L.evalGXR gsL grL!, and the tail is \lexlean!.trans (by rw [eq_b]; congr 1 <;> decide +kernel)!;
otherwise the right side is \lexlean!L.evalGX gsL!. With no terms the body is \lexlean!rfl!. It returns
the text and the right-side evaluation.

\paragraph{Slot data.} \lexpy!groups(d)! lists, per \bbox with groups in box order, its id and each
group's sorted nodes. \lexpy!slot_data(D, gs)! mirrors \lexlean!slotData!: the output and input counts
and per box its id, multiplicity index, parent and the nonempty slot blocks (the boundary
positions whose node is in each group, then in the box's ungrouped nodes). \lexpy!nested(sd)! holds
when each block of a box lies in one block of its parent. \lexpy!anchors(sd, i)! mirrors
\lexlean!(anch sd i, chainS sd i)!: walking the boxes holding slot $i$ from the outermost, per box the
least slot of its block on $i$'s side (outputs or inputs), then $i$; and the multiplicity indices.
The theorems read digits at the flat positions \lexlean!flatAt (L.slotData G) ks y! (the proof
kernel's boundary splice).

\subsection{Peel ends}\label{sec:lean-export-peel-ends}

\begin{signature}
\lexpy!PeelEnds(box, B, fresh, cp, ends_l, ends_r, held)!: the peeled box, its sorted nodes, the copy
ids, the copy map, the identity ids of both sides by port, and the boundary link ports of
\lexcode!before! whose identity \lexcode!after! holds in a box of its own.\\
\lexpy!peel_ends(before, after, tables, first=False) -> PeelEnds | None!
\end{signature}
\begin{algorithm}[\texttt{peel\_ends}]\label{alg:lean-export-peel-ends}
\begin{enumerate}
\item Exactly one outermost \bbox of \lexcode!before! must exist in \lexcode!after! with
$\mu_{\mathit{before}} - \mu_{\mathit{after}} = 1$ (after expansion); $B$ its sorted nodes.
\item The new nodes are the ids of \lexcode!after! not in \lexcode!before!. \lexpy!_engine_ends(after)!
maps the wired port of every phaseless Z($1\to1$) (\lexpy!_phase_key! of \lexpy!repr(None)!) with one
port on the boundary and one wire to that node id. With $|\mathit{new}| = |B|$ no engine
identity is used; otherwise the engine identities among the new nodes must number
$|\mathit{new}| - |B|$. \lexcode!fresh! are the other new nodes; \lexcode!cp! pairs $B$ with
\lexcode!fresh! in order.
\item \lexpy!_boundary_ends(before, after, box, cp, first)!: one fresh id (from one more than the
largest node id of both sides) per link port on the boundary of \lexcode!before!, the same id on
\lexcode!after! at the same port or, for the first port of a link of the peeled box (with
\lexcode!first!, the second port), at its copy \lexpy!_cpp(cp, p)!.
\item \lexpy!_held(after, ends_r, tables, engine)!: each port of \lexcode!ends_r! off the boundary of
\lexcode!after! that an engine identity is wired to; the identity must lie in exactly one box of
\lexcode!after! and that box must be outermost, else \lexcode!an engine identity end in no single outermost box!
(caught: \lexpy!None!); the result maps the port to the identity id and to
\lexpy!LBox(box, mult, None, (id,))!.
\item Each such port's \lexcode!before! end takes the engine identity id, its \lexcode!held! entry is the
box, and the \lexcode!after! end is removed. The new nodes must be exactly \lexcode!fresh! and the held
identities.
\end{enumerate}
\end{algorithm}

\subsection{Peel conditions}\label{sec:lean-export-peel-ok}

\paragraph{\texttt{PeelOK} (21 fields, \texttt{peel\_ok}).} Boxes nonempty; wires within one chain
on both sides; $R$'s nodes are $L$'s followed by the copies of $B$'s nodes (\lexpy!cpn!); unique ids,
ports and boundary of $R$; chains of $B$ nodes $[m_b, r]$ in $L$ and $[c, r]$ in $R$ with $r$ the
copy's chain in $R$, other chains kept; wires within or off $B$, kept, and for $B$ copied; every
wire of $R$ kept, copied, or $(\mathit{cp}(l_1), e_1)$ for a link of the peeled box; the link
well-formedness \lexpy!_links_ok! (\lexlean!lwf!: four ports exist, $l_2$ has $l_1$'s chain, the box at
its depth with its multiplicity, the ends have the remaining chain); links of the peeled box with
$l_1, l_2 \in B$, $e_1, e_2 \notin B$, depth 0, multiplicity $m_b$, moved to
$\langle l_1, l_2, \mathit{cp}(l_2), e_2\rangle$ in $R$ at depth 0 and multiplicity $c$ with the
wire $(\mathit{cp}(l_1), e_1)$; inner links of $B$ kept (\lexpy!_same!) and copied one level
shallower; links off $B$ kept; every link of $R$ one of these; boundary kept and extended by
copies.

\paragraph{\texttt{RevOK} (16 fields, \texttt{rev\_ok(L, R, bid, mb, B)}).} $R$ equals $L$ but for
\lexlean!flipLinks bid!, no \lexcode!cph! on either side; unique ids, ports and boundary; the boundary
among the ports; wire chains agree; every node of $B$ has chain head $m_b$; wires within or off
$B$; \lexpy!_links_ok!; links of \lexcode!bid! from $B$ to outside at depth 0 with multiplicity $m_b$;
other links wholly in $B$ (depth $\geq 1$) or wholly off it.

\paragraph{\texttt{PeelX} (\texttt{PeelOK} on the stripped sides plus 7 fields,
\texttt{peel\_x\_ok}).} \lexpy!peel_ok!; every node of $B$ at depth 0 in the peeled box; for every
node, its \lexcode!cph! reading kept in $R$ at the same depth in its box; for every node of $B$,
\lexpy!_rdc(L, R, bid, c, x, cp[x], o, o2)!; the same two conditions for \lexcode!cdim! over every port.
\lexpy!_rdc! mirrors \lexlean!rdC!: both readings absent, or for a reading $o$ of the peeled box
$o_1 = 0$, the copy's reading has $(k, m) = (1, c)$ and the copy lies at depth
$|\mathit{chain}|$ (outside) of its box; for another box $o_2$ keeps $(k, m)$ and lies one level
shallower.

\paragraph{\texttt{FirstX} (9 fields, \texttt{first\_x\_ok(L, R, bid, m1, B, cp)}).} No aligned
wires on either side; the copy of a $B$ node is outside $B$; off $B$, the \lexcode!cph! and
\lexcode!cdim! readings kept at equal depth; on $B$, \lexpy!_rds! between a node's readings in $L$ and
$R$ and \lexpy!_rdc0! between its reading in $L$ and its copy's in $R$ (for phases and for every
port). \lexpy!_rds! mirrors \lexlean!rdS!: for the peeled box, depth 0 in $L$ and $R$, $o_1 = 0$ and the
kept reading $(1, m_1)$; for another box, kept and at equal depth $\geq 1$. \lexpy!_rdc0! mirrors
\lexlean!rdC0!: for the peeled box, depth 0 in $L$, both coefficients 0 and the copy outside its box;
for another box, kept and one level shallower.

\paragraph{\texttt{PeelN} (\texttt{PeelX} plus 2 fields, \texttt{peel\_n\_ok}).} For every node its
\lexcode!cnd! reading kept at equal depth, and for a node of $B$ \lexpy!_rdc! between its \lexcode!cnd!
reading and its copy's.

\paragraph{\texttt{PeelFlat} (10 fields) and \texttt{PeelFirst} (11 fields),
\texttt{peel\_flat\_ok(L, R, B, cp, mb, c, sd, sd2, first)}.} The slot counts match the
boundaries; each slot's anchor chain equals its leg chain (\lexpy!_leg_chain!); each slot has an
image in $R$, kept slots keep their chain, slots of $B$ have chains $[m_b, r]$ and $[c, r]$ with
$r$ the copy slot's chain; images (kept and copies) are one to one and onto; for every pair of
slots \lexpy!_agree! on their anchors, \lexpy!_copy_vs! (with \lexcode!first!, \lexpy!_copy_vs0!) between a
copied slot and every kept one, \lexpy!_copy_copy! between two copied slots, and with \lexcode!first!
two slots with one head anchor lie both in $B$ or both off it. \lexpy!_copy_vs0! differs from
\lexpy!_copy_vs! only for equal heads, where it requires \lexcode!b! and $E_0 < F_0$ (instead of
$F_0 < E_0$).

\paragraph{Field counts.} \lexpy!_FIELDS = 21! (\lexlean!PeelOK!), \lexpy!_FLAT_FIELDS = 10!
(\lexlean!PeelFlat!; \lexlean!PeelFirst! has one more), \lexpy!_REV_FIELDS = 16!,
\lexpy!_PEELX_FIELDS = 7!, \lexpy!_GHOSTX_FIELDS = 3!, \lexpy!_FIRSTX_FIELDS = 9!,
\lexpy!_DROPX_FIELDS = 4!, \lexpy!_DROPN_FIELDS = 6!, \lexpy!_GHOSTN_FIELDS = 6!; \lexlean!DropOK! has 18 and
\lexlean!GhostOK! 23 fields.

\subsection{Copy readings}

\begin{itemize}
\item \lexpy!_read_copies(L, R, tables, sym, value, bid, cp, reading, phases)!: for each node $x$ of
$L$ in \lexcode!cp! and each \lexcode!cdim! entry of $x$ for box \lexcode!bid!, the copy's port dimension
text must equal $x$'s with \lexcode!sym! at \lexcode!value! (\lexpy!_equal!); the copy then takes $x$'s
dimension index and the reading $(p', \mathit{bid}, \mathit{reading})$. A \lexcode!cph! entry of $x$
for \lexcode!bid! requires the same of the phase index for B and S, and \lexpy!phases(x, copy, value)!
for other kinds; the copy takes $x$'s phase index and the reading. Any mismatch gives \lexpy!None!.
\item \lexpy!_shift_kept(L, R, tables, sym, bid, B, m1, phases)!: removes $R$'s readings of
\lexcode!bid! on $B$, then for each node of $B$ does as above with \lexcode!sym! at $\mathit{sym}+1$,
the kept node taking the reading $(\cdot, \mathit{bid}, 1, m_1)$.
\item \lexpy!_read_ends(L, R, ids)!: each identity of \lexcode!ids! in $R$ (a Z with
\lexlean!outs = ins = [dim]! and no \lexcode!cnd! or \lexcode!cdim! entry yet) wired to a port $P$ with a
\lexcode!cdim! reading $o$ takes $P$'s dimension index, the \lexcode!cdim! readings $o$ on both its ports
and the \lexcode!cnd! reading $o$, on both sides; \lexpy!None! when there is none or a changed identity
is absent from $L$.
\item \lexpy!_apart(*Ds, loose=False)!: over the dimension indices (node dimensions read when the
node has a \lexcode!cnd! entry, port dimensions when the port has a \lexcode!cdim! entry, phase indices
of B and S read when the node has a \lexcode!cph! entry) and separately over the phase indices of the
other kinds, the largest plain index lies below the least read index $k$; with \lexcode!loose!,
below the least $k^2 + k$ (the least \lexlean!Nat.pair k v!).
\end{itemize}

\subsection{Peel \texttt{export\_peel}}

\begin{signature}
\lexpy!export_peel(before: Diagram, after: Diagram) -> str | None!
\end{signature}
When some box of \lexcode!before! is a pair box, the result is
\lexpy!_export_pair_peel(before, after, False)! or else with \lexpy!True!
(\cref{sec:lean-export-pair}). Otherwise it is the first non-\lexpy!None! of
\lexpy!_export_peel(before, after, first)! for \lexcode!first! false (last copy), true (first copy),
and false with \lexcode!reads! (identities read as their legs).

\begin{signature}
\lexpy!_export_peel(before, after, first, reads=False, pairs=None)!
\end{signature}
\begin{algorithm}[\texttt{\_export\_peel}]\label{alg:lean-export-peel}\sloppy
\begin{enumerate}
\item \lexcode!pairs! with \lexcode!reads! gives \lexpy!None!. The mode \lexcode!xm! holds when either side has
a box index; then \lexpy!_seed([before, after], tables)!.
\item \lexpy!peel_ends! (\cref{alg:lean-export-peel-ends}); \lexpy!None! gives \lexpy!None!.
\item \lexpy!L = linked(before, tables, ends_l, held=held, readings=xm)!,
\lexpy!R = linked(after, tables, ends_r, readings=xm)!. $R$'s node ids must be $L$'s plus
\lexcode!fresh!; $R$'s nodes are reordered as $L$'s followed by the copies in $L$ order. $c$ is the
peeled box's multiplicity index in $R$.
\item With \lexcode!xm! and a box index \lexcode!sym!: $R$ becomes \lexpy!_read_copies! with \lexcode!sym! at
$0$ (\lexcode!first!) or at the after multiplicity, reading $(0,0)$ (\lexcode!first!) or $(1, c)$; with
\lexcode!first! then \lexpy!_shift_kept! with $m_1$ the multiplicity index of \lexcode!"1"!. The phase test
substitutes \lexcode!sym! in the before phase by the point (its names replaced by the after
multiplicity's symbols of the same name, else nonnegative integer symbols) and compares with
\lexpy!==! or by \lexpy!repr!.
\item With \lexcode!reads! (requires \lexcode!xm! and not \lexcode!first!): \lexpy!_read_ends! on the identity
ids common to both sides.
\item Last copy: \lexpy!peel_ok!; with \lexcode!xm! \lexpy!peel_x_ok!; with \lexcode!reads!
\lexpy!peel_n_ok!. First copy: on the stripped sides $L_s$, $R_s$ and their flips $L_r$, $R_r$,
\lexpy!rev_ok! in both directions on both sides, \lexpy!peel_ok(Lr, Rr, ...)!, every boundary port of a
$B$ node has a copy outside $B$ (\lexlean!hcp!); with \lexcode!xm! \lexpy!first_x_ok!.
\item $G$ is one more than the largest node id of $L$ and $R$.
\lexpy!_reduce(L, ids_l, set(), G, xm, reads)! and \lexpy!_reduce(R, ids_r, ids_r, G, xm, reads)!
must leave no identity. With \lexcode!pairs!, \lexcode!xm! must hold and no side may have a step or
ghost.
\item The ghost forms \lexpy!linked(before, ghost=G, held=held)! and
\lexpy!linked(after, ghost=G)! (the latter read as in step 4) must equal the reduced diagrams up to
node order (\lexpy!_same_up_to_order!, readings compared as sorted lists).
\item With \lexcode!xm!, \lexpy!_apart(L, R, Dl, Dr, loose=reads)!.
\item The slot data of each side must equal that of its reduced form and be nested;
\lexpy!peel_flat_ok(..., first)! must hold.
\end{enumerate}
\end{algorithm}

The file declares \lexlean!L!, \lexlean!L'! (reduced forms), \lexlean!zero!, \lexlean!Lb!, \lexlean!Lk!
(\lexpy!_alias!: the literal, or \lexlean!L! when the side has no step), both chains (\lexpy!_chain!
with \lexcode!x = xm!, or \lexpy!_chain_n! with \lexcode!reads!), \lexlean!cp!, with \lexcode!first! the flipped
\lexlean!Lbr!, \lexlean!Lkr! and \lexlean!rb!, \lexlean!rb'!, \lexlean!rk!, \lexlean!rk'! (each \lexlean!RevOK!, the
unflipped side \lexlean!Lb.strip! with \lexcode!xm!), the theorem
\lexlean!ok : <K> <pre> <post> bid mb c B cp! with $K$ \lexlean!PeelN! (\lexcode!reads!), \lexlean!PeelX!
(\lexcode!xm!, last copy) or \lexlean!PeelOK!, and $\langle\mathit{pre}, \mathit{post}\rangle$
\lexlean!Lb Lk! or (\lexcode!first!) \lexlean!Lbr Lkr!, then \lexlean!G!, \lexlean!G'!, \lexlean!sdb!, \lexlean!sdk!,
\lexlean!flat : PeelFlat ...! (\lexlean!PeelFirst! with \lexcode!first!), with \lexcode!first! and \lexcode!xm!
\lexlean!okx : FirstX Lb Lk bid m1 B cp! and \lexlean!fo : FirstOK Lb Lbr Lk Lkr bid mb c B cp := ⟨rb, rb', rk, rk', ok⟩!,
with \lexcode!first! alone the theorem \lexlean!hcp!, and \lexlean!main!:
\begin{lstlisting}[style=lexlean]
theorem main : ∀ (E : Env), (∀ k, E.phase zero k = 0) →
    ∀ (ks : ℕ → ℕ), ks mb = ks c + 1 → ∀ y : ℕ → ℕ,
      L'.evalG gsL' E ks (flatAt (L'.slotData G') ks y) =
        L.evalG gsL E ks (flatAt (L.slotData G) ks y) := by
  intro E hE ks hk y
  rw [← sdb, ← sdk]
  exact (chain_k E hE ks _).symm.trans
    ((evalL_peel_flat ok flat E ks hk y).trans (chain_b E hE ks _))
\end{lstlisting}
The variants:
\begin{center}
\begin{tabular}{@{}p{3.2cm}p{4.4cm}p{6.2cm}@{}}
\toprule
Form & Sums & Step term \\
\midrule
last, no index & \lexlean!evalG gsL! & \lexlean!evalL_peel_flat ok flat E ks hk y! \\
last, index & \lexlean!evalGX gsL! & \lexlean!evalX_peel_flat ok flat E ks hk y! \\
first, no index & \lexlean!evalG gsL! & \lexlean!evalL_peel_first rb rb' rk rk' ok flat hcp E ks hk y! \\
first, index & \lexlean!evalGX gsL!; extra \lexlean!ks m1 = 1 →! & \lexlean!evalX_peel_first fo okx flat E ks hk hm y! \\
reads & \lexlean!evalGXR gsL grL! or \lexlean!evalGX gsL! & \lexlean!evalN_peel_flat ok flat E ks hk y! \\
\bottomrule
\end{tabular}
\end{center}

\subsection{Pair-box peels}\label{sec:lean-export-pair}

\lexpy!_strip_pairs(d, owner=None)! returns \lexcode!(d0, owner, specs)! or \lexpy!None!. Without
\lexpy!owner! the diagram must have exactly one box that is no pair box, which becomes the owner; a
given \lexpy!owner! must be one of the boxes that are no pair boxes (the align-merge export passes
its root, \cref{sec:lean-export-merge}). At least one pair box must exist. Each pair box must
pair the owner, have no parent, no copy scalar and multiplicity \lexcode!1!, hold one node, one left
and one right port; the node is a \lexcode!Zp! with two ports, each on exactly one wire, the two
wires ending at the distinct left and right ports; each of these is a port of a Z node of the
owner, off the boundary and on exactly one wire. A spec is
\lexcode!(box, left, right, phase key, swap)! with \lexcode!swap! true when the node's first port
(outputs first) meets the right port; $d_0$ is $d$ without the pair boxes and their nodes.

\lexpy!_export_pair_peel(before, after, first)! requires \lexcode!after! to equal
(\lexpy!kernel._same!) the engine peel \lexpy!peel_one(before, owner, first=first)!, a
\lexpy!BangBoxError! giving \lexpy!None!, and returns
\lexpy!_export_peel(d0, peel_one(d0, owner, first=first), first, pairs=specs)!. The file has no
chains and adds
\begin{lstlisting}[style=lexlean]
def ps : List LPair := [{ box := 3, left := ⟨5, true, 0⟩, right := ⟨6, false, 0⟩,
  ph := 2, swap := false }]
def ps' : List LPair := ps.flatMap fun p => [p, p.fixRight (cpP cp p.right) c]
theorem pf : PairLast Lb Lk ps ps' bid mb c B cp :=
  ⟨by decide +kernel, by decide +kernel, by decide +kernel, rfl⟩
\end{lstlisting}
(for the first copy \lexlean![p.resid m1, p.fixLeft (cpP cp p.left) m1]! and
\lexlean!PairFirst Lb Lk ps ps' bid mb c m1 B cp!). The sums are \lexlean!evalP ps! and
\lexlean!evalP ps'!, and the proof is \lexlean!exact (evalP_peel_flat ok pf flat E ks hk y)! or
\lexlean!exact (evalP_peel_first_flat fo okx pf flat E ks hk hm y)!. The \lexlean!form! field of
\lexlean!PairFirst! and \lexlean!PairLast! (each pair at the peeled box, unfixed, leg chains $[m_b]$)
has no Python mirror.

\subsection{Identity slide \texttt{export\_slide}}

\lexpy!export_slide(before, after)! attaches identities behind the boundary link ports of both sides
(fresh ids from one more than the largest node id), requires equal box lists, reduces both to
their ghost forms (\lexpy!_reduce! without \lexcode!xm!), and searches the nodes of $L$ also in $R$
(equal \lexpy!LNode!) that are phaseless Z nodes other than the identities: the first $x$ for which
dropping $x$ and the identities gives one diagram $C$ with equal ghosts on both sides is the slid
identity. The slot data must agree and be nested. The file declares \lexlean!L!, \lexlean!L'!,
\lexlean!G!, \lexlean!G'!, \lexlean!zero!, \lexlean!Lb!, \lexlean!Lk!, \lexlean!C!, the chains \lexlean!chain_b!,
\lexlean!chain_k!, \lexlean!chain_sb!, \lexlean!chain_sk! (the last two to \lexlean!C! with shared ghosts
\lexlean!gsC!), \lexlean!layout : L'.slotData G' = L.slotData G!, the theorem \lexlean!sums! (the ghost
sums agree) and \lexlean!main! by \lexlean!evalG_flatAt layout!.

\subsection{Regroup \texttt{export\_regroup}}

\lexpy!export_regroup(new, old)! (module \module{archytaszx.semantics.lean_regroup}) requires the two
diagrams to be one linked diagram $L$ (equal \lexpy!linked! images, which carry no groups) with equal
scalars. \lexpy!blocks(d, tables)! lists the slot blocks over \lexcode!outs ++ ins!: an unboxed slot is
the block \lexcode!((), (i,))!; slots of one box group (or of the ungrouped nodes of a box) form one
block \lexcode!((m,), slots)! in order of first slot; \lexpy!None! when a port box exists, a slot lies
in more than one box, or in a box with a parent. \lexpy!covers(L, G)! mirrors \lexlean!Covers!: the
blocks partition the boundary positions and each slot's leg chain (\lexpy!_leg_chain!: the node's
chain extended by the partner's longer chain) equals its block's chain. The theorem is
\begin{lstlisting}[style=lexlean]
theorem main : ∀ (E : Env) (ks : ℕ → ℕ) (x : ℕ → ℕ),
    (layout ks G').Perm (layout ks G) ∧ Function.Bijective (lperm (layout ks G) (layout ks G')) ∧
      L.evalL E ks (readAt (layout ks G') x) =
        L.evalL E ks (readAt (layout ks G) (x ∘ lperm (layout ks G) (layout ks G'))) :=
  regroup_layout L G G' (by decide +kernel) (by decide +kernel)
\end{lstlisting}

\subsection{Align merges \texttt{export\_merge}}\label{sec:lean-export-merge}

\module{archytaszx.semantics.lean_merge} exports a kernel \lexcode{align_merge} step (\cref{ch:kernel}) over \lexlean{LDiagram.evalP} (\cref{ch:lean-sym}); the kernel's \lexpy{merged} calls it and passes the source to \lexpy{linked_net} with the tag \lexcode{M}.

\begin{algorithm}[\texttt{export\_merge(before, after)}]\label{alg:lean-export-merge}
Return \lexpy{None} at the first failing step.
\begin{enumerate}
\item The two sides must have equal scalars. Find the root: the first alignment root $r$ (sorted) of \lexpy{before}'s \bboxes{} that exists, passes the kernel's \lexpy{_align_group}, and whose \lexpy{align_merged(before, r)} is \lexpy{_same} as \lexpy{after} (a \lexpy{KernelError} or \lexpy{BangBoxError} skips $r$).
\item For each side: if it has pair boxes, strip them at $r$ (\lexpy{_strip_pairs}); both sides must yield the same pair specs, and a side without pair boxes is refused when the other had them. Then strip the \lexcode{Zp} nodes (\lexpy{_strip_zp}): each must have one output and one input, lie in exactly one \bbox{} that has no parent, no pair partner, and no links, have each leg on exactly one wire, the two far ends one inside that box and one outside every box, neither on the boundary nor on the node itself, each on one wire. Each becomes $(z,\mathit{box},\mathit{outer},\mathit{inner},\mathit{key},\mathit{out})$ with the phase key and whether the output meets the inner end, and is removed.
\item The two sides' \lexcode{Zp} lists must agree except for the box id. Seed the tables (\lexpy{_seed}) and build both linked mirrors (\lexpy{_xlinked}); a \lexpy{LinkedUnsupported} or any port box refuses. Both sides must have the same node ids; the right side's nodes are put in the left side's order.
\item The multiplicity pairs are $(\mathrm{mult}_R(r),\mathrm{mult}_L(r))$ per pair spec and $(\mathrm{mult}_R(b'),\mathrm{mult}_L(b))$ per \lexcode{Zp} pair; \lexpy{merge_ok(L, R, mults)} checks every field of \lexlean{MergeOK} (\cref{ch:lean-sym}) in Python, with every multiplicity pair equal.
\item The slot data of both sides (\lexpy{slot_data} of their \lexpy{groups}) must have equal counts and equal anchors at every slot.
\item Emit the Lean file: \lexlean{L}, \lexlean{L'} (\lexpy{_xliteral}); \lexlean{ps}, \lexlean{ps'} (the pair specs as \lexlean{LPair}s of box $r$, then each \lexcode{Zp} node as an \lexlean{LPair} with \lexlean{fixL := true} on its own box, phase rows from the phase table); \lexlean{G}, \lexlean{G'} (\lexpy{_lgroups}); \lexlean{ok : MergeOK L L' ps ps'} with sixteen fields, each \lexcode{by decide +kernel}; \lexlean{ss : SameSlots (L.slotData G) (L'.slotData G')} with three such fields; and
\begin{lstlisting}[style=lexlean]
theorem main : ∀ (E : Env) (ks : ℕ → ℕ) (y : ℕ → ℕ),
    L'.evalP ps' E ks (flatAt (L'.slotData G') ks y) =
      L.evalP ps E ks (flatAt (L.slotData G) ks y) :=
  evalP_merge_flat ok ss
\end{lstlisting}
in the namespace \lexlean{ArchytasZX.Export}, after \lexlean{import ArchytasZX}.
\end{enumerate}
\end{algorithm}

\subsection{Power split \texttt{export\_split}}\label{sec:lean-export-split}

\lexpy!export_split(before, after)! (module \module{archytaszx.semantics.lean_split}) exports the
split of one outermost node box of multiplicity $bh$ into $b$ boxes of multiplicity $h$, the
$c$-th new box reading the box index raised by $c\cdot h$, as $b-1$ splits in two over
\lexlean!evalX!.

\paragraph{Extended mirror.} The split mirror extends \lexpy!LDiag! by these types.
\begin{signature}
\lexpy!Align = (a, b, perm)!: an aligned wire, \lexcode!perm! the places of $a$'s leg chain that $b$
reads. \lexpy!Tail = (leg, appended leg)!.\\
\lexpy!PBox(id, mult, parent, port, src)!: a port box holding one port, mirroring node box
\lexcode!src!; ids start at \lexpy!_PBOX = 1000!.\\
\lexpy!XDiag(d: LDiag, cph=[], cdim=[], aligned=[], pboxes=[])!.
\end{signature}
\lexpy!XDiag.partner(q)! is the other end of the first wire at $q$. \lexpy!tail(q)! mirrors
\lexlean!Diagram.tail!: with a partner $r$, the part of $r$'s chain beyond $q$'s node chain (empty
when not longer); without, the port boxes of $q$ walked from $q$'s innermost node box (at most
$|\mathit{boxes}| + |\mathit{pboxes}|$ steps). \lexpy!leg_chain(q)! is the node chain followed by
the tail; \lexpy!port_dim(p)! the dimension index (0 when absent); \lexpy!mates(q)! the other ends of
wires and aligned wires at $q$. \lexpy!_pick(t, perm)! is $[t_i]_{i \in \mathit{perm}}$ (0 out of
range); \lexpy!_set(t, i, x)! replaces entry $i$ in range. \lexpy!_ext(T, n)! mirrors \lexlean!extT!: $n$
with, per tail entry whose leg is on $n$, that leg's dimension appended on its side.
\lexpy!_copied(L, B, cp, cb)! mirrors \lexlean!cph'! and \lexlean!cdim'!: $L$'s readings followed by, for
each reading of a node (port) of $B$, $(\mathit{cp}(x), \mathit{cb}(y), k + u, m)$ with $u = 1$
when $x$ lies at depth 0 in $y$, else $0$.

\paragraph{Readings and alignment.} \lexpy!_seed(ds, tables, names)! numbers, per diagram and node,
every port dimension reading none of \lexcode!names! and every phase (\lexpy!repr(phase)!) with no free
symbol among them, then the empty text. \lexpy!_readings(d, skip, tables)! is as in
\cref{sec:lean-export-readings} except that each entry is $(x, y, 0, m)$ with $m$ the multiplicity
index of $y$, the readings of boxes in \lexcode!skip! are left out, and each rejection returns
\lexpy!None!. \lexpy!_strip(d)! clears every box index and alignment, removes every aligned wire (in
\lexpy!sort_key! order) and records it with the sorted members of the first alignment group whose
boxes hold both ends. \lexpy!_aligned(D, gone, base)! turns each removed wire $(p, q)$ into an
aligned wire:
\begin{enumerate}
\item both ends' node-box walks must be nonempty with their innermost box among the members; the
innermost boxes must differ and share the multiplicity index;
\item $k$ is the common prefix length (at most $\min - 1$); the extra boxes $d_a$, $d_b$ are the
walks between the prefix and the innermost box; both nonempty gives \lexpy!None!; when $d_a$ is
nonempty the ends swap;
\item with $n = |\mathit{walk}_b|$, $\mathit{perm}_i = i$ for $i < k$, $i + 1$ for $k \le i < n-1$
and $k$ for $i = n - 1$;
\item each box of $d_b$ becomes a port box of $p$, ids from \lexcode!base!, the first with parent
$a$'s innermost box and each next with the previous port box as parent.
\end{enumerate}
\lexpy!_xlinked(d, tables, skip)! is \lexpy!_readings!, \lexpy!_strip!, \lexpy!linked! and \lexpy!_aligned(D, gone, 1000)!;
a failure raises \lexcode!a box index read outside the linked form! or
\lexcode!an aligned wire outside the linked form!. \lexpy!_split_reads(d, bid)! holds when some port
dimension of a node of box \lexcode!bid! reads its index.

\paragraph{Renaming and leg order.} \lexpy!_appended(R, fixed, copies)! moves, for every node off
\lexcode!fixed!, the ports of each side whose mates include a node of \lexcode!copies! after the other
ports of that side (orders kept); it returns $R$ relabeled to that order and the map \lexcode!rn!
from the new port to $R$'s port. \lexpy!_renamed(L, R, fixed)! matches the nodes of $L$ missing from
$R$ (\lexcode!gone!) one to one with the nodes of $R$ off $L$ and off \lexcode!fixed! (\lexcode!new!);
unequal counts give \lexpy!None!. A pair $(x, y)$ fits when kind, phase and dimension indices agree,
$y$'s port dimensions start with $x$'s, and for every port of $x$ the boundary membership on
each side agrees, both have equal mate counts at most 1, and a mate of $x$ is the same port (off
\lexcode!gone!), the matched port, or a port of a \lexcode!new! node. Rounds over \lexcode!gone! in order
match every node with exactly one fitting free candidate; a round without progress gives
\lexpy!None!. $R$ is then rewritten with the matched ids in nodes, wires, boundary, box node lists
(sorted), links, readings, aligned wires and port boxes.

\paragraph{Stages and tails.} \lexpy!_below(R, tops)! is the set of boxes in or under \lexcode!tops!.
\lexpy!_stage(L, R, bid, mult, keep, drop_boxes, drop_nodes)! is $R$ without \lexcode!drop_nodes! and
\lexcode!drop_boxes! (and their links), box \lexcode!bid! (and its port boxes) at \lexcode!mult!, each node
of $L$ cut to its own ports plus the appended ports whose mates lie in \lexcode!keep!; those must
form a leading run of the appended ports, else \lexpy!None!; wires, boundary and aligned wires keep
only live ports; the readings are empty. \lexpy!_tails(prev, nxt, B, cp)! lists, for every wire or
aligned wire of \lexcode!prev! from a port $p$ off $B$ to a port $q$ of $B$, the entry $(p, p')$ with
$p'$ the first mate of \lexpy!_cpp(cp, q)! in \lexcode!nxt! on $p$'s node (\lexpy!None! when absent),
sorted by $p'$. \lexpy!_box_map(prev, nxt, src, dst, cp)! maps \lexcode!src! to \lexcode!dst! and each box
strictly under \lexcode!src! to the box in or under \lexcode!dst! whose node tuple is the sorted copies
of its nodes; a missing box gives \lexpy!{}!.

\paragraph{\texttt{BoxSplitOK} (28 fields, \texttt{split\_ok}).}
A port $p$ is \emph{split} when its node is in $B$ or it is a tail key; its image \lexpy!cq(p)! is
its copy for a $B$ node, else its tail value; its split position \lexpy!spos(p)! is 0 for a $B$ node,
else the length of its node chain.
\begin{enumerate}
\item \lexlean!cph!, \lexlean!cdim!: each reading's node (port, not a tail key) exists and keeps its depth
in $L'$; for $B$, the copy lies at the same depth in $\mathit{cb}(y)$, and at depth 0 the
multiplicity index is $c_1$.
\item \lexlean!cph'!, \lexlean!cdim'!: $L'$'s readings equal \lexpy!_copied(L, B, cp, cb)!.
\item \lexlean!nodes!: $L'$'s nodes are \lexpy!_ext(T, n)! of $L$'s nodes followed by the copies of the
$B$ nodes; \lexlean!ids!, \lexlean!legsL!, \lexlean!legs!, \lexlean!bnd!: no repeats.
\item \lexlean!chB!, \lexlean!chO!: a $B$ node has chain $[m_b, r]$ in $L$, $[c_1, r]$ in $L'$, its copy
$[c_2, r]$, depth 0 in \lexcode!bid! on both sides and its copy depth 0 in \lexcode!bid₂!; other chains
kept.
\item \lexlean!tl!: each tail key is a port of a Z or X node off $B$, its value a port of the same
node and side.
\item \lexlean!legS!, \lexlean!legO!, \lexlean!legL!, \lexlean!legs'!, \lexlean!cqinj!, \lexlean!cqnew!: every port
of $L$ stays with its dimension; an unsplit port keeps its leg chain; a split port has entry
$m_b$ at its split position, becomes $c_1$ there in $L'$, and its new image has the chain with
$c_2$ there and its dimension; images are distinct, new, and cover every new port of $L'$.
\item \lexlean!wires!, \lexlean!wires'!: wires join ports of equal leg chain and split status, are kept,
and split wires have both split positions 0 and are copied; every wire of $L'$ is kept or a copy.
\item \lexlean!lwf!, \lexlean!lin!, \lexlean!lout!, \lexlean!links'!: link well-formedness over leg chains; a
split link is wholly split, at depth $\geq 1$, split positions 0, kept and copied; an unsplit link
is wholly unsplit and kept; every link of $L'$ is kept or a copy.
\item \lexlean!al!, \lexlean!al'!: each aligned wire joins ports of equal split status, its places lie
in $a$'s leg chain, $b$'s leg chain is \lexpy!_pick! of $a$'s, it is kept, and a split one has
$\mathit{spos}(b) < |\mathit{perm}|$, $\mathit{perm}_i = \mathit{spos}(a)$ exactly at
$i = \mathit{spos}(b)$, and is copied; every aligned wire of $L'$ is kept or a copy.
\item \lexlean!bndL!, \lexlean!bnd'!: the boundary is kept and extended by the images of split ports.
\end{enumerate}

\paragraph{\texttt{LegPermOK A R rn} (22 fields, \texttt{perm\_ok}).} Equal nodes; unique ids, ports
and boundary of $R$; \lexcode!rn! permutes $A$'s ports, keeps node and side, and moves only ports of
Z and X nodes; leg chains, port dimensions and node chains kept; wires, boundary and links mapped
by \lexcode!rn!; links keep depth and multiplicity and are well formed; wires join equal leg chains;
\lexcode!cph! equal with equal depths; \lexcode!cdim! mapped with ports and equal depths; aligned wires
mapped and well formed; the boundary among the ports. \lexpy!_PERM_FIELDS = 22!,
\lexpy!_FIELDS = 28! (\lexlean!BoxSplitOK!).

\paragraph{\texttt{\_plain\_apart(Xs)}.} Over all stages: node dimensions, port dimensions without
\lexcode!cdim! entry and phase indices of B and S without \lexcode!cph! entry are plain dimensions; port
dimensions with an entry and B, S phase indices with one are read dimensions; the phase indices
of the other kinds split likewise. The largest plain index must lie below the least read index,
separately for dimensions and phases.

\begin{algorithm}[Split export]\label{alg:lean-export-split}
\begin{enumerate}
\item \lexcode!tops!: the outermost boxes new in \lexcode!after!, by decreasing id; \lexcode!split!: the
outermost boxes of both whose multiplicity changed. \lexcode!tops! must be nonempty, \lexcode!split! one
box, the scalars equal. $b = |\mathit{tops}|+1$, $h$ the new multiplicity; the old multiplicity
must expand to $bh$ and every new box to $h$; the index must be kept and present on all new boxes
or on none. No port dimension of the split or new boxes may read its index (\lexpy!_split_reads!).
\item \lexpy!_seed! over both sides with every box index; $L$ = \lexpy!_xlinked(before, tables, {bid})!,
$R_0$ = \lexpy!_xlinked(after, tables, {bid} | tops)!; \lexpy!LinkedUnsupported! gives \lexpy!None!.
\item $B$ the sorted nodes of the split box; every new scope must have $|B|$ nodes; copy maps pair
$B$ with each scope in order. \lexpy!_appended(R0, B | copies, copies)! and \lexpy!_renamed! give $R$
and the leg map \lexcode!rn!; $R$'s ids must be $L$'s plus the copies.
\item Multiplicity indices $m_c$ of $c\cdot h$ for $c = 1..b$; the old multiplicity's index must be
$m_b$. Each node of $B$ must keep its phase in \lexcode!after!; a node whose phase reads the index
(and has no other \lexcode!cph! reading) gets $(x, \mathit{bid}, 0, m_1)$; its copy in the $c$-th new
box must carry the phase with the index replaced by that box's index symbol (\lexpy!_symbol_in!)
plus $c h$, and then takes $x$'s phase index. $R$'s nodes are ordered as $L$'s followed by the
copies, scope by scope.
\item Owners $o_0 = \mathit{bid}$, $o_c = \mathit{tops}_c$. For $k = 1..b-1$: stage $k < b-1$ is
\lexpy!_stage(L, R, tops[k-1], m_{b-k}, ...)! keeping scopes $0..k$ and dropping the boxes under
$\mathit{tops}_{k..}$ and their nodes; stage $b-1$ is $R$ without readings. The step copy map
pairs scope $k-1$ with scope $k$; \lexpy!_box_map! from $o_{k-1}$ to $o_k$ must be nonempty; the
stage readings are \lexpy!_copied! of the previous stage.
\item The last stage's \lexcode!cdim! must equal $R$'s and its \lexcode!cph! outside the owners equal
$R$'s (sorted); \lexpy!_plain_apart! over all stages.
\item With a nonempty \lexcode!rn!, the last stage is named \lexlean!P!, $L'$ is \lexpy!_moved(P, rn)!, and
\lexpy!perm_ok(P, L', rn)! must hold.
\item Per step $k$: \lexpy!_tails! must exist, and \lexpy!split_ok! must hold with
$\mathit{bid} = o_{k-1}$, $\mathit{bid}_2 = o_k$, $m_b = m_{b-k+1}$, $c_1 = m_1$ and
$c_2 = m_{b-k}$.
\end{enumerate}
\end{algorithm}
The file declares each stage \lexlean!L!, \lexlean!M1!, \ldots, then \lexlean!L'! (or \lexlean!P!, \lexlean!L'!,
\lexlean!rn! and \lexlean!okP : LegPermOK P L' rn!), each by \lexpy!_xliteral! (node boxes, then port boxes
\lexlean!⟨id, mult, parent, [], [port]⟩!, and always \lexlean!cph!, \lexlean!cdim!,
\lexlean!aligned := [⟨a, b, [perm]⟩, ...]!); per step \lexlean!cp<k>!, \lexlean!cb<k>! (lookup tables
defaulting to 0), \lexlean!T<k> : List (Port × Port)! and
\lexlean!ok<k> : BoxSplitOK <src> <dst> bid bid₂ mb c₁ c₂ B cp<k> cb<k> T<k>!; and
\begin{lstlisting}[style=lexlean]
theorem main : ∀ (E : Env) (ks : ℕ → ℕ),
    ks 2 = ks 1 + ks 1 →
    ∀ σ' : ℕ → List ℕ → ℕ,
      L'.evalX E ks σ' = L.evalX E ks (bndMap L.base L'.base (splitφ L [4, 5] cp1 T1 (ks 1)) σ') :=
  fun E ks h1 σ' => evalX_split ok1 E ks h1 _
\end{lstlisting}
(shown for $b=2$; for larger $b$ the hypotheses are \lexlean!ks m_{b-k+1} = ks m_1 + ks m_{b-k}!,
the boundary maps nest from \lexlean!L! through the stages, and the proof is
\lexlean!(evalX_split ok<k> E ks h<k> _).trans (...)! down to \lexlean!ok1!; with \lexcode!rn! the innermost
map is \lexlean!(bndMap P.base L'.base (rnφ rn) σ')! and the proof starts with
\lexlean!(evalX_legPerm okP E ks σ').trans!).

% ---------------------------------------------------------------------------
\section{Inductions over Chain and Tree Families}\label{sec:lean-export-chain}

The module \module{archytaszx.semantics.lean_chain} exports a proof-kernel induction
(\cref{ch:kernel}) whose step was settled through a chain cut region (\lexcode!chain_window!) or a
tree window. A family of diagrams indexed by a multiplicity $k$ is written as a Lean
\lexlean!chainNet! (\file{ArchytasZX/ChainInduct.lean}) or \lexlean!treeNet!
(\file{ArchytasZX/TreeInduct.lean}): a base network and copies attached level by level. The
base pair and one period of copy pairs are proved equal by their path sums (the network bodies of
\cref{sec:lean-export-net}), and \lexlean!chain_equal! or \lexlean!tree_equal! gives both families
equal at every $k$. At each checked $k$ the instantiated diagrams are linked to the chain
networks by a label map found by search. Every unsupported input raises
\lexpy!PathSumUnsupported!.

\subsection{Constants}

\begin{center}
\begin{tabular}{@{}lll@{}}
\toprule
Name & Value & Meaning \\
\midrule
\lexpy!W! & $2^{20}$ & label stride between chain levels \\
\lexpy!_FREE! & $2^{19}$ & offset of a copy's own boundary labels within its level \\
\lexpy!_BASE! & $2^{18}$ & shift of the base labels over mixed wire dimensions \\
\lexpy!MAX_PERIOD! & 64 & longest copy period searched \\
\bottomrule
\end{tabular}
\end{center}

\subsection{Splitting a side into base and copies}

\begin{signature}
\lexpy!@dataclass class _Split!: \lexpy!base: Diagram!, \lexpy!copies: list[Diagram]! (one period),
\lexpy!base_cut!, \lexpy!copy_in!, \lexpy!copy_out! (boundary positions of the cut ports in the base, of
the link-in and link-out ports in a copy), \lexpy!ratio! (copy $i+P$'s scalar over copy $i$'s,
default 1).
\end{signature}

\begin{signature}
\lexpy!_split(exposed, family, copy_ids, values, step, offset=None) -> _Split!
\end{signature}
\begin{algorithm}[\texttt{\_split}]\label{alg:lean-export-chain-split}\sloppy
\begin{enumerate}
\item Exactly one \bbox of the exposed successor has a multiplicity containing \lexcode!step!
(\lexcode!no single chain box carries the step index!); it must have links and not be cyclic
(\lexcode!the chain box is cyclic or has no links!). Its link ports map to the copy by the sorted
node-scope correspondence; $q$ are the second ports, $p$ the first.
\item Instantiate at $k=0$ (with \lexcode!offset!: $k=0$ and offset $j$). The copy nodes must survive
(\lexcode!the copy does not survive instantiation at k = 0!). The cut is the set of copy ports wired
to the rest. If the cut equals the $p$ ports and not the $q$ ports, swap $p$ and $q$; the cut
must equal the $q$ ports (\lexcode!the copy meets the rest off its link ports!).
\item The base: the non-copy nodes, with the original boundary on them plus the far ends of the
cut as new boundary ports, its scalar the family's at $k=0$.
\item Copy $j$: the copy nodes of the exposed successor at $k=j$ with their boundary ports plus
the $q$ ports, its scalar the exposed scalar over the family's (a zero family scalar raises
\lexcode!the family at k = j has scalar zero!); a different boundary raises
\lexcode!the copy at k = j has another boundary!.
\item Period: generate copies $j = 1, \ldots, 2\cdot$\lexpy!MAX_PERIOD!; at the first odd $j$ where
the first half of the keys (\lexpy!_key!: nodes, phases, wires, boundary, without scalar) equals the
second, the period is the least divisor $q$ of the half with every key equal to the key at
$t \bmod q$. None raises \lexcode!the copies have no period up to 64!. Extend to $3P+1$ copies.
\item Scalars: periodic, or geometric with ratio $s_P/s_0$ (a zero first scalar raises
\lexcode!the first copy has scalar zero!; otherwise \lexcode!the copy scalars are neither periodic nor geometric!).
\item With \lexcode!offset!, the copy at $k=1$, offset 0 must equal copy 0 with its boundary
(\lexcode!the offset copy depends on k!). Every $p$ port must be on the copy's boundary
(\lexcode!a link-out port is not on the boundary!).
\end{enumerate}
\end{algorithm}
\lexpy!_scaled(split, j)! is copy $j \bmod P$ with its scalar times the ratio to the power
$\lfloor j/P\rfloor$.

\subsection{Bodies over one environment}

\lexpy!_bodies(pairs, names, atoms=False, prefers=())! proves each pair by a network body in one
shared environment: \lexpy!qnet_body! when the first pair is a qudit sum; otherwise \lexpy!net_body!
over one phase modulus (rebuilt at the lcm when the moduli differ). With \lexcode!atoms!, each pair
uses \lexpy!net_amp_body! when its sums carry atoms, \lexpy!net_quarter_body! when quartered, else
\lexpy!net_body!, all over one \lexpy!Coefs! at the lcm modulus (two passes; differing moduli raise
\lexcode!the bodies take different phase moduli!); the coefficient definition becomes the first
body's \lexpy!head!, the atom tables its \lexpy!tabs!.

\subsection{The chain export \texttt{export\_chain}}

\begin{signature}
\lexpy!export_chain(theorem, values=None, name="main", checks=(2,)) -> str!
\end{signature}
\begin{algorithm}[\texttt{export\_chain}]\label{alg:lean-export-chain}
\begin{enumerate}
\item The theorem's inference must be \lexcode!induction! or \lexcode!induction_offset! with
\lexcode!window = "cut"! (\lexcode!not an induction settled through a chain cut region!). It supplies
the step and index symbols, the exposed successor, both families at $k$, the region and the side
(\lexcode!START! or goal).
\item With an offset, the families are the induction hypotheses of both sides and the other side
is offset by the offset symbol. The other side's successor (multiplicity $k+1$) must peel
(\lexcode!the other side's chain box does not peel!).
\item Split both sides; their cut and link positions must agree (\lexcode!the two sides' cut and link positions differ!). The period is the lcm of the two; the ratios raised to the period must agree
(\lexcode!the two sides' copy scalars grow at different rates!).
\item Prove the base pair (\lexlean!baseEq!) and each copy pair $j < P$ (\lexlean!copyEq!) by
\lexpy!_bodies!. Over qubits, the labels' digit counts decide the dimensions: one count $c$ gives
\lexlean!2 ^ c!; several give the mixed export (below).
\item Relabel each copy body (\lexpy!_relabel!): link-in position $t$ to $\mathit{cut}_t$, link-out
position $t$ to $\mathit{cut}_t + W$, any other boundary position $j$ to $W + \_FREE + j$, internal
labels $x$ to $W+x$; the permutation list closes the image by pairing moved-to labels with freed
ones (non-injective raises \lexcode!the copy relabelling is not injective!).
\end{enumerate}
\end{algorithm}
The file imports \lexlean!ArchytasZX.PathSumQuditNet! or \lexlean!ArchytasZX.PathSumNetW!, and
\lexlean!ArchytasZX.ChainInduct!, opens namespace \lexlean!ArchytasZX.ChainProof.<name>!, and contains the
environment, \lexlean!hD : ∀ x y, labDim E x = labDim E y!, namespaces \lexlean!Base! and
\lexlean!Copy<j>! holding the bodies, then
\begin{lstlisting}[style=lexlean]
def W : ℕ := 1048576
def cut : List ℕ := [...]
def πC0 : List (ℕ × ℕ) := [...]
def cOt : List Net := [Copy0.NL.mapL (ren πC0)]
def cXt : List Net := [Copy0.NR.mapL (ren πC0)]
noncomputable def aOt : List ℂ := [(s)]
def cO (i : ℕ) : Net := perIdx cOt ⟨[], [], []⟩ i
noncomputable def aO (i : ℕ) : ℂ := perIdx aOt 1 i      -- r ^ (i / P) * ... when geometric
theorem copyEq0 : Equal E (s) (Copy0.NL.mapL (ren πC0)) (s') (Copy0.NR.mapL (ren πC0)) :=
  Equal.mapL E _ (ren_injective πC0 (by decide) (by decide)) (fun _ => hD _ _) Copy0.copyEq
theorem copyEq : ∀ i, Equal E (aO i) (cO i) (aX i) (cX i) := ...   -- interval_cases on i % P
theorem hcO : ∀ i, copyOK W cut (cO i) = true := ...
theorem main : ∀ k, Equal E
    ((sL) * ∏ i ∈ Finset.range k, aO i) (chainNet W cut Base.NL cO k)
    ((sR) * ∏ i ∈ Finset.range k, aX i) (chainNet W cut Base.NR cX k) :=
  chain_equal E W cut _ _ _ _ hD (by decide) hcO (by decide) hcX _ _ _ _
    Base.baseEq copyEq
\end{lstlisting}
A geometric family adds the lemma \lexlean!copyS! (scaling both sides of an \lexlean!Equal!). With an
offset, \lexlean!main : ∀ s k, ...! gives level $j$ of the chain at $k$ the copy $s+k-1-j$.

\paragraph{Mirrors.} \lexpy!_parse_net(lines, tag)! reads back \lexlean!N<tag>! from a body's lines;
\lexpy!_ren!, \lexpy!_map_net! and \lexpy!_chain_net(cut, base, copies, k, offset)! mirror \lexlean!ren!,
\lexlean!Net.mapL! and \lexlean!chainNet!: level $\ell$ adds copy $\ell$ (or $k-1-\ell$) shifted by
$\ell W$, its nodes before the previous ones, the cut labels at that level summed and removed from
the boundary.

\paragraph{Instance links (\texttt{\_link\_lines}).} For each $k$ in \lexcode!checks!, namespace
\lexlean!At<k>!: each side instantiated at $k$ as a network \lexlean!DL!, \lexlean!DR!
(\lexpy!_instance_net!: boundary port $i$ is label $i$, wire $w$ label $(\text{boundary size or base}) +
w$; only Z, X, F, W nodes, no \bboxes (\lexcode!the instance keeps a bang box!), no W on qudits); a
label map from the chain network onto the instance by backtracking search over nodes in
breadth-first order of shared labels, matching kind, dimension, phase vector and ports in order
(\lexpy!_iso!; failures raise \lexcode!the chain and the instance have different node counts! or
\lexcode!the chain at k is not the instance's network!); the boundary must map onto
$0..b-1$; theorems \lexlean!linkL!, \lexlean!linkR! (the mapped chain network's nodes, internal and
boundary labels are permutations of the instance's, by \lexlean!decide!), \lexlean!instK! by
\lexlean!Equal.mapL2! and \lexlean!reorderL!/\lexlean!reorder!, \lexlean!scalL!, \lexlean!scalR! (the chain
scalar product equals the instance scalar, by \lexpy!prove!), and
\lexlean!inst := instK.congr_scalar scalL scalR!.

\paragraph{Mixed wire dimensions (\texttt{\_export\_mixed}).} When the qubit labels carry
different digit counts: labels must stay below \lexpy!_BASE! (\lexcode!a body label reaches the base shift!) and have digit counts (\lexcode!a network label has no digit count!); per body a local
environment \lexlean!{ E with labs := labsN }!; the base networks shifted by \lexlean!OB = 2^18!; copy
labels in $[W, 2W)$ with one count each (\lexcode!the copies give one chain label two dimensions!);
\lexlean!baseEq! and \lexlean!copyEq<j>! transported by \lexlean!relabs! and \lexlean!mapL_on! to
\lexlean!lvEnv E W C B 1!; \lexlean!main! by \lexlean!chain_equal_lv! over \lexlean!lvEnv E W C B k!; the
instances use \lexlean!EK := { E with labs := instLabs ++ lvls W C k ++ B }! with the right
instance's wires labeled from $2^{17}$. The file imports \lexlean!ArchytasZX.ChainMixed!.

\subsection{Tree families}

\begin{signature}
\lexpy!export_tree(bases, copies, leaf, root, kids, name="main", checks=()) -> str!\\
\lexpy!export_tree_theorem(theorem, values=None, name="main", checks=(2,)) -> str!\\
\lexpy!@dataclass class TreeCheck!: \lexpy!depth!, \lexpy!left!, \lexpy!right!, \lexpy!fan! (\lexcode!"L"! or
\lexcode!"R"!, the side whose instance is the fused tree).
\end{signature}
Copy $n$ is plugged in heap order: copy 0's port \lexcode!root! to the base's boundary position
\lexcode!leaf!, copy $t+1$'s to child \lexcode!kids[t mod b]! of copy $\lfloor t/b\rfloor$; depth $d$
holds $\sum_{j<d} b^{j}$ copies (\lexpy!_tree_count!). \lexpy!_heap_plug! and \lexpy!_tree_net! mirror
\lexlean!heapPlug W leaf kids! and \lexlean!treeNet W r p base copies n!; \lexpy!_grow_tree! builds the
grown diagram; \lexpy!_tree_pis! relabels the copy bodies (root to 0, boundary $j$ to $W+\_FREE+j$,
internal $x$ to $W+x$).

\lexpy!export_tree! proves the pairs by \lexpy!_bodies(..., atoms=True)!; for each check it grows the
fan side's base to the depth, reorders its boundary as the tree side's instance orders the
family's labels (\lexpy!_fused_tree!, \lexpy!_reorder!: an output landing among inputs raises
\lexcode!the fused tree's boundary does not match the instance's!), and adds a fusion pair (the grown
tree against the fan instance) whose left nodes keep the family's phase indices
(\lexcode!prefer!). More than one digit count raises \lexcode!the tree export needs one wire dimension!.
The file imports the path-sum modules the bodies need (\lexlean!PathSumNetAmp!, \lexlean!PathSumSupport!,
\lexlean!PathSumChain! when used) and \lexlean!ArchytasZX.TreeInduct!; it declares \lexlean!W!,
\lexlean!r := 0!, \lexlean!leaf0!, \lexlean!kids!, \lexlean!p := heapPlug W leaf0 kids!, the copy tables,
\lexlean!hcO! (\lexlean!tcopyOK!), \lexlean!hpO : PlugOK W p Base.NL cO! by \lexlean!heap_plugOK!, and
\lexlean!main : ∀ n, Equal E (... ) (treeNet W r p Base.NL cO n) (...) (treeNet W r p Base.NR cX n)!
by \lexlean!tree_equal!. Each \lexlean!At<k>! (\lexpy!_tree_link!) holds the fusion body in namespace
\lexlean!Fuse!, the tree instance \lexlean!D!, the links, \lexlean!instK!, \lexlean!scalL!, \lexlean!scalR! and
\lexlean!inst! (composed with \lexlean!Fuse.fuseEq!).

\lexpy!export_tree_theorem! requires an \lexcode!induction! with \lexcode!window = "tree"!
(\lexcode!not an induction settled through a tree window!) with exactly one side a tree box
(\lexcode!not exactly one side is a tree box!). The tree side (\lexpy!_tree_side!) must be one
top-level tree box of multiplicity in the index only, with a constant copy scalar, no level,
position or index, an input parent port and output child ports, no other boundary ports in the
copy, a scalar that is the product of copy scalars at depth 2, and one plug on the base boundary
(messages \lexcode!the tree side is not one top-level tree box!, \lexcode!the tree box has another multiplicity!, \lexcode!the tree box's copy scalar is not a constant!, \lexcode!the tree box reads an index!, \lexcode!the tree's parent port is no input!, \lexcode!a tree child port is no output!,
\lexcode!the copy has boundary ports off its child ports!, \lexcode!the tree family's scalar is not the copy scalars' product!, \lexcode!the copy meets the base off its parent port!, \lexcode!the plug is not on the base's boundary!, \lexcode!the copy's child ports are not its outputs!). The fan side
(\lexpy!_fan_side!) must be one Z or X spider (or a W node with one output and one input) with one
output port boxed $b^{n}$ times, $b$ the tree's branching; its copy is a phase-free spider of the
same kind with $b$ outputs and one input, or for a W node $c\cdot\mathrm{W}(1\to b)$ after a
weighted Z read off the symbolic contraction of the tree copy (\lexpy!_w_fan_copy!). The plug
positions must agree (\lexcode!the two sides plug at different boundary positions!); the call is
\lexpy!export_tree(pairs[0], [pairs[1]], leaf, b, list(range(b)), name, links)!.

% ---------------------------------------------------------------------------
\section{Periodic Theorems}\label{sec:lean-export-periodic}

The module \module{archytaszx.semantics.lean_periodic} exports a theorem of the proof
kernel's \lexcode!periodic! inference: one side holds a chain \bbox whose index runs over all
naturals, the other side is fixed, and the equality holds at every index because it holds at
$0, \ldots, p-1$ and $p$ further copies multiply the chain side by a constant $c$. The chain side
is written as \lexlean!famNet W u tl cp hd bd k! (\file{ArchytasZX/PeriodicChain.lean}) and the other
as a fixed network \lexlean!G!; \lexlean!fam_periodic! gives the two equal at every index.

\begin{signature}
\lexpy!export_periodic(theorem, values=None, name="main", checks=(2, 3)) -> str!
\end{signature}
\begin{algorithm}[\texttt{export\_periodic}]\label{alg:lean-export-periodic}
\begin{enumerate}
\item The inference must be \lexcode!periodic! (\lexcode!not a periodic theorem!); it gives the index and
the step $p$. Both sides are instantiated at \lexcode!values! when given.
\item \lexpy!chain_period! (\module{archytaszx.semantics.periodic}) must find a chain on exactly one
side (\lexcode!the export needs a chain on exactly one side!), with its period and chain factor
$c_{\mathrm{chain}}$. The chain \bbox (the one whose multiplicity holds the index) must carry no
copy scalar other than 1 (\lexcode!the chain box carries a copy scalar!).
\item \emph{Sides of the chain.} For every port of a box node: a link tail port ($p$ of a link)
must be wired to a node outside the box (\lexcode!a link's tail end is not a node outside the box!); any other non-link port must be internal to the box (\lexcode!a box node meets the outside off its links!). The head side is every node outside the box reachable from the partners of the
links' $q$ ports without entering the box; the tail side likewise from the $p$ ports; they must
be disjoint (\lexcode!the two ends of the chain meet outside the box!); every other outside node joins
the head side.
\item Labels: link $t$ is $W+t$ ($W = 2^{20}$), $u = [W, W+1, \ldots]$; boundary slot $i$ is
$W + t$ when its port is the $q$ port of link $t$, else $i$ (the list \lexlean!bd!).
\item $c = c_{\mathrm{chain}}^{p/\mathrm{period}}$; $A$ = \lexpy!shift_ratio(chain, index, p)!.
Instantiate the chain side at $0..p$ and the goal side at $0..p-1$ (renamed nodes outside the
box raise \lexcode!instantiation renamed a node outside the box!).
\item \emph{Open parts}: the instance without its head side, its own boundary ports then its cut
ports (ports wired to the head; a head end wired to the head raises \lexcode!a head end is not wired off the head side!). The block pair is (open part at $p$ with scalar 1, open part at $0$ with
scalar $c$); their input directions must agree (\lexcode!the open parts' boundaries differ in direction!).
\item Prove the pairs (chain at $r$, goal at $r$) for $r<p$ (\lexlean!baseEq!) and the block pair
(\lexlean!blkEq!) by \lexpy!_bodies! (\cref{sec:lean-export-chain}); qudit sums or weights raise
\lexcode!the periodic export covers unweighted qubit networks!; several digit counts raise
\lexcode!the periodic export covers one wire dimension!. Equal phase vectors are canonicalized to their
first index in every body line (\lexpy!_canon_phases!).
\item \emph{Templates} (\lexpy!piece!): \lexlean!tl! (head side, its ends at link labels), \lexlean!hd!
(tail side), \lexlean!cp! (box nodes, $q$ ports at $W+t$, $p$ ports at $2W+t$), internal wires on
fresh labels (from the slot count, and from $W + |\mathit{links}|$ for the copy); a port leaving
its piece raises \lexcode!leg <ref> of <id> leaves its piece!; a phase vector not registered by a
body raises \lexcode!node <id> has a phase vector no proof registers!. Only Z, X, F nodes.
\item Relabel each body onto its template by \lexpy!_pull! (a backtracking label map seeded with the
boundary): \lexlean!π<r>! for \lexlean!famNet ... r!, \lexlean!ρ<r>! for \lexlean!G! (the goal network at 0
relabelled by \lexlean!ρ!), \lexlean!πB! for \lexlean!blkNet W u cp hd p!, \lexlean!πH! for
\lexlean!headNet u hd!.
\end{enumerate}
\end{algorithm}
\lexpy!_fam!, \lexpy!_blk! and \lexpy!_sh! mirror \lexlean!famNet!, \lexlean!blkNet! and \lexlean!sh W k x! ($x$
for $x<W$, else $x+kW$): \lexlean!famNet! is \lexlean!tl!, then $k$ copies of \lexlean!cp! shifted by
$jW$, then \lexlean!hd! shifted by $kW$, with the copies' internal labels and the shifted link labels
summed and the boundary \lexlean!bd!.

The file imports \lexlean!ArchytasZX.PathSumNetW! and \lexlean!ArchytasZX.PeriodicChain!, opens
namespace \lexlean!ArchytasZX.PeriodicProof.<name>!, and declares the environment, \lexlean!hD!,
namespaces \lexlean!Base<r>! and \lexlean!Blk!, the templates and permutations, the scalar tables
\lexlean!sT!, \lexlean!gT!, \lexlean!A!, \lexlean!c!, and
\begin{lstlisting}[style=lexlean]
theorem ok : famOK W u tl cp hd bd = true := by decide
theorem h0 : Equal E 1 (blkNet W u cp hd 2) c (headNet u hd) := ...
theorem base0 : Equal E (s0) (famNet W u tl cp hd bd 0) (g0) G := ...
theorem base : ∀ r < 2, Equal E (perS sT A 2 r) (famNet W u tl cp hd bd r)
    (perS gT (A * c) 2 r) G := ...
theorem main : ∀ k, Equal E (perS sT A 2 k) (famNet W u tl cp hd bd k)
    (perS gT (A * c) 2 k) G :=
  fam_periodic E hD ok 2 (by norm_num) c A h0 G _ _ (perS_step _ _ _ (by norm_num))
    (perS_step _ _ _ (by norm_num)) base
\end{lstlisting}
(shown for $p=2$). \lexlean!perS tab r p k! is the table entry at $k \bmod p$ times $r^{\lfloor
k/p\rfloor}$. For each $k$ of \lexcode!checks!, namespace \lexlean!At<k>! (\lexpy!_check_lines!) writes both
instances \lexlean!DL!, \lexlean!DR!, the maps \lexlean!πL!, \lexlean!πR!, \lexlean!linkL!, \lexlean!linkR!,
\lexlean!instK! from \lexlean!main k!, \lexlean!scalL!, \lexlean!scalR! (the table scalar at $k$ equals the
instance's, proved from the base scalar at $k \bmod p$ times $A^{\lfloor k/p\rfloor}$, resp.\
$(Ac)^{\lfloor k/p\rfloor}$) and \lexlean!inst!.

% ---------------------------------------------------------------------------
\section{Adder Families}\label{sec:lean-export-adder}

\module{archytaszx.semantics.lean_adder} proves a theorem in one multiplicity $n$ for every $n$ when both sides are the flat networks of the Fourier constant adder or the Draper adder of \cref{sec:boxes-fca-eng,sec:boxes-draper-eng}.

\begin{signature}
\py!export_family(lhs: Diagram, rhs: Diagram, symbols: tuple[str, ...], name: str = "main") -> str | None!
\end{signature}

\py!FAMILIES! lists, in order, the families \code!FcaEng! (tag \code!Fca!, scalars $1$) and \code!DraperEng! (tag \code!Draper!, scalars equal by \py!sympy.simplify!), each with Python mirrors of \lean!engL!, \lean!engR!, \lean!engLabs!, and the phase table (\py!eng_l!, \py!eng_r!, \py!eng_labs!, \py!eng_phases!; \py!drp_l!, \py!drp_r!, \py!drp_labs!, \py!drp_phases!).

\begin{algorithm}[\texttt{export\_family}]\label{alg:lean-export-adder}
\begin{enumerate}
\item With other than exactly one symbol, return \py!None!.
\item For each family in order (\py!_matches!): instantiate both sides at $k=1,2,\ldots,\texttt{CHECKED},0$ (\py!CHECKED = 5!); require the scalar condition, the network texts of the network exporter's labelling equal to the mirrors at $k$, the dimension table equal to \py!eng_dims(k, d)! and the phase table to the mirror's ($d$ the first dimension at $k=1$; at $k=0$ the dimension table \code![1]!), and the listed labels of nonzero index equal to the mirror's. The first match gives $d$; an exporter exception counts as no match; none matching returns \py!None!.
\item Write the namespace \lean!ArchytasZX.Export.<tag>! with, for $k=0,\ldots,5$, \lean!litL<k>! and \lean!litR<k>! (the exports) and \lean!engL_lit<k>!, \lean!engR_lit<k>! by \lean!decide +kernel!, and the theorem \lean!name!: for \code!FcaEng!, \lean!∀ n, Equal (FcaEng.engE d n) 1 (FcaEng.engL n) 1 (FcaEng.engR n)! by \lean!FcaEng.engFam!; for \code!DraperEng!, the same with a common scalar \lean!s : ℂ! by \lean!DraperEng.engFamS!. A symbolic $d$ adds \lean!∀ (d : ℕ), 0 < d →!; a numeric $d$ is discharged by \lean!norm_num!.
\item Write both boxed sides as \lean!FamDiagram! literals \lean!famL!, \lean!famR! (\py!fam_literal!: node boxes, chain links, pair boxes, and the phase, dimension and leg indices, each dimension read as a power of $d$), the checks \lean!famL_eq : famL = <F>.D0! and \lean!famR_eq : famR = <F>.D1! by \lean!rfl! (\lean!<F>! is \lean!FamFca! or \lean!FamDraper!), and \lean!<name>_diagram!: \lean!∀ n, DEqual (<ns>.engE d n) 1 (famL.famExpand n) 1 (famR.famExpand n)! by \lean!FamFca.famFca! or \lean!FamDraper.famDraper! (\cref{sec:boxes-famdraper}). The corollary reader of \cref{sec:lean-export-restrict} still reads \lean!name!.
\end{enumerate}
\end{algorithm}

\section{Restricted Corollaries}\label{sec:lean-export-restrict}

The module \module{archytaszx.semantics.lean_restrict} restates the final theorem of any
export at the clauses of a restriction (\module{archytaszx.semantics.restrict},
\cref{ch:proofs}). It parses Lean text; it builds no proof of its own beyond applying the final
theorem. Failures raise \lexpy!RestrictUnsupported!.

\paragraph{Reading the final theorem (\texttt{\_final}).} The final theorem is the last declared
\lexcode!theorem! among \lexcode!<name>_diagrams_original!, \lexcode!<name>_diagrams! and \lexcode!<name>!
(none raises \lexcode!no theorem <name> in the export!). Its text up to the first \lexcode!:=! outside
brackets (\lexcode!the theorem has no proof!; brackets \lexcode!([{⟨! and \lexcode!)]}⟩!, unbalanced ones
raise \lexcode!unbalanced brackets in the theorem!) splits into leading binder groups
\lexcode!(names : type)!, \lexcode!{...}!, \lexcode![...]! (a group without \lexcode!:! raises
\lexcode!a binder without a type: ...!), then \lexcode!:! (\lexcode!the theorem's statement does not follow its binders!) and the body. Leading guards \lexcode!c ≤ n → ! of the body (an integer, then an identifier) become further binders \lexcode!(h_n : c ≤ n)!. A body starting with \lexcode!∀! contributes further binders, bracketed or
one plain \lexcode!names : type,! group (\lexcode!a ∀ statement this reader does not parse!, \lexcode!a ∀ statement without a comma!).

\begin{signature}
\lexpy!corollary(source, clauses, name="main", var=lambda s: s) -> str!
\end{signature}
\begin{algorithm}[\texttt{corollary}]\label{alg:lean-export-corollary}
Per clause, with $v$ = \lexcode!var(symbol)!:
\begin{enumerate}
\item A side clause with a domain adds its hypotheses after $v$'s binder (\lexpy!_hypotheses!:
\lexcode!h_v_side_lo : lo ≤ v! or \lexcode!<!, \lexcode!h_v_side_hi!, \lexcode!h_v_side_ne<k> : v ≠ x!).
\item A clause that does not substitute is skipped.
\item A shift clause binds the fresh $m$ and writes $v = (c + a\cdot m)$, with
\lexlean!(m : ℝ)! and an integer binder \lexlean!(m : ℤ)! when the clause has a lattice $q$; the fresh
domain's hypotheses and the excluded points (on the value) follow.
\item An equality writes \lexlean!(c : T)!, $T$ the type \lexlean!ℕ! or \lexlean!Nat! of $v$'s binder when it has one, else
\lexcode!c!.
\item A non-dimension $\ge$ clause adds \lexlean!(h_v_ge : c ≤ v)! after $v$.
\item Otherwise (an affine clause or a dimension floor) bind the fresh $m$ and write
$v = (a * m + c)$.
\end{enumerate}
Then each binder group's type is substituted; a round group whose type changed and holds no letters (typed numerals \lexcode!(c : ℕ)! read as 0) is closed and discharged by \lexlean!(by norm_num)!; the names keep, drop or replace their
binders and pass their values as arguments (round groups only). The body is substituted. The
theorem is \lexlean!<name>_at! when every substituting clause is an equality, else
\lexlean!<name>_restricted!, with the doc comment \lexcode!The final theorem at <clauses>.! and the proof
\lexcode!<final> <args>!; it is inserted before the source's closing \lexcode!end! line.
\end{algorithm}
Substitution (\lexpy!_sub!) replaces whole identifiers only (not preceded by a word character,
apostrophe, dot or subscript digit, nor followed by a word character, apostrophe or subscript
digit). \lexpy!_num! parenthesizes values containing \lexcode!/! or starting with \lexcode!-!.

\paragraph{Shift corollaries.} \lexpy!shift_corollary(source, clauses, name="main")! restates the
final theorem, stated at $n = k + c$, over $n$ for each non-equality clause: with one clause the
step binder is the last binder (it must be a single round binder, else \lexcode!the final theorem does not end in its step binder!); with several, each clause's binder is the unique plain (\lexlean!ℕ! or \lexlean!Nat!)
\lexlean!ℕ! binder whose value \lexcode!(k + c)!, \lexcode!(a * k)! or \lexcode!((a * k) + c)! the body holds
(\lexpy!_step_binders!; otherwise \lexcode!no single step binder for <clause>!). For $a=1$, $c\ne0$ the body
must hold \lexcode!(k + c)! (\lexcode!the final theorem is not stated at k + c!), which becomes $n$
(primed when taken); the binder becomes \lexlean!(n : T) (h_n_ge : c ≤ n)! ($T$ the step binder's type) and the proof
\lexlean!obtain ⟨k, rfl⟩ : ∃ k, n = k + c := ⟨n - c, by omega⟩!, \lexlean!exact final args!. Clauses
with $a>1$ keep their binder. The source is returned unchanged when every clause is an equality.
\lexpy!header(source, note)! appends \lexcode!note! as the last lines of the leading \lexcode!/- -/!
comment, or prepends a new comment.

